% Propietario conservado: /Users/ruben/Documents/New project/output/EXTREMA_MEDIA_RAZON_RECIPROCIDAD_ES_EN_20260918/ARTICULO/ES/source/nuclear/antecedentes_fisicos.tex
\chapter{Canonical realization of covariance}
\label{x_reciprocidad:nuclear:covarianza-antecedente}

% BEGIN OWNER /Users/ruben/Documents/New project/output/EXTREMA_MEDIA_RAZON_RECIPROCIDAD_ES_EN_20260918/ARTICULO/ES/source/antecedentes_ii/sections/07c_covarianza_accion.tex
% Propuesta de inserción después de 07b_geometria_elipse.tex.
% RESULTADO_RECUPERADO: c52_06_01_incertidumbre.tex, líneas 10--91.
% Dependencias: c52_04_01_ley_areal.tex, líneas 137--198;
% c52_06_02_intercambio_constitutivo.tex, líneas 5--29.
% Explicitación local: dominio operatorio, prueba de covarianza,
% transporte recíproco y distinción dimensional posición--momento.
% Este fragmento no introduce una constante ni un generador adicional.

\section{Canonical covariance and reciprocity of the action ellipse}
\label{x_reciprocidad:sec:ii-area-covarianza}

The preceding construction determines two complementary invariants: the product of the semiaxes fixes the action area and their quotient preserves the anisotropy of the angular pair. We now specify the canonical realization in which these semiaxes are two standard deviations. The starting structure is the ellipse already obtained from the angular and action expressions of APP--TRIT--TPK; the operator realization is established after that construction.

In this subsection we write \(\eta=\eta_{\rm el}
=\operatorname{artanh}(C^*/A)\), with \(|C^*/A|<1\) and \(\hbar>0\). The angular quotient uses the same unit for both its components. In the balanced coordinates \((Q,P)\), both of square-root-of-action dimension, we define
\begin{equation}
 V_\eta:=\frac14
 \begin{pmatrix}a_S^2&0\\0&b_S^2\end{pmatrix}
 =\frac{\hbar}{2}
 \begin{pmatrix}e^\eta&0\\0&e^{-\eta}\end{pmatrix}.
 \label{x_reciprocidad:eq:ii-covarianza-geometrica}
\end{equation}
This positive matrix has entries of action dimension. Its interpretation as a quantum covariance is specified by the following realization.

\begin{proposition}[Gaussian realization of the action covariance]
\label{x_reciprocidad:prop:ii-covarianza-gaussiana}
Consider the canonical representation on \(L^2(\mathbb R,dQ)\), with
\(\widehat Qf(Q)=Qf(Q)\) and
\(\widehat Pf(Q)=-i\hbar f'(Q)\), initially on the Schwartz
space \(\mathcal S(\mathbb R)\). The normalized state
\begin{equation}
 \psi_\eta(Q)=\bigl(\pi\hbar e^\eta\bigr)^{-1/4}
 \exp\!\left(-\frac{Q^2}{2\hbar e^\eta}\right)
 \label{x_reciprocidad:eq:ii-gaussiana-accion}
\end{equation}
has zero means and covariance matrix \(V_\eta\). In particular,
\begin{equation}
 (\Delta Q)^2=\frac{\hbar e^\eta}{2},\qquad
 (\Delta P)^2=\frac{\hbar e^{-\eta}}{2},\qquad
 \operatorname{Cov}(Q,P)=0,\qquad
 \det V_\eta=\frac{\hbar^2}{4}.
 \label{x_reciprocidad:eq:ii-varianzas-accion}
\end{equation}
This realization saturates the Robertson--Schrödinger inequality.
\end{proposition}

\begin{proof}
The preceding operators are essentially self-adjoint on
\(\mathcal S(\mathbb R)\), which is a common invariant core;
on this space,
\([\widehat Q,\widehat P]=i\hbar I\). For
\(\widehat Y=(\widehat Q,\widehat P)\), the symmetric covariance is
\[
 V_{jk}=\frac12\left\langle
 (\widehat Y_j-\langle\widehat Y_j\rangle)
 (\widehat Y_k-\langle\widehat Y_k\rangle)
 +(\widehat Y_k-\langle\widehat Y_k\rangle)
 (\widehat Y_j-\langle\widehat Y_j\rangle)
 \right\rangle.
\]
With \(s=\hbar e^\eta>0\), let \(I_s=\int_{\mathbb R}e^{-Q^2/s}\,dQ\). The product of two equal integrals, evaluated in polar coordinates, gives \(I_s^2=2\pi\int_0^\infty r e^{-r^2/s}\,dr=\pi s\). Moreover, the integral of \(\frac{d}{dQ}(Qe^{-Q^2/s})\) is zero, owing to decay at both endpoints; hence \(\int_{\mathbb R}Q^2e^{-Q^2/s}\,dQ=sI_s/2\). These identities and parity prove normalization, the zero mean and \(\langle\widehat Q^2\rangle=s/2\). Moreover, \(\widehat P\psi_\eta=i e^{-\eta}Q\psi_\eta\). Therefore, \(\langle\widehat P\rangle=0\), \(\|\widehat P\psi_\eta\|^2=\hbar e^{-\eta}/2\) and \(\operatorname{Re}\langle\widehat Q\psi_\eta,
\widehat P\psi_\eta\rangle=0\).

For any normalized state in that domain, Cauchy--Schwarz applied
to \((\widehat Q-\langle\widehat Q\rangle)\psi\) and
\((\widehat P-\langle\widehat P\rangle)\psi\) gives
\[
 (\Delta Q)^2(\Delta P)^2
 \geq \operatorname{Cov}(Q,P)^2
       +\frac14\left|\langle[\widehat Q,\widehat P]\rangle\right|^2
 =\operatorname{Cov}(Q,P)^2+\frac{\hbar^2}{4}.
\]
The computed variances attain equality. Equivalently, the
operator
\[
 a_\eta=\frac{e^{-\eta/2}\widehat Q
                  +i e^{\eta/2}\widehat P}{\sqrt{2\hbar}}
\]
satisfies \([a_\eta,a_\eta^\dagger]=I\) and
\(a_\eta\psi_\eta=0\), by direct substitution.
\end{proof}

\begin{proposition}[Area, dispersion, and reciprocal transport]
\label{x_reciprocidad:prop:ii-area-covarianza-reciprocidad}
The action ellipse is exactly the covariance sublevel set
\begin{equation}
 \mathcal E_\eta
 =\left\{y=(Q,P)^{\mathsf T}:y^{\mathsf T}V_\eta^{-1}y\leq4\right\}
 =\left\{\frac{Q^2}{a_S^2}+\frac{P^2}{b_S^2}\leq1\right\}.
 \label{x_reciprocidad:eq:ii-elipse-covarianza}
\end{equation}
Its area is \(h=2\pi\hbar\), its semiaxes are
\(a_S=2\Delta Q\), \(b_S=2\Delta P\), and
\begin{equation}
 \frac{\Delta P}{\Delta Q}=e^{-\eta}
 =\frac{b_S}{a_S}=R_{\rm el}^2.
 \label{x_reciprocidad:eq:ii-covarianza-forma}
\end{equation}
The canonical rotation \(J_2\) of Section~\ref{x_reciprocidad:iv:ellipse:section} transports
\(V_\eta\) to \(V_{-\eta}\) and preserves the symplectic form.
\end{proposition}

\begin{proof}
Inverting the diagonal matrix
\eqref{x_reciprocidad:eq:ii-covarianza-geometrica} gives
\eqref{x_reciprocidad:eq:ii-elipse-covarianza}. Its semiaxes are the positive roots of
four times its variances, and
\[
 \operatorname{Area}(\mathcal E_\eta)
 =4\pi\sqrt{\det V_\eta}=2\pi\hbar=h,
\]
in agreement with the areal law~\eqref{x_reciprocidad:iv:ellipse:area-fase}.
The ratio of the positive roots proves
\eqref{x_reciprocidad:eq:ii-covarianza-forma}.
If \(M_\eta=\operatorname{diag}(e^{\eta/2},e^{-\eta/2})\), then
\[
 V_\eta=M_\eta V_0M_\eta^{\mathsf T},\qquad
 M_\eta^{\mathsf T}J_2M_\eta=J_2,\qquad
 J_2V_\eta J_2^{\mathsf T}=V_{-\eta}.
\]
These identities follow by multiplication. In operator terms,
\((\widehat Q',\widehat P')=(-\widehat P,\widehat Q)\) preserves
\([\widehat Q',\widehat P']=i\hbar I\). The permutation
\(S_2:(Q,P)\mapsto(P,Q)\) also exchanges the variances but reverses
the sign of the commutator and of the symplectic form. Thus covariance
transport respects the distinction between canonical reciprocity and
orientation reversal already established in Section~\ref{x_reciprocidad:iv:ellipse:section}.
\end{proof}

\subsubsection{Position, momentum, and dimensional transduction.}
Let \(\lambda_{xp}>0\) be a factor of the dimensional chart, with the dimension of momentum per unit length. The transformation
\[
 Q=\sqrt{\lambda_{xp}}\,x,\qquad
 P=\frac{p_x}{\sqrt{\lambda_{xp}}}
\]
realizes the balanced pair as a position \(x\) and its conjugate momentum
\(p_x\); it preserves \(dQ\wedge dP=dx\wedge dp_x\) and
\([\widehat x,\widehat p_x]=i\hbar I\). By transformation of
covariance,
\begin{equation}
 V_\eta^{(x,p_x)}=\frac{\hbar}{2}
 \begin{pmatrix}e^\eta/\lambda_{xp}&0\\
                 0&\lambda_{xp}e^{-\eta}\end{pmatrix},\qquad
 \Delta x\,\Delta p_x=\frac{\hbar}{2},\qquad
 \frac{\Delta p_x}{\Delta x}=\lambda_{xp}e^{-\eta}.
 \label{x_reciprocidad:eq:ii-covarianza-dimensiones}
\end{equation}
The dispersion ratio acquires the dimension of momentum per length here. By contrast, \eqref{x_reciprocidad:eq:ii-covarianza-forma} belongs to the balanced chart. For the inductive--capacitive realization already defined, \eqref{x_reciprocidad:hmtII:eq:ii-elipse-impedancia} supplies the dimensionless identity \(Z_\eta/Z_*=e^{-\eta}=\Delta P/\Delta Q\). Each reading uses its declared transduction and preserves the same action determinant.

\subsubsection{Contour energy and mean energy of the state.}
The inductive--capacitive realization has Hamiltonian
\[
 H_\eta(Q,P)=\frac{\omega}{2}
 \left(e^{-\eta}Q^2+e^\eta P^2\right),\qquad \omega>0.
\]
On \(\partial\mathcal E_\eta\), its value is
\(H_\eta=\hbar\omega\), by the semiaxis equations. In the
quantum realization of Proposition~\ref{x_reciprocidad:prop:ii-covarianza-gaussiana},
\begin{equation}
 \langle\widehat H_\eta\rangle_{\psi_\eta}
 =\frac{\omega}{2}
 \left(e^{-\eta}\frac{\hbar e^\eta}{2}
       +e^\eta\frac{\hbar e^{-\eta}}{2}\right)
 =\frac{\hbar\omega}{2}.
 \label{x_reciprocidad:eq:ii-energia-contorno-covarianza}
\end{equation}
The two-standard-deviation contour therefore has twice the
mean energy of this Gaussian state. The areal law fixes action
normalization; the angular pair fixes its distribution between the two coordinates;
the canonical representation and explicit state realize that distribution
as covariance. Saturation of uncertainty is thus established in
the realization specified by its operators and its state.

% END OWNER


% BEGIN OWNER /Users/ruben/Documents/New project/output/EXTREMA_MEDIA_RAZON_RECIPROCIDAD_ES_EN_20260918/ARTICULO/ES/source/antecedentes_iii/sections/03_hojas_y_orientacion.tex
\chapter[Sheet return, areal--volumetric incidence and constitutive orientation]{Sheet return, areal--volumetric incidence\\and constitutive orientation}
\label{x_reciprocidad:iii:sec:hojas}
\label{x_reciprocidad:nuclear:constitutiva-antecedente}

The sheet structure enters before electromagnetic quantities
are evaluated. Its function is to distinguish closure of a projection,
restoration of an orientation and evolution of the transported register.
These operations act on related objects of different types: a
projected trajectory may have closed while its lift remains
on the opposite sheet; orientation may have been restored while memory
preserves the two windings performed. The constitutive construction must receive
this information together with the structure determining it.

\section{Projected return and restoration of the oriented cover}

Let $\ell$ be a return loop of a surviving extension, and let
$\varepsilon:\langle\ell\rangle\to\{+1,-1\}$ be the character of its
lift, with $\varepsilon(\ell)=-1$. The associated interval bundle
has the presentation
\begin{equation}
 \mathcal M=([0,1]\times[-1,1])/{(0,u)\sim(1,-u)}.
 \label{x_reciprocidad:iii:eq:mobius}
\end{equation}
The decisive datum is the sign of transport around the specified loop. The
identification of endpoints geometrically realizes that character.

\begin{proposition}[Two closure orders]
One winding of the loop closes the base and reverses the fibre orientation. Two
windings restore orientation. In the angular parametrization of the boundary
readout, these returns correspond, respectively, to the areal closure $2\pi$ and
the volumetric closure $4\pi$ of the oriented cover.
\end{proposition}
\begin{proof}
The equivalence relation of \eqref{x_reciprocidad:iii:eq:mobius} identifies the terminal fibre
with the initial one by $u\mapsto-u$. Composing two transports gives
$u\mapsto-(-u)=u$. Parametrizing one winding of the base by $2\pi$ assigns $4\pi$
to its square. The statement concerns this lift and this readout.
\end{proof}

TPK memory also preserves the trajectory traversed and the register
updates. The square of the sign character is the identity, but it does not impose
vanishing of these updates. The spinorial representation, when
used, must transport this datum through its own map.

\section{Descent of incidence from the tritic calendar}

In the primitive block $(+1,-1,0)$, the TPK updates the arithmetic mode through
$+1:m\mapsto\Sigma$, $-1:m\mapsto\Pi$ and $0:m\mapsto m$. Its cyclic orbit
is $(\Sigma,\Pi,\Pi)$. The sheet readout assigns $\Sigma$ to the areal sheet and
$\Pi$ to the volumetric sheet. In the basis order $(e_{\rm ar},e_{\rm vol})$,
edge incidence is represented by
\begin{equation}
 F_{\rm av}=\begin{pmatrix}0&1\\1&1\end{pmatrix}.
 \label{x_reciprocidad:iii:eq:fav}
\end{equation}
The three consecutive pairs are $\Sigma\to\Pi$, $\Pi\to\Pi$ and
$\Pi\to\Sigma$, each with multiplicity one. The calendars of lengths
$9$, $27$ and $108$ contain $3$, $9$ and $36$ copies of the primitive block.
Their incidence matrices are therefore these multiples of $F_{\rm av}$.

This construction preserves the distinction between chronological frequency and the measure
of the continuation space. The former assigns weights $(1/3,2/3)$ to the mode
cycle. The latter is defined on the minimal memory-one symbolic envelope
admitting exactly the preceding pairs. Its adjacency matrix is
$F_{\rm av}$; trajectory counts belong to that envelope.

\Needspace{10\baselineskip}
\begin{proposition}[Return weights of the two sheets]
Let $F_0=0$, $F_1=1$ and $F_{n+1}=F_n+F_{n-1}$. For $n\geq1$,
\begin{equation}
 F_{\rm av}^{n}=
 \begin{pmatrix}F_{n-1}&F_n\\F_n&F_{n+1}\end{pmatrix}.
 \label{x_reciprocidad:iii:eq:fav-powers}
\end{equation}
The Parry measure of this envelope assigns the sheets weights proportional
to $(1,\varphi^2)$, and the ratio of diagonal counts converges to $\varphi^{-2}$.
\end{proposition}
\begin{proof}
The power formula follows by induction using the recurrence.
The square of $F_{\rm av}$ is strictly positive. Its positive Perron
vector is $(1,\varphi)$, with eigenvalue $\varphi$, since
$\varphi^2=\varphi+1$. As the matrix is symmetric, the Parry weights are
proportional to the squares of that vector's components. The other root
of the characteristic polynomial is $-\varphi^{-1}$; its modulus is strictly
less than $\varphi$, giving the limit of the diagonal ratios.
\end{proof}

Spectral identification of this incidence follows its construction
from the calendar. The regional coordinate $\varphi$ retains its earlier coinductive
genealogy: recognition of the eigenvalue does not select the calendar.

\section{Boundary marking and areal--volumetric ratio}

Let $\Pi_{\rm HMT}(X_\infty)$ be the value of the regional closure readout already
constructed. With the diagonal sheet projectors, define
\begin{equation}
 D_\pi=\Pi_{\rm HMT}(X_\infty)P_{\rm ar}+P_{\rm vol},\qquad
 \Theta_n=
 \frac{\langle e_{\rm ar},D_\pi F_{\rm av}^{n}e_{\rm ar}\rangle}
 {\langle e_{\rm vol},F_{\rm av}^{n}e_{\rm vol}\rangle}.
 \label{x_reciprocidad:iii:eq:marked-return}
\end{equation}
Volumetric normalization is unity; regional closure marks the
areal return only once. Applying \eqref{x_reciprocidad:iii:eq:fav-powers},
\begin{equation}
 \Theta_n=\Pi_{\rm HMT}(X_\infty)\frac{F_{n-1}}{F_{n+1}}
 \longrightarrow\frac{\pi}{\varphi^2}.
 \label{x_reciprocidad:iii:eq:av-limit}
\end{equation}
The return ratio is determined by incidence, measure and marking.
The quotient $2\pi/4\pi$, by contrast, describes the cover's two closure orders.
Both relations participate in the sheet readout without replacing each other.

\section{Transport of orientation to the constitutive channels}

In a two-channel representation, let $\mathcal R$ be a self-adjoint
involution and $\mathcal S$ a unitary interchange operator satisfying
$\mathcal S\mathcal R\mathcal S^{-1}=-\mathcal R$. Write the angular coordinates
already constructed as $x=A_{\rm rad}$ and $y=C^*_{\rm rad}$, with $x>|y|$.
Then
\begin{equation}
 B=xI+y\mathcal R,\qquad T=e^{-B},\qquad
 \mathcal ST(y)\mathcal S^{-1}=T(-y).
 \label{x_reciprocidad:iii:eq:channel-conjugation}
\end{equation}
Indeed, conjugation transforms $B(y)$ into $B(-y)$ and commutes with its convergent
exponential series. The eigenvalues $q_\pm=e^{-x\mp y}$ are interchanged.
To fix the response order, the chamber $x>y>0$ is used, where
$0<q_+<q_-<1$. The map $(x,y)\mapsto(x,-y)$ transports it to the
opposite chamber $x>-y>0$, where $0<q_-<q_+<1$. Both belong to the
contractive domain $x>|y|$. On the boundary $y=0$ the two channels coincide.
Conjugation interchanges the chambers; it does not preserve the strict ordering of
their labels. It acts on the projectors through
$\mathcal SP_\pm\mathcal S^{-1}=P_\mp$, with
$P_\pm=(I\pm\mathcal R)/2$.
This concrete map transmits orientation to the constitutive operator
developed next.

% END OWNER


% BEGIN OWNER /Users/ruben/Documents/New project/output/EXTREMA_MEDIA_RAZON_RECIPROCIDAD_ES_EN_20260918/ARTICULO/ES/source/antecedentes_iii/sections/03b_vacancias_prefactor.tex
% Ampliación propia de III; no modifica el núcleo común copiado de I.
% Residencia propuesta: antes de Retorno de hoja, incidencia areal--volumétrica.
% Procedencia y controles: VACANCIAS_PREFACTOR_PROCEDENCIA.md, en esta carpeta.
\chapter{Discrete polarization and the tritic mode of arithmetic projections}
\label{x_reciprocidad:iii:sec:vacancias-prefactor}

Refinement and comparison between the APP sheets provide two
complementary observations. The first follows the accumulated discrepancy
between the number of transitions and the capacity acquired; the second preserves
the mode of the tritic selector in the pairing of the additive and
multiplicative fibres. Discrete polarization and the electronic prefactor
are obtained from these respective operations. Their articulation with the vacuum
requires preserving both the temporal information and the spaces on which
the operators act.

\section{Refinement discrepancy and endpoint convention}

The cardinalities of the ternary block and its decimal representation fix
\[
 \rho=\log_{1000}729=\log_{10}9,\qquad
 \delta=1-\rho=\log_{10}(10/9).
\]
We have $0<\delta<1$. Moreover, $\delta$ is irrational: an equality
$\delta=p/q$, with $q>0$, would imply $10^{q-p}=9^q$, which is incompatible
with prime factorization. This mismatch already enters the
encoding of refinement. Here its centred readout is made explicit,
preserving the phase origin and the endpoint convention.

\begin{definition}[Lower and upper encodings]
For $\theta\in\mathbb R$ and $n\in\mathbb Z$, define
\begin{align}
 z_n^{\downarrow}(\theta)
 &=\lfloor(n+1)\delta+\theta\rfloor-\lfloor n\delta+\theta\rfloor,
 \label{x_reciprocidad:iii:vp:lower}\\
 z_n^{\uparrow}(\theta)
 &=\lceil(n+1)\delta+\theta\rceil-\lceil n\delta+\theta\rceil.
 \label{x_reciprocidad:iii:vp:upper}
\end{align}
Both take values in $\{0,1\}$. The arrow distinguishes only the
endpoint convention for the intervals; tritic orientation retains
its earlier definition.
\end{definition}

The lower encoding coincides with the one used in the development
of refinement. The upper encoding recovers the formulation of
polarization using the ceiling function. The two expressions are related
by an exactly determined boundary term.

\begin{proposition}[Change of convention and boundary term]
Let $b_n(\theta)=\mathbf1_{\mathbb Z}(n\delta+\theta)$. Then
\begin{equation}
 z_n^{\uparrow}(\theta)-z_n^{\downarrow}(\theta)
 =b_n(\theta)-b_{n+1}(\theta).
 \label{x_reciprocidad:iii:vp:endpoint}
\end{equation}
In particular, for $N\geq0$,
\[
 \sum_{n=0}^{N-1}(z_n^{\uparrow}-z_n^{\downarrow})
 =b_0-b_N.
\]
\end{proposition}
\begin{proof}
For every real $x$, $\lceil x\rceil=\lfloor x\rfloor+1-
\mathbf1_{\mathbb Z}(x)$. Applying this identity to both endpoints
of each increment gives \eqref{x_reciprocidad:iii:vp:endpoint}; its sum telescopes.
The irrationality of $\delta$ implies that, for fixed $\theta$, at most
one integer index satisfies $n\delta+\theta\in\mathbb Z$.
\end{proof}

With $\theta=0$, the first upper symbol is $1$ and the lower one is $0$;
for $n\geq1$ they coincide. The initial term belongs to the register and is
preserved when the convention is changed. The densities and discrepancy
bounds are transported with that term, without modifying the TPK update
or restarting the history at each phase return.

\section{Bounded polarization and exact temporal balance}

\begin{definition}[Centred discrete polarization]
The polarization of the upper encoding, relative to the density
$\delta$, is
\begin{equation}
 P_N^{\uparrow}(\theta)
 =\sum_{n=0}^{N-1}\bigl(z_n^{\uparrow}(\theta)-\delta\bigr),
 \qquad P_0^{\uparrow}=0.
 \label{x_reciprocidad:iii:vp:polarization}
\end{equation}
Its temporal increment is
\begin{equation}
 J_N^{\uparrow}(\theta)
 =P_{N+1}^{\uparrow}(\theta)-P_N^{\uparrow}(\theta)
 =z_N^{\uparrow}(\theta)-\delta.
 \label{x_reciprocidad:iii:vp:current}
\end{equation}
\end{definition}

Centring separates uniform growth of density $\delta$ from
phase information. Thus, $P_N^{\uparrow}$ measures an accumulated imbalance,
not the total number of vacancies. The increment $J_N^{\uparrow}$ preserves
the local pulse and the background relative to which it is evaluated.

\begin{theorem}[Uniform bound and discrete temporal continuity]
For $N\geq0$ and every phase $\theta$,
\begin{equation}
 P_N^{\uparrow}(\theta)
 =\lceil N\delta+\theta\rceil-\lceil\theta\rceil-N\delta,
 \qquad |P_N^{\uparrow}(\theta)|<1.
 \label{x_reciprocidad:iii:vp:bound}
\end{equation}
For any integers $0\leq a<b$,
\begin{equation}
 \sum_{n=a}^{b-1}J_n^{\uparrow}(\theta)
 =P_b^{\uparrow}(\theta)-P_a^{\uparrow}(\theta),
 \qquad
 \left|\sum_{n=a}^{b-1}J_n^{\uparrow}(\theta)\right|<1.
 \label{x_reciprocidad:iii:vp:continuity}
\end{equation}
The density of the vacancy sequence is $\delta$.
\end{theorem}
\begin{proof}
Summing the ceiling increments in \eqref{x_reciprocidad:iii:vp:upper} yields the
first identity. Writing $r(x)=\lceil x\rceil-x\in[0,1)$
gives $P_N^{\uparrow}=r(N\delta+\theta)-r(\theta)$; its absolute
value is less than one. On the interval $[a,b)$, one similarly obtains
$r(b\delta+\theta)-r(a\delta+\theta)$, proving
the second bound, which is sharper than the sum of two individual bounds.
Finally,
$N^{-1}\sum_{n=0}^{N-1}z_n^{\uparrow}
=\delta+N^{-1}P_N^{\uparrow}\to\delta$.
\end{proof}

The lower encoding admits the same proofs with
$r(x)$ replaced by $\lfloor x\rfloor-x\in(-1,0]$. In particular,
$P_N^{\uparrow}-P_N^{\downarrow}=b_0-b_N$. The bounded discrepancy
previously used in the Dirichlet readouts and the present polarization
therefore preserve an explicit relationship of origin and endpoints.
Nor does aperiodicity introduce secular drift in this imbalance:
the history can extend indefinitely while the centred discrepancy
remains uniformly bounded.

\section{Realization on the charge and time lines}

To realize the balance dimensionally, specify a charge line
$\mathcal L_Q$ with a nonzero section $u_Q$ and an increasing sequence of
times $t_N$ on an oriented time line. With
$\Delta t_N=t_{N+1}-t_N>0$, the centred charge and the mean current of
each interval are defined by
\begin{equation}
 Q_N=P_N^{\uparrow}u_Q,\qquad
 I_N=\frac{Q_{N+1}-Q_N}{\Delta t_N}
     =\frac{u_Q}{\Delta t_N}\bigl(z_N^{\uparrow}-\delta\bigr).
 \label{x_reciprocidad:iii:vp:dimensional}
\end{equation}
Thus, $Q_N\in\mathcal L_Q$ and
$I_N\in\mathcal L_Q\otimes\mathcal L_T^{-1}$. The section $u_Q$ and
the time intervals are data of this realization; the preceding arithmetic
construction has already determined $P_N^{\uparrow}$ and
$J_N^{\uparrow}$.

\begin{proposition}[Preservation of balance under dimensional realization]
For $a<b$,
\begin{equation}
 Q_b-Q_a=\sum_{n=a}^{b-1}\Delta t_n I_n,
 \qquad
 \left|\frac{Q_b-Q_a}{u_Q}\right|<1.
 \label{x_reciprocidad:iii:vp:dimensional-balance}
\end{equation}
If all intervals have length $t_0$, the two possible currents
are $-\delta u_Q/t_0$ and $(1-\delta)u_Q/t_0$.
\end{proposition}
\begin{proof}
Multiplying each identity in \eqref{x_reciprocidad:iii:vp:dimensional} by
$\Delta t_n$ and summing telescopes the charges. The bound is
\eqref{x_reciprocidad:iii:vp:continuity}; the two local values follow from
$z_n^{\uparrow}\in\{0,1\}$.
\end{proof}

The continuity proved is temporal and refers to the centred charge.
A spatial realization also incorporates the localization of charges,
incidence between cells, and identification of boundary fluxes.
In a constitutive description by channels, the coupling is specified
through an excitation map
$\iota:\mathcal L_Q\to\mathcal H_{\rm ch}\otimes\mathcal L_Q$
and an observation map
$\ell:\mathcal H_{\rm ch}\otimes\mathcal L_Q\to\mathcal L_Q$,
where $\mathcal H_{\rm ch}$ is a real realization of the channel space.
If both are
linear and independent of the time index, the response of a fixed operator
$\mathcal V$ is transported by
\[
 Q_N^{\rm resp}=\ell(\mathcal V\otimes I_{\mathcal L_Q})\iota(Q_N),
 \qquad
 I_N^{\rm resp}=\frac{Q_{N+1}^{\rm resp}-Q_N^{\rm resp}}{\Delta t_N}.
\]
Linearity preserves the telescoping balance. This composition makes
explicit the data of a constitutive identification: the scalar sequence
determines the pulse, while $\iota$, $\mathcal V$, and $\ell$
determine its excitation, transport, and observation through channels. Identities
\eqref{x_reciprocidad:iii:vp:polarization}--\eqref{x_reciprocidad:iii:vp:dimensional-balance}
establish the first construction independently of that choice of
realization. The vacuum operator is constructed in
Section~\ref{x_reciprocidad:iii:sec:constitutiva} on its own channels and projectors.

\section{Pairing the additive and multiplicative fibres}

The second observation comes from the three-dimensional realization of APP,
$\mathcal A_3=D_9^3$, with uniform measure and
$D_9=\{1,\ldots,9\}$. The observables preserved are
\[
 \Sigma(x)=\rho_9(x_1+x_2+x_3),\qquad
 \Pi(x)=\rho_9(x_1x_2x_3),
\]
where $\rho_9(n)=1+[(n-1)\bmod9]$ for $n>0$. In
$\ell^2(\mathcal A_3)$, let $P_{\Sigma,3}$ and $P_{\Pi,3}$ be
the conditional expectations on their fibres. For example,
\[
 (P_{\Sigma,3}f)(x)=
 \frac1{|\Sigma^{-1}(\Sigma(x))|}
 \sum_{y:\,\Sigma(y)=\Sigma(x)}f(y).
\]
They are orthogonal projections defined on the same support. The incidence
matrix $N_{ab}=|\Sigma^{-1}(a)\cap\Pi^{-1}(b)|$ results from
enumerating the $729$ triples, adding and multiplying their coordinates before
applying $\rho_9$:
\begin{equation}
 N=9\begin{pmatrix}
 0&1&1&0&1&1&0&1&4\\
 1&0&1&1&0&1&1&0&4\\
 1&0&2&0&1&2&0&0&3\\
 0&1&1&0&1&1&0&1&4\\
 1&0&1&1&0&1&1&0&4\\
 0&0&2&1&0&2&0&1&3\\
 0&1&1&0&1&1&0&1&4\\
 1&0&1&1&0&1&1&0&4\\
 0&1&2&0&0&2&1&0&3
 \end{pmatrix}.
 \label{x_reciprocidad:iii:vp:joint-matrix}
\end{equation}
The rows sum to $81$, and the column sums are
\[
 (m_1,\ldots,m_9)=(36,36,108,36,36,108,36,36,297).
\]
Therefore, the pairing of the normalized indicator functions of the fibres
is $C_{ab}=N_{ab}/\sqrt{81m_b}$. Normalization incorporates the
multiplicities produced by APP, not a choice of singular values.

\begin{proposition}[Singular mode of the tritic selector]
In the canonical chart, the selector vector is
$m_{\rm ph}=(1,-1,0)^3$. We have
\begin{equation}
 Cm_{\rm ph}=C^{\mathsf T}m_{\rm ph}=-\tfrac12m_{\rm ph}.
 \label{x_reciprocidad:iii:vp:trit-mode}
\end{equation}
Thus, the corresponding singular value is $s=1/2$.
\end{proposition}
\begin{proof}
In the six columns where $m_{\rm ph}$ is nonzero, $m_b=36$ and
$C_{ab}=N_{ab}/54$. Multiplying the rows of
\eqref{x_reciprocidad:iii:vp:joint-matrix} by $(1,-1,0)^3$ yields
$-m_{\rm ph}/2$. For the transposed product, the signed sums
of columns $3,6,9$ are zero; the other six give the same
values $-m_{{\rm ph},b}/2$. This proves both identities and, by
composition, $CC^{\mathsf T}m_{\rm ph}=m_{\rm ph}/4$.
\end{proof}

\section{Commutator norm and electronic prefactor}

\begin{theorem}[Exact norm of arithmetic incompatibility]
The two projections satisfy
\begin{equation}
 \bigl\|[P_{\Sigma,3},P_{\Pi,3}]\bigr\|=\frac{\sqrt3}{4}.
 \label{x_reciprocidad:iii:vp:prefactor}
\end{equation}
The maximum is attained in the principal plane corresponding to the
tritic mode of \eqref{x_reciprocidad:iii:vp:trit-mode}.
\end{theorem}
\begin{proof}
The rational matrix $M=CC^{\mathsf T}$ has entries
$M_{ac}=\sum_bN_{ab}N_{cb}/(81m_b)$. Its complete spectrum is
verified through the following independent subspaces: the
vectors $(1,\ldots,1)$, $(1,-1,0)^3$, and $(1,1,-2)^3$ have
eigenvalues $1$, $1/4$, and $7/132$, respectively; the subspace of
vectors supported on $\{3,6,9\}$ whose sum is zero has dimension
$2$ and eigenvalue $1/18$; the zero-sum vectors supported on
$\{1,4,7\}$ or $\{2,5,8\}$ form a kernel of dimension $4$.
Substitution into the matrix verifies these actions, and the dimensions
sum to $9$. Consequently,
\[
 \operatorname{spec}M=
 \{1,1/4,7/132,1/18,1/18,0,0,0,0\}.
\]
For two orthogonal projections, a singular value $s\in(0,1)$ of the
pairing determines a principal plane in which their matrices are
\[
 P=\begin{pmatrix}1&0\\0&0\end{pmatrix},\qquad
 Q=\begin{pmatrix}s^2&s\sqrt{1-s^2}\\
 s\sqrt{1-s^2}&1-s^2\end{pmatrix}.
\]
This form is obtained by taking a unit vector $u$ in the first
image with $PQu=s^2u$ and completing it with
$(Qu-s^2u)/(s\sqrt{1-s^2})$. Their commutator is
$s\sqrt{1-s^2}\bigl(\begin{smallmatrix}0&1\\-1&0\end{smallmatrix}\bigr)$.
The sectors with singular value $0$ or $1$ contribute zero.
Therefore, the total norm is the maximum of
$\sqrt{\lambda(1-\lambda)}$ over the spectrum of $M$. That maximum
is $\sqrt{3/16}$, attained at $\lambda=1/4$.
\end{proof}

The norm is the prefactor entering the electronic composition.
Its derivation preserves the three-dimensional support, the fibres, and the
selector mode; scalar evaluation does not replace that structure. In
particular, knowing $\sqrt3/4$ gives the maximum magnitude of
the incompatibility, while the projectors preserve the directions
in which it acts.

\clearpage
\section{Quadratic defect and response of the 90/120 sectors}

The tritic mode simultaneously determines
\begin{equation}
 1-s^2=\frac34=\frac{90}{120},\qquad s=\frac12.
 \label{x_reciprocidad:iii:vp:defect-ratio}
\end{equation}
The first expression is the quadratic defect of the APP principal plane.
The last preserves the lengths of the two calendar sectors.
The constitutive response uses those sectors on another object: a
contractive channel $q\in(0,1)$. Its evaluation is
\[
 r(q)=\frac{1-q^{90}}{1-q^{120}}.
\]
The relationship between the sector ratio and this response can be
made precise without identifying their operators.

\begin{proposition}[Limit of the response ratio]
Para $0<q<1$,
\begin{equation}
 \frac34<r(q)<1,\qquad
 \lim_{q\uparrow1}r(q)=\frac34=1-s^2.
 \label{x_reciprocidad:iii:vp:response-limit}
\end{equation}
\end{proposition}
\begin{proof}
With $u=q^{30}\in(0,1)$, factoring $1-u^3$ and
$1-u^4$ gives
$r(q)=(1+u+u^2)/(1+u+u^2+u^3)$. The limit at $u=1$ is
$3/4$. Moreover,
$4(1+u+u^2)-3(1+u+u^2+u^3)
=(1-u)(3u^2+2u+1)>0$, and the denominator exceeds the numerator
by $u^3>0$. This yields both inequalities.
\end{proof}

The exact coincidence concerns the defect of the singular mode and the limit
of the sector response. In the contractive interior, the response
depends on $q$ and is strictly greater than $3/4$. The constitutive
development preserves its two channels $q_\pm$ and its projectors
$P_\pm$; the projectors $P_{\Sigma,3}$ and $P_{\Pi,3}$, for their
part, preserve the arithmetic fibres that determined the prefactor.
An operator equivalence between the two realizations requires a map
between their spaces that intertwines the corresponding actions.
Equations \eqref{x_reciprocidad:iii:vp:trit-mode}, \eqref{x_reciprocidad:iii:vp:prefactor}, and
\eqref{x_reciprocidad:iii:vp:response-limit} specify the shared information and
keep the data of each realization visible.

% END OWNER


% BEGIN OWNER /Users/ruben/Documents/New project/output/EXTREMA_MEDIA_RAZON_RECIPROCIDAD_ES_EN_20260918/ARTICULO/ES/source/antecedentes_iii/sections/04_respuesta_constitutiva.tex
\chapter{Constitutive operator and reconstruction of the vacuum channels}
\label{x_reciprocidad:iii:sec:constitutiva}

The constitutive quantities are obtained by an operation on angular
transport. The starting information comprises the channels, their
projectors, and the orientation distinguishing them. This organization makes it possible
to follow separately the construction of the operator, the evaluation of its
spectral functions, and its subsequent dimensional realization.

\section{Positive response on two sheets}

In the oriented chamber $x>y>0$ of Section~\ref{x_reciprocidad:iii:sec:hojas}, let
$P_\pm=(I\pm\mathcal R)/2$ and
$T=q_+P_++q_-P_-$, with $0<q_+<q_-<1$. The incidences constructed by the
calendar fix the sectors $90$ and $120$. The constitutive response is
\begin{equation}
 \mathcal V=(I-T^{90})(I-T^{120})^{-1}.
 \label{x_reciprocidad:iii:eq:vacuum-op}
\end{equation}
This definition receives the transport arising from the TPK and preserves the
sheet assignments. The response rule and its domain are an explicit part
of the construction.

The quotient arises from comparing the two calendar incidences on
the same transport, not from fitting four quantities separately. If
$S=T^{30}$ and $\Sigma_j(S)=I+S+\cdots+S^{j-1}$, factoring
$I-S^j$ gives
\begin{equation}
 \mathcal V\,\Sigma_4(S)=\Sigma_3(S),\qquad
 \mathcal V=\Sigma_3(S)\Sigma_4(S)^{-1}.
 \label{x_reciprocidad:iii:eq:completion-response}
\end{equation}
The factors $\Sigma_3$ and $\Sigma_4$ express, respectively, the
sectors $90=3\cdot30$ and $120=4\cdot30$. Since the spectrum of $S$ is
contained in $(0,1)$, $\Sigma_4(S)$ is invertible: the completion
equation determines a unique response for the fixed transport.
This uniqueness concerns the declared constitutive rule and its
incidences. Electromagnetic identification adds the oriented
assignment of sheets and the dimensional lines specified
below.

\begin{proposition}[Decomposition and positivity]
The response is well-defined and
\begin{equation}
 \mathcal V=r_+P_++r_-P_-,\qquad
 r_\pm=\frac{1-q_\pm^{90}}{1-q_\pm^{120}}.
 \label{x_reciprocidad:iii:eq:rchannels}
\end{equation}
Its eigenvalues satisfy $3/4<r_-<r_+<1$ in this chamber.
\end{proposition}
\begin{proof}
Each coefficient $1-q_\pm^{120}$ is positive. The inverse of $I-T^{120}$
is the sum of the projectors divided by those coefficients, proving
\eqref{x_reciprocidad:iii:eq:rchannels}. With $s=q^{30}$, one obtains
\begin{equation}
 f(s)=\frac{1+s+s^2}{1+s+s^2+s^3},\qquad
 f'(s)=-\frac{s^2(3+2s+s^2)}{(1+s+s^2+s^3)^2}<0.
 \label{x_reciprocidad:iii:eq:monotone}
\end{equation}
The limits at $0$ and $1$ are $1$ and $3/4$, respectively. Monotonicity
and the channel ordering complete the proof.
\end{proof}

\section{Four coordinates of a common response}

Define the normalized evaluations
\begin{equation}
 \widehat\varepsilon=r_+^2,\qquad
 \widehat\mu=r_-^2,\qquad
 \widehat Z=\frac{r_-}{r_+},\qquad
 \widehat c=\frac1{r_+r_-}.
 \label{x_reciprocidad:iii:eq:four}
\end{equation}
The first pair preserves the quadratic responses assigned to each sheet.
The product captures the common mode; the quotient distinguishes their relative orientation.
The following hold exactly:
\begin{equation}
 \widehat\mu\widehat\varepsilon=\widehat c^{-2},\qquad
 \frac{\widehat\mu}{\widehat\varepsilon}=\widehat Z^2,
 \qquad
 r_+=\frac1{\sqrt{\widehat Z\widehat c}},\quad
 r_-=\sqrt{\frac{\widehat Z}{\widehat c}}.
 \label{x_reciprocidad:iii:eq:four-identities}
\end{equation}
The positive roots fix the inversion. Under sheet conjugation,
read relative to the fixed markings $P_\pm$,
$\widehat\varepsilon$ and $\widehat\mu$ are interchanged,
$\widehat Z$ is inverted, and $\widehat c$ is preserved. The response is then transported
to the chamber $x>-y>0$: the formulas remain valid, with the orderings
$q_-<q_+$ and $r_+<r_-$ reversed. Consequently, common mode
and orientation are data recoverable through complementary observables.

\begin{proposition}[Determination of the quadratic evaluations]
Let $W=r_+P_++r_-P_-$ be a positive response on the two sheets.
With the product and quotient readouts fixed,
\[
 \widehat c^{-1}=\det W=r_+r_-,
 \qquad \widehat Z=r_-/r_+,
\]
there exists a unique positive pair $(\widehat\varepsilon,\widehat\mu)$
satisfying
$\widehat\mu\widehat\varepsilon=\widehat c^{-2}$ and
$\widehat\mu/\widehat\varepsilon=\widehat Z^2$.
It is precisely $(r_+^2,r_-^2)$.
\end{proposition}
\begin{proof}
Positivity permits taking roots unambiguously:
\[
 \widehat\mu=\frac{\widehat Z}{\widehat c}
 =\frac{r_-}{r_+}r_+r_-=r_-^2,\qquad
 \widehat\varepsilon=\frac1{\widehat Z\widehat c}
 =\frac{r_+}{r_-}r_+r_-=r_+^2.
\]
These expressions satisfy both relations and are their only positive
roots.
\end{proof}

Thus the squares are not two independent choices added
to the channels: they are determined by the multiplicative and oriented
readouts once the constitutive relations are fixed. This proposition
establishes uniqueness within that system of rules. Interpreting
the rules as an electromagnetic response retains its additional
physical-representation content.

The determinantal readout of transport and the constitutive response
must also preserve their order. For $T=e^{-xI-y\mathcal R}$, one has
$\det T=e^{-2x}$, whereas
\[
 \det\mathcal V=f(q_+^{30})f(q_-^{30})=\widehat c^{-1}.
\]
Both arise from the same two-sheet operator but evaluate different
functions of its spectrum. In general, applying the determinant first
and then a rational function is not equivalent to taking the determinant of
that function of the operator. Keeping both eigenvalues until the
response is completed preserves orientation information that an isolated
determinant no longer distinguishes.

\section{Cubic inversion and angular recovery}

\begin{theorem}[Constitutive inversion in the contractive domain]
For each $r\in(3/4,1)$, there exists a unique $s\in(0,1)$ such that $f(s)=r$.
This $s$ is the unique root in that interval of
\begin{equation}
 r s^3+(r-1)(1+s+s^2)=0.
 \label{x_reciprocidad:iii:eq:cubic}
\end{equation}
The labeled response reconstructs $s_\pm$, $q_\pm=s_\pm^{1/30}$, and
\begin{equation}
 x=-\frac12\log(q_+q_-),\qquad
 y=\frac12\log\frac{q_-}{q_+}.
 \label{x_reciprocidad:iii:eq:recover-angular}
\end{equation}
\end{theorem}
\begin{proof}
The derivative and limits in \eqref{x_reciprocidad:iii:eq:monotone} prove that $f$ is a
decreasing bijection between the stated intervals. Multiplying $f(s)=r$
by its denominator produces \eqref{x_reciprocidad:iii:eq:cubic}. The positive root of order
$30$ recovers the channel. Finally, adding and subtracting
$\log q_\pm=-x\mp y$ gives \eqref{x_reciprocidad:iii:eq:recover-angular}.
\end{proof}

Reconstruction preserves the sheet labels. The complete operator
$\mathcal V$ also preserves the projectors, and the inverse function is applied
by spectral calculus. Recovering two angular coordinates is an
operation defined on this representation; the complete historical
TPK trace preserves additional data.

The interval $r\in(3/4,1)$ and the inverse function of this theorem belong
to the internal normalization of \eqref{x_reciprocidad:iii:eq:four}. If coordinates
expressed in another units chart are available, this normalization
is first recovered through the transformation in
equation~\eqref{x_reciprocidad:iii:eq:dimensions}; cubic inversion is then applied.

\section{Logarithmic coordinates and refinement}

Let $\mathcal E=\log(r_+r_-)$ and $\mathcal B=\log(r_-/r_+)$. Then
\begin{equation}
 \log\mathcal V=\tfrac12\mathcal E I-\tfrac12\mathcal B\mathcal R.
 \label{x_reciprocidad:iii:eq:logvac}
\end{equation}
This decomposition displays the invariant and oriented components.
In the amplification $\mathcal V_m=\mathcal V\otimes I_{9^m}$, normalized
partial expectations satisfy
$E_m(\mathcal V_{m+1})=\mathcal V_m$. For every polynomial $p$,
$p(\mathcal V_m)=p(\mathcal V)\otimes I_{9^m}$; hence
$E_m(p(\mathcal V_{m+1}))=p(\mathcal V_m)$. Uniform continuity
extends the identity to continuous functions on the spectrum.
Thus projectors, powers, and the logarithm are preserved in the stated tower.

\section{Dimensional types of the four evaluations}

With the action, charge, and velocity lines, equipped with sections
$u_S,u_Q,u_C$, the dimensional realization reads
\begin{align}
 Z_{\rm vac}&=\widehat Z u_Su_Q^{-2},&
 c_{\rm vac}&=\widehat c u_C,\\
 \mu_{\rm vac}&=\widehat\mu u_Su_C^{-1}u_Q^{-2},&
 \varepsilon_{\rm vac}&=\widehat\varepsilon u_S^{-1}u_C^{-1}u_Q^2.
 \label{x_reciprocidad:iii:eq:dimensions}
\end{align}
Multiplying and dividing these expressions preserves
$\mu_{\rm vac}\varepsilon_{\rm vac}=c_{\rm vac}^{-2}$ and
$\mu_{\rm vac}/\varepsilon_{\rm vac}=Z_{\rm vac}^2$.
The SI chart subsequently evaluates the sections. Physical comparison requires
presenting their construction and normalization alongside the evaluation; the four
normalized coordinates of \eqref{x_reciprocidad:iii:eq:four} retain their dimensionless type.

% END OWNER


% BEGIN OWNER /Users/ruben/Documents/New project/output/EXTREMA_MEDIA_RAZON_RECIPROCIDAD_ES_EN_20260918/ARTICULO/ES/source/antecedentes_iii/sections/04b_reconstruccion_forma_lc.tex
% INSERCIÓN FOCAL AUTORIZADA: no introduce un capítulo ni el concepto de radión.
% Incluir dentro de la sección constitutiva, después de la inversión cúbica
% y de la declaración de las cartas dimensionales. Etiquetas fuera de grupos
% de prefijado automático. Conserva íntegramente el bloque geométrico anterior.
\section{Reconstruction of the action shape from the constitutive response}
\label{x_reciprocidad:iii:subsec:puente-forma-LC}

The constitutive response and the inductive--capacitive realization receive
the same angular pair constructed from APP--TRIT--TPK. The former evaluates
rational functions of the two channels; the latter uses the ratio of
their exponents to determine the shape of the action ellipse. Cubic
inversion allows the two readouts to be composed without identifying their
numerical ratios.

Let $x=A_{\rm rad}$ and $y=C^*_{\rm rad}$, with $x>|y|$, be the lifted
coordinates already constructed, and preserve their sheet labels. Write
\begin{equation}
 q_\pm=e^{-(x\pm y)},\qquad s_\pm=q_\pm^{30},\qquad
 r_\pm=f(s_\pm),\qquad
 f(s)=\frac{1+s+s^2}{1+s+s^2+s^3}.
 \label{x_reciprocidad:iii:eq:forma-canales}
\end{equation}
The contractive domain gives $s_\pm\in(0,1)$ and
$r_\pm\in(3/4,1)$. In the chamber $x>y>0$, we have
$s_+<s_-$ and $r_+>r_-$; sheet interchange transports this chamber
to the opposite orientation.

The image of the two normalized constitutive coordinates is the domain
\begin{equation}
 \mathcal U=\left\{(z,v)\in(0,\infty)^2:
   (zv)^{-1/2}\in(3/4,1),\quad
   (z/v)^{1/2}\in(3/4,1)\right\}.
 \label{x_reciprocidad:iii:eq:forma-dominio}
\end{equation}
Here $(z,v)=(\widehat Z,\widehat c)$. The restrictions express
membership of the reconstructed channels in the domain of the inversion
theorem; they do not constitute a subsequent choice of their values.

\Needspace{23\baselineskip}
\begin{proposition}[Constitutive composition of the shape and LC impedance]
\label{x_reciprocidad:iii:prop:forma-LC}
For $(\widehat Z,\widehat c)\in\mathcal U$, let $s_\pm\in(0,1)$ be the
unique roots of
\begin{equation}
 r_\pm s_\pm^3+(r_\pm-1)(1+s_\pm+s_\pm^2)=0,
 \qquad
 r_+=\frac1{\sqrt{\widehat Z\widehat c}},\quad
 r_-=\sqrt{\frac{\widehat Z}{\widehat c}}.
 \label{x_reciprocidad:iii:eq:forma-inversion}
\end{equation}
The angular ratio, the anisotropy of the action ellipse, and the relative LC
impedance are reconstructed exactly by
\begin{align}
 \frac{y}{x}
 &=\frac{\log s_+-\log s_-}{\log s_++\log s_-},
 \label{x_reciprocidad:iii:eq:forma-razon}\\
 \eta_{\rm el}
 &=\frac12\log\frac{\log s_+}{\log s_-},
 \label{x_reciprocidad:iii:eq:forma-eta}\\
 \frac{Z_{\rm LC}}{Z_*}
 &=e^{-\eta_{\rm el}}
  =\sqrt{\frac{\log s_-}{\log s_+}}.
 \label{x_reciprocidad:iii:eq:forma-LC}
\end{align}
In the last identity, $Z_{\rm LC}=Z_{\eta_{\rm el}}$ is the impedance
of the LC chart of \eqref{x_reciprocidad:hmtII:eq:ii-elipse-lc}.
\end{proposition}
\begin{proof}
The cubic inversion in \eqref{x_reciprocidad:iii:eq:cubic} uniquely determines $s_\pm$,
because $f$ is a decreasing bijection from $(0,1)$ onto
$(3/4,1)$. Their positive roots of order $30$ recover $q_\pm$.
The identities
\[
 \log s_+=-30(x+y),\qquad \log s_-=-30(x-y)
\]
prove \eqref{x_reciprocidad:iii:eq:forma-razon} by difference and sum. Both logarithms
are negative; their ratio is positive and their sum is nonzero. Therefore,
\[
 \frac12\log\frac{\log s_+}{\log s_-}
   =\frac12\log\frac{x+y}{x-y}
   =\operatorname{artanh}(y/x)=\eta_{\rm el}.
\]
The relation $Z_{\eta_{\rm el}}/Z_*=e^{-\eta_{\rm el}}$, proved in
\eqref{x_reciprocidad:hmtII:eq:ii-elipse-impedancia}, gives \eqref{x_reciprocidad:iii:eq:forma-LC}.
\end{proof}

The composition preserves the information needed for each reconstruction.
With $\hbar$, one recovers
$a_S=\sqrt{2\hbar}\,e^{\eta_{\rm el}/2}$ and
$b_S=\sqrt{2\hbar}\,e^{-\eta_{\rm el}/2}$; with the chart
$(\omega,Z_*)$, one recovers the inductance and capacitance already defined.
Action, clock and dimensional reference accompany the form as
typed data. Inverting two dimensionless coordinates does not
reconstruct their units or the entire memory of the state.

Normalization enters before solving the cubic equations.
If $Z_{\rm vac}$ and $c_{\rm vac}$ are received in a dimensional chart,
the relations \eqref{x_reciprocidad:iii:eq:dimensions} first determine
\[
 \widehat Z=\frac{Z_{\rm vac}}{u_Su_Q^{-2}},\qquad
 \widehat c=\frac{c_{\rm vac}}{u_C}.
\]
These are the coordinates to which
\eqref{x_reciprocidad:iii:eq:forma-inversion} applies. The reference $Z_*$ belongs to the
LC realization and is explicitly preserved in \eqref{x_reciprocidad:iii:eq:forma-LC}.

The two ratios are related through the channels, but evaluate
different functions:
\begin{equation}
 \widehat Z=\frac{f(s_-)}{f(s_+)},\qquad
 \frac{Z_{\rm LC}}{Z_*}
       =\sqrt{\frac{\log s_-}{\log s_+}}.
 \label{x_reciprocidad:iii:eq:forma-dos-cocientes}
\end{equation}
The preceding reconstruction provides the map between the two readouts.
In particular, interchanging the sheets inverts both ratios, changes
$\eta_{\rm el}$ to $-\eta_{\rm el}$, and preserves $\widehat c$.
In the symmetric limit $y=0$, the channels coincide and both ratios
equal $1$. Outside that case, a common origin does not impose equality
of the two functions. The action shape is recoverable from the
constitutive response while preserving orientation and normalization.

% END OWNER


% BEGIN OWNER /Users/ruben/Documents/New project/output/EXTREMA_MEDIA_RAZON_RECIPROCIDAD_ES_EN_20260918/ARTICULO/ES/source/antecedentes_iii/sections/04c_velocidad_y_reloj.tex
% RESULTADO_RECUPERADO: funtor dimensional integral, eq:cZ-internos,
% eq:longitud-ruta; composición reunida en Artículo V REV02, sección 69.
% Incorporación autónoma a III: mismo transporte, mismas hojas y misma carta.
\section{Constitutive speed and kinematic realization of the clock}
\label{x_reciprocidad:iii:sec:enlace-velocidad-constitutiva}

The constitutive response and the kinematic reader realize the same
velocity section. Their identification preserves the order of construction:
the APP sheets, their TRIT orientation and the angular transport published
by TPK precede the velocity readings and the choice of units.
In the two-sheet representation of
section~\ref{x_reciprocidad:iii:sec:hojas}, with $x>|y|$ and rank-one projectors
$P_\pm$, the transport is
\[
 T=e^{-xI-y\mathcal R}=q_+P_++q_-P_-,
 \qquad q_\pm=e^{-x\mp y},\qquad 0<q_\pm<1.
\]
Its temporal blocks of lengths $90$ and $120$ produce the response
$\mathcal V$ in \eqref{x_reciprocidad:iii:eq:vacuum-op}. The coordinates $x,y$ and the
channels are outputs of the preceding angular development; this composition
applies its readouts to that already constructed transport.

The symmetric logarithmic readout of temporal transport is
\begin{equation}
 \mathcal E=
 \log\!\bigl[(1-q_+^{90})(1-q_-^{90})\bigr]
 -\log\!\bigl[(1-q_+^{120})(1-q_-^{120})\bigr].
 \label{x_reciprocidad:iii:eq:velocidad-lector-simetrico}
\end{equation}
On the dimensional velocity line, with positive section $u_C$, the
kinematic reader publishes $c_{\mathrm{int}}=\exp(-\mathcal E)u_C$.
The constitutive response publishes
$c_{\mathrm{vac}}=\widehat c u_C$ through \eqref{x_reciprocidad:iii:eq:four}.

\begin{proposition}[Constitutive and kinematic identity]
\label{x_reciprocidad:iii:prop:identidad-velocidad}
In the same dimensional chart, we have
\begin{equation}
 c_{\mathrm{int}}=c_{\mathrm{vac}}=\widehat c\,u_C,
 \qquad \widehat c=(r_+r_-)^{-1}=\exp(-\mathcal E).
 \label{x_reciprocidad:iii:eq:identidad-velocidad}
\end{equation}
The identification preserves sheet interchange.
\end{proposition}
\begin{proof}
All factors in \eqref{x_reciprocidad:iii:eq:velocidad-lector-simetrico} are
positive. The logarithm law and \eqref{x_reciprocidad:iii:eq:rchannels} give
\[
 \mathcal E=\log(r_+r_-)=\log\det\mathcal V,
 \qquad e^{-\mathcal E}=(\det\mathcal V)^{-1}.
\]
The determinant corresponds to the rank-one two-sheet realization.
The two definitions therefore publish the same section. Exchanging
the sheets permutes $r_+$ and $r_-$ and preserves their product. The determinant
$\det T=e^{-2x}$ corresponds to elementary transport, whereas
$\det\mathcal V=r_+r_-$ corresponds to its temporal response;
\eqref{x_reciprocidad:iii:eq:identidad-velocidad} uses the latter.
\end{proof}

\subsection{Path length and elementary duration}

The positive dimensional frame $(u_S,u_C,u_T,u_Q)$ induces the length
basis $u_L=u_Cu_T$. For an enriched path $\gamma$ of
$n=|\gamma|$ transitions and a constant channel, the additive readers are
\begin{equation}
 T_{\mathrm{int}}(\gamma)=nu_T,
 \qquad L_{\mathrm{int}}(\gamma)=n\widehat c\,u_L.
 \label{x_reciprocidad:iii:eq:velocidad-longitud-reloj}
\end{equation}
Concatenation adds the numbers of transitions and their realizations.
Enriched routes also preserve their orientation and memory.

\begin{corollary}[Normalization of the elementary length]
For $n>0$, $L_{\mathrm{int}}(\gamma)/T_{\mathrm{int}}(\gamma)
=c_{\mathrm{vac}}$. In the normalization $t_0=u_T$,
\begin{equation}
 \ell_0=c_{\mathrm{int}}t_0=\widehat c\,u_L.
 \label{x_reciprocidad:iii:eq:longitud-elemental-constitutiva}
\end{equation}
In a temporal chart $t_0=\tau_0u_T$, with $\tau_0>0$, we obtain
$\ell_0=\tau_0\widehat c\,u_L$.
\end{corollary}
\begin{proof}
Dividing the two expressions in \eqref{x_reciprocidad:iii:eq:velocidad-longitud-reloj}
cancels $n$ and uses $u_L/u_T=u_C$. Multiplying the speed section
by $t_0$ proves both elementary formulas.
\end{proof}

When the channel depends on the edge, the readout retains the additive form
$L_{\mathrm{int}}(\gamma)=\sum_{e\in\gamma}\widehat c(e)u_L$.
Its quotient by $nu_T$ is the mean of the speeds of those edges.
The constant-channel expression is the homogeneous specialization of
this readout.

\subsection{Change of chart and adapted normalization}

\begin{proposition}[Covariance of the identification]
Let $u_C'=d u_C$ and $u_T'=\eta u_T$, with $d,\eta>0$.
The coordinates of the same sections satisfy
\begin{equation}
 \widehat c'=d^{-1}\widehat c,
 \qquad \tau_0'=\eta^{-1}\tau_0,
 \qquad u_L'=d\eta u_L.
 \label{x_reciprocidad:iii:eq:velocidad-cambio-carta}
\end{equation}
If another publication writes $c'=\widehat c' u_C'$ with $u_C'=d u_C$,
equality of sections is exactly equivalent to $d\widehat c'=\widehat c$.
\end{proposition}
\begin{proof}
Equality of each section in the two bases gives
$\widehat c'u_C'=\widehat c u_C$ and $\tau_0'u_T'=\tau_0u_T$.
In particular,
\[
 \tau_0'\widehat c'u_L'
 =\frac{\tau_0}{\eta}\frac{\widehat c}{d}\,d\eta u_L
 =\tau_0\widehat c u_L=\ell_0.
\]
The final equivalence follows from the same substitution for speed.
\end{proof}

The adapted choice $d=\widehat c$ gives $u_C'=c_{\mathrm{vac}}$ and
$\widehat c'=1$: it expresses the constructed section in an adapted basis.
The coefficient of the internal chart remains $(r_+r_-)^{-1}$, and this
is the normalization used by cubic inversion. Keeping the
action and charge bases, \eqref{x_reciprocidad:iii:eq:dimensions} produces
\[
 \widehat\mu'=d\widehat\mu,\qquad
 \widehat\varepsilon'=d\widehat\varepsilon,\qquad
 \widehat\mu'\widehat\varepsilon'=(\widehat c')^{-2}.
\]
In the adapted chart, we have $\widehat\mu'=r_-/r_+$ and
$\widehat\varepsilon'=r_+/r_-$, by substitution of
$\widehat\mu=r_-^2$ and $\widehat\varepsilon=r_+^2$.

\subsection{Preservation of the readout under amplification}

For the refinement $\mathcal V_m=\mathcal V\otimes I_{9^m}$, the
normalized trace $\tau_m=\operatorname{Tr}/(2\cdot9^m)$ satisfies
\begin{equation}
 \exp[-2\tau_m(\log\mathcal V_m)]
 =\exp[-\log r_+-\log r_-]=\widehat c.
 \label{x_reciprocidad:iii:eq:velocidad-traza-normalizada}
\end{equation}
Indeed, each eigenvalue appears $9^m$ times; that multiplicity
cancels against the denominator of the normalized trace. The ordinary
amplified determinant is $(r_+r_-)^{9^m}$. The two-sheet normalization,
equivalently expressed by \eqref{x_reciprocidad:iii:eq:velocidad-traza-normalizada},
preserves the constitutive readout during amplification. This
realization preserves the enriched transport from which it originates.

% END OWNER


% BEGIN OWNER /Users/ruben/Documents/New project/output/EXTREMA_MEDIA_RAZON_RECIPROCIDAD_ES_EN_20260918/ARTICULO/ES/source/antecedentes_iii/sections/07_carga_y_cuantos.tex
\chapter{Charge section and electrical quantization relations}
\label{x_reciprocidad:iii:sec:carga}

The constitutive response establishes a relation between the two sheets and a
relative propagation scale. The charge section is obtained by composing
this response with the fine-structure constant and the action, both
constructed in the preceding chapters. This order has a precise
consequence: charge does not retrospectively determine the dodecaphase record.
It is the record, together with the action and the vacuum response, that
determines the charge section on the corresponding line.

\section{Positive determination of the charge section}

Consider the positive sections $\hbar_a$, with
$a\in\{\mathrm{pre},\mathrm{ret}\}$, and $h_a=2\pi\hbar_a$.
The impedance $Z_{\rm vac}$ has already been constructed through
\eqref{x_reciprocidad:iii:eq:dimensions}. The quotient $h_a/Z_{\rm vac}$ belongs to
the quadratic charge line. This dimensional compatibility allows
the following definition to be formulated without yet choosing an SI chart.

\begin{definition}[Positive charge section]
The charge section corresponding to action orientation $a$ is
\begin{equation}
 e_a=\left(\frac{2\alpha h_a}{Z_{\rm vac}}\right)^{1/2}.
 \label{x_reciprocidad:iii:eq:charge}
\end{equation}
The root is taken on the positive charge half-line. Its conjugate section
is $-e_a$; the sign of the latter is not identified with the action
return index.
\end{definition}

\begin{proposition}[Existence, uniqueness, and compatibility of the coupling]
The equation $Z_{\rm vac}e_a^2=2\alpha h_a$ determines a unique
positive section. In the constitutive realization, one also has
\begin{equation}
 \alpha=\frac{Z_{\rm vac}e_a^2}{2h_a}
 =\frac{e_a^2}{4\pi\varepsilon_{\rm vac}\hbar_a c_{\rm vac}}.
 \label{x_reciprocidad:iii:eq:coupling}
\end{equation}
\end{proposition}
\begin{proof}
Positivity of $\alpha$, $h_a$, and $Z_{\rm vac}$ guarantees a unique
positive root. The constitutive identities give
$\varepsilon_{\rm vac}c_{\rm vac}=Z_{\rm vac}^{-1}$.
Substituting this equality and $h_a=2\pi\hbar_a$ into the first expression
of \eqref{x_reciprocidad:iii:eq:coupling} yields the second.
\end{proof}

The last expression is the usual form of electromagnetic coupling
in a physical chart. Here it appears as a subsequent identity of the constructed
sections. The proof does not recalculate $\alpha$ from a measured
charge: it establishes the algebraic consistency of the preceding result with the
section just defined. Physical identification additionally requires
retaining the chart of units and the coupling regime used in the
comparison. An identity between quantities and an experimental comparison
have different functions in that recognition.

\section{Resistance, conductance, flux, and frequency}

The following four quantities use the same section $e_a$. Their
dependencies can be fully resolved before any decimal is evaluated.

\begin{theorem}[Composition of the electrical quanta]
The sections
\begin{align}
 R_K&=\frac{h_a}{e_a^2}=\frac{Z_{\rm vac}}{2\alpha},
 &G_0&=\frac{2e_a^2}{h_a}=\frac{4\alpha}{Z_{\rm vac}},
 \label{x_reciprocidad:iii:eq:rk-g0}\\
 \Phi_{0,a}&=\frac{h_a}{2e_a}
 =\sqrt{\frac{h_aZ_{\rm vac}}{8\alpha}},
 &K_{J,a}&=\frac{2e_a}{h_a}=\Phi_{0,a}^{-1}
 \label{x_reciprocidad:iii:eq:flux-josephson}
\end{align}
satisfy exactly
\begin{equation}
 R_KG_0=2,\qquad G_0Z_{\rm vac}=4\alpha,\qquad
 K_{J,a}\Phi_{0,a}=1,\qquad
 K_{J,a}^{2}R_K=\frac4{h_a}.
 \label{x_reciprocidad:iii:eq:quanta-identities}
\end{equation}
\end{theorem}
\begin{proof}
Substituting $e_a^2=2\alpha h_a/Z_{\rm vac}$ into the first two
definitions cancels $h_a$ and gives \eqref{x_reciprocidad:iii:eq:rk-g0}. Squaring
the positive section $h_a/(2e_a)$ gives
$h_aZ_{\rm vac}/(8\alpha)$; positivity fixes the root in
\eqref{x_reciprocidad:iii:eq:flux-josephson}. Its inverse is $2e_a/h_a$.
The first three products in \eqref{x_reciprocidad:iii:eq:quanta-identities}
reduce, respectively, to $2$, $4\alpha$, and $1$. Finally,
$(4e_a^2/h_a^2)(h_a/e_a^2)=4/h_a$.
\end{proof}

The physical realization recognizes in these sections the von
Klitzing resistance, the conductance quantum, the flux quantum, and the
Josephson constant. The factor $2$ in the conductance belongs to the normalization
of the quantity defined here; identifying it with a multiplicity of
channels in a material system requires specifying that representation.
The four equations do not constitute four independent generators:
they share $\alpha$, the action, and the vacuum response. Precisely this
dependency produces the cross-identities and makes it possible to check an
evaluation without altering the operations from which it arose.

\section{Action return and transport of the sections}

The quotient $R_{\rm act}=h_{\rm pre}/h_{\rm ret}$, obtained before the
angular pair, preserves the distinction between the two action sections.
It does not change the sheet assignments of the constitutive response. These are
two distinct operations: action return and channel interchange.

\begin{proposition}[Covariance under return]
With common $\alpha$ and $Z_{\rm vac}$, one has
\begin{align}
 \frac{e_{\rm pre}}{e_{\rm ret}}&=R_{\rm act}^{1/2},
 &R_K^{\rm pre}&=R_K^{\rm ret},
 &G_0^{\rm pre}&=G_0^{\rm ret},\label{x_reciprocidad:iii:eq:charge-return}\\
 \frac{\Phi_{0,\rm pre}}{\Phi_{0,\rm ret}}&=R_{\rm act}^{1/2},
 &\frac{K_{J,\rm pre}}{K_{J,\rm ret}}&=R_{\rm act}^{-1/2}.
 \label{x_reciprocidad:iii:eq:flux-return}
\end{align}
\end{proposition}
\begin{proof}
The first identity follows by dividing the two versions of
\eqref{x_reciprocidad:iii:eq:charge} and taking the positive root. In
\eqref{x_reciprocidad:iii:eq:rk-g0}, the action has already cancelled. The other two
identities follow from the powers $h_a^{1/2}$ and $h_a^{-1/2}$ in
\eqref{x_reciprocidad:iii:eq:flux-josephson}.
\end{proof}

Resistance and conductance preserve the constitutive quotient, whereas
flux and frequency also retain the action orientation. This
separation gives a structural reading of their dependencies: some quantities
are invariants of the sheet quotient, while others additionally measure the
action section used. The metrological comparison must specify which section it evaluates;
the numerical proximity of the two actions does not authorize their identification.

% END OWNER


% BEGIN OWNER /Users/ruben/Documents/New project/output/EXTREMA_MEDIA_RAZON_RECIPROCIDAD_ES_EN_20260918/ARTICULO/ES/source/antecedentes_vii/sections/v_funcional_barbero.tex
% RESULTADO_RECUPERADO desde II REV09. Véase MAPA_PROCEDENCIA.json.
% Selección de líneas y cambios mecánicos explícitos; ninguna prueba nueva.
\chapter{Hexadic--octadic incidence and the Barbero--Immirzi functional}
\label{x_reciprocidad:v:ant:sec:barbero}
Let $H\hookrightarrow O$ be the marked inclusion of a hexad in an octad,
$|H|=6$, $|O|=8$, constructed through the binary residue of
Section~\ref{x_reciprocidad:exc:residuo-binario}. The hexad, dodecad
and regional chart determining that inclusion are preserved.
The incidence objects
\[
\mathcal F_6=H\times\binom H2,\qquad
\mathcal F_8=O\times\binom H2
\]
have cardinalities $90$ and $120$. Their normalized difference gives
\[
\frac{|\mathcal F_6|-|\mathcal F_8|}{24\binom62}=-\frac1{12}.
\]
\begin{figure}[htbp]
\centering
\begin{tikzpicture}[>=Latex,every node/.style={font=\small}]
\draw[rounded corners,draw=ink] (-1.8,-1.35) rectangle (1.8,1.35);
\foreach \a in {0,60,...,300}{
\fill[ink] (\a:0.83) circle(2pt);}
\node at (0,-1.7) {$H:\ 6$ puntos};
\draw[->,thick,ink] (2.1,0)--(3.6,0);
\node[above] at (2.85,1.5) {marked inclusion};
\draw[rounded corners,draw=ink] (4,-1.35) rectangle (7.6,1.35);
\foreach \a in {0,60,...,300}{
\fill[ink] ($(5.8,0)+(\a:0.83)$) circle(2pt);}
\fill[accent] (4.5,0) circle(2pt);
\fill[accent] (7.1,0) circle(2pt);
\node at (5.8,-1.7) {$O:\ 8$ puntos};
\node at (0,-2.35) {$6\binom62=90$};
\node at (5.8,-2.35) {$8\binom62=120$};
\end{tikzpicture}
\caption{Marked incidence from six to eight points. The drawing represents the sets and inclusion; it does not attempt to turn the faces or vertices of a cube into Steiner blocks without their incidence map.}
\end{figure}

\section{Dodecaphase fundamental chain and pole coefficient}
In the free module of oriented events, fix the fundamental
chain of the revolution
\[
c_{12}^{+}=
\sum_{j\in\Z/12\Z}[j,\mathcal F_6]
-[\partial_{\rm ret},\mathcal F_8].
\]
Multiplicity twelve corresponds to the sectors and multiplicity one to the oriented seam. The involution reverses orientation of the complete chain. Let
\[
S_m(q)=\frac{q^m}{1-q^{3m}},\qquad 0<q<1.
\]
The linear reader assigns $S_{90}$ to every sector and $S_{120}$ to the seam; consequently,
\[
\mathscr R(c_{12}^{+})=12S_{90}-S_{120}.
\]

\begin{definition}[Oriented evaluation of hexadic--octadic incidence]
\label{x_reciprocidad:paper:barbero-definicion}
The parameter \(\gamma_{\HMT}\) is the conjugate evaluation of the chain
\(c_{12}^{+}\), completed by its angular component. On the domain
\(0<q_\pm<1\), it is defined by
\[
\gamma_{\HMT}=\frac{\pi AC^*}{180}
+[\mathscr R(c_{12}^{+})](q_-)-[\mathscr R(c_{12}^{+})](q_+).
\]
The marked incidence determines the event spaces; the oriented
chain fixes the multiplicities, and the channels of
Section~\ref{x_reciprocidad:iv:angular:inversion} determine the evaluation.
\end{definition}

\begin{theorem}[Pole coefficient and conjugate evaluation]
The image of the fundamental chain has coefficient $1/24$ relative to $(1-q)^{-1}$ and residue $-1/24$ relative to $(q-1)^{-1}$. Its evaluation on the previously generated channels determines
\begin{equation}
\gamma_{\HMT}=
\frac{\pi AC^*}{180}
+12\bigl[S_{90}(q_-)-S_{90}(q_+)\bigr]
-\bigl[S_{120}(q_-)-S_{120}(q_+)\bigr].
\label{x_reciprocidad:v:ant:eq:gamma}
\end{equation}
\end{theorem}
\begin{proof}
The factorization $1-q^{3m}=(1-q)\sum_{j=0}^{3m-1}q^j$ implies
$\lim_{q\uparrow1}(1-q)S_m(q)=1/(3m)$. By linearity,
\[
\frac{12}{270}-\frac1{360}=\frac1{24}.
\]
The coordinate $q-1$ reverses the sign of the residue.
The conjugate evaluation and the angular component of Definition~\ref{x_reciprocidad:paper:barbero-definicion}
yield \eqref{x_reciprocidad:v:ant:eq:gamma}.
\end{proof}

The Diophantine certificate $4a-3b=45$ admits the pairs $(12+3t,1+4t)$, $t\ge0$; it certifies minimality of $(12,1)$ but does not replace the origin of multiplicities in the chain. The evaluation of \eqref{x_reciprocidad:v:ant:eq:gamma} is
\[
\gamma_{\HMT}
=0.27400767269445927474725526077398041940\ldots.
\]
This numerical value is an output of the oriented functional, rather than an external value
used to select $A$, $C^*$, or their coefficients.

\section{Parity and structural sensitivity}
Reversing $C^*$ exchanges $q_+$ and $q_-$ and changes the sign of the angular component. Therefore
\[
\gamma_{\HMT}(A,-C^*)=-\gamma_{\HMT}(A,C^*).
\]
The functional preserves orientation, whereas other readers, such as $q_+q_-$, are even. This difference precedes its physical interpretation.

% END OWNER


% BEGIN OWNER /Users/ruben/Documents/New project/output/EXTREMA_MEDIA_RAZON_RECIPROCIDAD_ES_EN_20260918/ARTICULO/ES/source/antecedentes_vii/sections/v_unidad_areal.tex
% RESULTADO_RECUPERADO desde II REV09. Véase MAPA_PROCEDENCIA.json.
% Selección de líneas y cambios mecánicos explícitos; ninguna prueba nueva.
\chapter{Action, gravity, and thermal information}
\label{x_reciprocidad:v:ant:sec:gravity}
\section{Dodecaphasic period and circular realization of action}
The already constructed calendar determines $T_{108}=108t_0$.
The constitutive speed, duration and trajectory length are
obtained through the identities of
Proposition~\ref{x_reciprocidad:iii:prop:identidad-velocidad} and
equations~\eqref{x_reciprocidad:iii:eq:velocidad-longitud-reloj} and
\eqref{x_reciprocidad:iii:eq:longitud-elemental-constitutiva}.
In its circular realization,
\[
\omega_{108}=\frac{2\pi}{108t_0},\qquad
R_{108}=\frac{c_{\rm int}}{\omega_{108}}
=\frac{54}{\pi}\ell_0,\qquad \ell_0=c_{\rm int}t_0.
\]
The unit $\ell_0$ and the time and velocity charts remain
explicit; the SI chart does not select the cycle.
For each oriented action section $a\in\{\mathrm{pre},\mathrm{ret}\}$,
\[
E_{108,a}=\hbar_a\omega_{108},\qquad
m_{108,a}=\frac{E_{108,a}}{c_{\rm int}^2}.
\]
The reduced and full lengths are, respectively,
\[
\frac{\hbar_a}{m_{108,a}c_{\rm int}}=R_{108},
\qquad
\frac{h_a}{m_{108,a}c_{\rm int}}=2\pi R_{108}.
\]
The same action and duration enter into both readings.

\begin{proposition}[Radius-coincidence selector]
In the dimensional family
\[
G_a(\vartheta)=\vartheta\frac{c_{\rm int}^5t_0^2}{\hbar_a},
\qquad \vartheta>0,
\]
the circular-realization condition
$G_a(\vartheta)m_{108,a}/c_{\rm int}^2=R_{108}$
has the unique solution
\[
\vartheta^\circ_{108}=(54/\pi)^2.
\]
\end{proposition}
\begin{proof}
The first radius is
$r_g(\vartheta)=\vartheta(\pi/54)\ell_0$.
Comparing it with $(54/\pi)\ell_0$, with $\ell_0>0$, reduces the
equality to a strictly increasing linear equation.
\end{proof}
The coincidence condition is an explicit part of this realization;
the theorem solves for its coefficient without using a metrological value of $G$.
The radius used here is $Gm/c^2$, so the factor two from
the convention $2Gm/c^2$ is not inadvertently transferred.

\begin{proposition}[Reciprocal action and gravity, conserved area]
In the two preceding charts,
\[
\frac{G_{\rm pre}}{G_{\rm ret}}=R_{\rm act}^{-1},\qquad
\hbar_aG_a=\vartheta^\circ_{108}c_{\rm int}^5t_0^2,
\qquad
\ell_P^2=\frac{\hbar_aG_a}{c_{\rm int}^3}
=\vartheta^\circ_{108}\ell_0^2.
\]
\end{proposition}
\begin{proof}
The definition of $G_a$ contains $\hbar_a^{-1}$ and the remaining factors
are common to the two charts. Taking the quotient and the product gives
the three equalities.
\end{proof}
This reciprocity explains why orientation information can
vary without changing the area. A precise relation is thus available
between the reciprocal action sections, the Planck scale and the areal operator
of Section~\ref{x_reciprocidad:v:ant:sec:area}.

% END OWNER


% BEGIN OWNER /Users/ruben/Documents/New project/output/EXTREMA_MEDIA_RAZON_RECIPROCIDAD_ES_EN_20260918/ARTICULO/ES/source/antecedentes_vii/sections/v_area_informacion.tex
% RESULTADO_RECUPERADO desde II REV09. Véase MAPA_PROCEDENCIA.json.
% Selección de líneas y cambios mecánicos explícitos; ninguna prueba nueva.
\chapter{Area operator and reconstruction of sectorial information}
\label{x_reciprocidad:v:ant:sec:area}
Regrouping six trits constructs a bijection between
$(\Z/9\Z)^3$ and $(\Z/27\Z)^2$, both of cardinality $729$.
The bijection is not presented as an additive isomorphism. On a
$9\times9$ boundary block, the source decomposes 36 links into two
sets of 18, labeled by channels $90$ and $120$.

\section{Tritic regrouping and boundary incidence}
\label{x_reciprocidad:v:ant:sec:ii-area-incidencia-explicita}

The correspondence preserves an ordered word, in addition to its cardinality.
Write each coordinate \(x_j\in\mathbb Z/9\mathbb Z\), for
\(j=1,2,3\), through its unique positional representatives
\(x_j=r_{2j-1}+3r_{2j}\), with \(r_k\in\{0,1,2\}\). The map is
\begin{equation}
 \beta_{729}(x_1,x_2,x_3)
 =\bigl(r_1+3r_2+9r_3,\ r_4+3r_5+9r_6\bigr)
 \in(\mathbb Z/27\mathbb Z)^2.
 \label{x_reciprocidad:v:ant:eq:ii-area-beta729}
\end{equation}

\begin{proposition}[Positional correspondence between interior and boundary]
\label{x_reciprocidad:v:ant:prop:ii-area-beta729}
The map \(\beta_{729}\) is bijective and preserves the order of the
six trits. Its inverse extracts the two triples of digits and regroups them
into three pairs.
\end{proposition}
\begin{proof}
The positional expansion of the representatives \(0,\ldots,8\) and
\(0,\ldots,26\) is unique. Extracting the six digits, regrouping them, and
extracting them again returns the same word. The inverse operation
thus restores each of the three initial coordinates.
\end{proof}

Carry preserves its position in the enriched word. The additive
operations of the two packagings have different boundaries: for example,
\(\beta_{729}(0,2,0)=(18,0)\) and
\(\beta_{729}(0,1,0)=(9,0)\), whose sum in the order-twenty-seven
torus is \((0,0)\), whereas
\(\beta_{729}(0,3,0)=(0,1)\). This distinction allows the positional
correspondence to be reused while keeping the operations typed.

In the torus \((\mathbb Z/27\mathbb Z)^2\), the block
\(B_\partial=\{0,\ldots,8\}^2\) has the oriented link sets
\begin{align}
 \partial B_{90}
 &=\{((-1,j),(0,j)),\ ((8,j),(9,j)):0\le j\le8\},\\
 \partial B_{120}
 &=\{((i,-1),(i,0)),\ ((i,8),(i,9)):0\le i\le8\}.
 \label{x_reciprocidad:v:ant:eq:ii-area-enlaces}
\end{align}
All coordinates are read modulo twenty-seven. Each system contains eighteen links and the two systems are disjoint. Their labels preserve the incidence channels; the cardinality of the links is \(18+18\).

\begin{theorem}[Minimum capacity of the cubic cut]
\label{x_reciprocidad:v:ant:thm:ii-area-corte-minimo}
In the cubic graph with vertices \(\{0,\ldots,N-1\}^3\), with
\(N\ge2\), links between neighbors, and unit capacity per link,
the minimum cut between two opposite faces has capacity \(N^2\).
For \(N=9\), a maximally mixed product state on the
trits of a minimum layer has information
\begin{equation}
 \operatorname{cap}_{\min}=81,\qquad s_{\rm corte}=81\log3.
 \label{x_reciprocidad:v:ant:eq:ii-area-corte-informacion}
\end{equation}
\end{theorem}
\begin{proof}
The \(N^2\) lines parallel to the normal of the faces form edge-disjoint
paths. Every separating cut intercepts each path, so
it contains at least \(N^2\) links. The links crossing
a transverse layer constitute a cut of that cardinality. Each uniform trit
contributes \(\log3\) nats; entropy is additive for the product
state of the eighty-one trits of the layer.
\end{proof}

The capacity of the cubic cut and the state of the thirty-six links
in \eqref{x_reciprocidad:v:ant:eq:ii-area-enlaces} belong to observables defined on
different cuts. The genealogy preserves both domains and their measures.

\section{Transport of the collective incidence modes}
\label{x_reciprocidad:v:ant:sec:ii-area-transporte-colectivo}

The spaces \(\mathcal F_6\) and \(\mathcal F_8\) of
Section~\ref{x_reciprocidad:v:ant:sec:barbero} provide the uniform vectors
\[
 u_{90}=\frac1{\sqrt{90}}\sum_{f\in\mathcal F_6}(\delta_f,0),
 \qquad
 u_{120}=\frac1{\sqrt{120}}\sum_{f\in\mathcal F_8}(0,\delta_f)
\]
in \(\ell^2(\mathcal F_6)\oplus\ell^2(\mathcal F_8)\).
In \(\ell^2(\partial B_{90})\oplus\ell^2(\partial B_{120})\),
construct
\[
 v_{90}=\frac1{\sqrt{18}}\sum_{e\in\partial B_{90}}(\delta_e,0),
 \qquad
 v_{120}=\frac1{\sqrt{18}}\sum_{e\in\partial B_{120}}(0,\delta_e).
\]
The two pairs are orthonormal. Denote their respective spanned planes
by \(\mathscr H_{\rm inc}^{\rm coll}\) and
\(\mathscr H_\partial^{\rm coll}\).

\begin{theorem}[Collective intertwining of the sectorial labels]
\label{x_reciprocidad:v:ant:thm:ii-area-entrelazamiento-colectivo}
There is a unique surjective linear isometry between these planes such that
\begin{equation}
 \mathcal J_{6\to8}u_m=v_m\quad(m=90,120).
 \label{x_reciprocidad:v:ant:eq:ii-area-j-colectivo}
\end{equation}
For the operators
\[
 \begin{aligned}
 D_{\rm inc}&=90|u_{90}\rangle\langle u_{90}|
             +120|u_{120}\rangle\langle u_{120}|,\\
 D_\partial&=90|v_{90}\rangle\langle v_{90}|
             +120|v_{120}\rangle\langle v_{120}|,
 \end{aligned}
\]
we have
\begin{equation}
 \mathcal J_{6\to8}D_{\rm inc}
 =D_\partial\mathcal J_{6\to8}.
 \label{x_reciprocidad:v:ant:eq:ii-area-j-entrelaza}
\end{equation}
The sectorial projectors and their oriented combination with
coefficients \(12:-1\) are also transported.
\end{theorem}
\begin{proof}
The image of a basis fixes the unique linear map. Norms and
inner products are preserved because both bases are orthonormal.
Identity \eqref{x_reciprocidad:v:ant:eq:ii-area-j-entrelaza} holds on each vector
\(u_m\), since its two sides are \(m v_m\). The projectors are
expressed as \((120I-D)/30\) and \((D-90I)/30\); their linear functional
calculus and any combination of the two preserve intertwining.
\end{proof}

This transport acts on the two uniform modes and their spectral
labels. Individual incidences and fluctuations within
each sector remain in their complete spaces. The link preserves
exactly the collective datum that the areal weighting will use.

\section{Angular character and area operator}
\label{x_reciprocidad:v:ant:sec:ii-area-caracter}

Let $N_e=\operatorname{diag}(0,1,2)$ on each link, and let
$N_{90},N_{120}$ be the sector sums. With the scale
\[
\theta_{\rm clk}=\frac{C^*}{6}\frac{\pi}{180},\qquad
a_0=\frac{1+2\cos(10^\circ)}3\theta_{\rm clk},
\]
fixed, the angular factor is the normalized real character of the triplet
\begin{equation}
 \chi_{10}=\frac13\Tr\operatorname{diag}
       (1,e^{i\pi/18},e^{-i\pi/18})
 =\frac{1+2\cos(10^\circ)}3,
 \qquad a_0=\chi_{10}\theta_{\rm clk}.
 \label{x_reciprocidad:v:ant:eq:ii-area-caracter-triplete}
\end{equation}
The sum of the conjugate phases proves the equality. On the positive branch
considered, \(\theta_{\rm clk}>0\), \(\chi_{10}>0\), and
\(\gamma_{\HMT}>0\); hence \(a_0\) and the areal scale are
positive. The angular triplet and its normalization precede the
entropic distribution on the links. With these quantities,
the normalized area operator is
\[
\widehat{\mathcal A}/\ell_P^2
=\gamma_{\HMT}a_0(N_{90}+N_{120}).
\]
Its spectrum has values $n\gamma_{\HMT}a_0$, $0\le n\le72$,
with multiplicities given by $(1+z+z^2)^{36}$.
The proof consists in summing the eigenvalues $0,1,2$ of the
36 tensor factors.

\begin{theorem}[Recovery of the boundary sectors]
Let
\[
N=N_{90}+N_{120},\qquad D=90N_{90}+120N_{120}.
\]
Then
\begin{equation}
N_{90}=4N-\frac D{30},\qquad
N_{120}=\frac D{30}-3N.
\label{x_reciprocidad:v:ant:eq:recover}
\end{equation}
Thus the total and the weighted reader jointly recover
the two occupations.
\end{theorem}
\begin{proof}
La matriz $\left(\begin{smallmatrix}1&1\\90&120\end{smallmatrix}\right)$
has determinant $30$. Its inversion gives \eqref{x_reciprocidad:v:ant:eq:recover}.
On the domain of integer occupations, the image satisfies
$D\in30\Z$ and $90N\le D\le120N$.
\end{proof}

The modular Hamiltonian of the source is
$K_A=-\log\rho_A=cI+\Delta_{\rm mod}D$.
For $\Delta_{\rm mod}\ne0$ and known normalization,
$D=(K_A-cI)/\Delta_{\rm mod}$ publishes the second observable.
This operator describes the spectral structure of the reduced state;
it is not identified by notation with the generator of time evolution.
The area preserves the total occupation and the pair adds the provenance.
For example, $(N_{90},N_{120})=(1,0)$ and $(0,1)$ share $N=1$,
but have $D=90$ and $D=120$.


% BEGIN OWNER /Users/ruben/Documents/New project/output/EXTREMA_MEDIA_RAZON_RECIPROCIDAD_ES_EN_20260918/ARTICULO/ES/source/antecedentes_vii/sections/v_estado_modular.tex
% RESULTADO_RECUPERADO desde II REV09. Véase MAPA_PROCEDENCIA.json.
% Selección de líneas y cambios mecánicos explícitos; ninguna prueba nueva.
\section{State normalization and the meaning of the modular Hamiltonian}
The second observable is made explicit through the state that defines it.
The preceding tensor decomposition is retained: eighteen qutrits
per sector, each with number operator $\operatorname{diag}(0,1,2)$;
$N_m$ is the sum of these operators and the sectors act on distinct
factors. Let
\[
 z_m=1+e^{-m\Delta_{\rm mod}}+e^{-2m\Delta_{\rm mod}},
 \qquad m\in\{90,120\}.
\]
The family of source states has the product realization
\[
 \rho_A=Z^{-1}e^{-\Delta_{\rm mod}(90N_{90}+120N_{120})},
 \qquad Z=z_{90}^{18}z_{120}^{18}.
\]
Commutativity of the number operators and tensor factorization
prove the expression for $Z$. The state has full rank for
real, finite $\Delta_{\rm mod}$, and its spectral logarithm gives
\[
 K_A=-\log\rho_A=(\log Z)I+\Delta_{\rm mod}D.
\]
For $\Delta_{\rm mod}\ne0$, $K_A$ weights the excitations
of the two sectors differently. Joint reading with $N$ recovers the two occupations
by equation~\eqref{x_reciprocidad:v:ant:eq:recover}.

The need for joint reading is seen in two complementary examples.
The occupation states $(1,0)$ and $(0,1)$ have the same total $N=1$
and different modular weights. Conversely, $(4,0)$ and $(0,3)$ have
the same $D=360$ and different totals. The pair $(N,D)$ distinguishes both
cases. This recovery concerns sectoral occupations, not
each microscopic vector within a degeneracy of those occupations.

When $\Delta_{\rm mod}=0$, the state is proportional to the identity
and $K_A$ is scalar: that reading ceases to distinguish the sectors.
For $\Delta_{\rm mod}\ne0$, the known normalization allows one to recover
$D$ and compose it with the area. The role of the modular Hamiltonian
is thereby defined: it characterizes the reduced state and its information distribution.
Its inclusion in the article is motivated by that mathematical function,
whereas the evolution Hamiltonian belongs to the dynamical development
of the action.

% END OWNER


\section{Reduced state, Schmidt decomposition, and entanglement entropy}
If $\rho_A=\sum_\nu p_\nu|\nu\rangle\langle\nu|$,
the purification
\[
|\Psi\rangle=\sum_\nu\sqrt{p_\nu}\,
|\nu\rangle_A|\nu\rangle_B
\]
satisfies $\Tr_B|\Psi\rangle\langle\Psi|=\rho_A$.
Thus $-\Tr(\rho_A\log\rho_A)$ is the entanglement entropy
of this particular pure bipartition. The enlarged state declares
which system carries that interpretation.

For a two-sheet involution with one eigenvector of each sign,
\[
\rho_y=\frac{e^{-yR}}{2\cosh y},\qquad R^2=I,
\]
diagonalization gives
\[
S(\rho_y)=\log(2\cosh y)-y\tanh y,\qquad
D(\rho_y\Vert\rho_{-y})=2y\tanh y.
\]
Entropy is even in $y$; the distance explicitly compares
the two orientations. This two-level entropy is not identified
with the capacity $81\log3$ of the triadic cut.

% END OWNER


% BEGIN OWNER /Users/ruben/Documents/New project/output/EXTREMA_MEDIA_RAZON_RECIPROCIDAD_ES_EN_20260918/ARTICULO/ES/source/antecedentes_vii/sections/v_transduccion_termica.tex
% RESULTADO_RECUPERADO desde II REV09. Véase MAPA_PROCEDENCIA.json.
% Selección de líneas y cambios mecánicos explícitos; ninguna prueba nueva.
% Fragmento aditivo. No sustituye 10b_boltzmann.tex.
% Procedencia y distinción entre recuperación y formalización:
% 10c_boltzmann_procedencia.md. Etiquetas locales prefijadas ii-kb-aut.

\chapter{Genealogical memory, informational reduction and thermal transduction}
\label{x_reciprocidad:v:ant:sec:ii-kb-aut-desarrollo}

The thermal construction uses two publications of the same
APP--TRIT--TPK development. The first preserves the states, admissible routes and
register information; the second provides the oriented action
and the duration of the return. Their composition determines an energy per
tritic decision. The temperature and energy coordinates are subsequently
introduced through explicit dimensional maps. This sequence
allows one to identify what information is preserved, what information a reader
publishes and what quantity converts that information into dimensional entropy.

\section{Reader, memory, and state recovery}

Fix a finite level of the development and denote by \(\Omega\) the
set of admissible enriched states. The register preserves sheet,
orientation, residue, quotient, carry, route, phase and memory in the
domains where each coordinate enters. Let
\[
 q:\Omega\longrightarrow Y,\qquad
 M:\Omega\longrightarrow\mathcal M
\]
be the visible publication and the retained memory, respectively. For a
distribution \(p\) over \(\Omega\), we write
\(\mathsf H(p)=-\sum_xp_x\log p_x\), with \(0\log0=0\).
The logarithm is applied to dimensionless probabilities.

\begin{definition}[Genealogical recovery of a publication]
The pair \(\widehat q=(q,M)\) is a recoverable reader on the support
\(\Omega_p\) when there exists
\[
 d:\widehat q(\Omega_p)\longrightarrow\Omega_p,
 \qquad d\circ\widehat q=\operatorname{id}_{\Omega_p}.
\]
The fibre information of \(q\) is
\(\mathsf H_{\mathrm{fib}}(q;p)=\mathsf H(X\mid q(X))\), where
\(X\) has law \(p\).
\end{definition}

\begin{proposition}[Recovery identity and cost of reduction]
For a recoverable reader,
\begin{align}
 \mathsf H(X\mid q(X),M(X))&=0,\\
 \mathsf H(X)&=\mathsf H(q(X))+\mathsf H(X\mid q(X)),\\
 \mathsf H(X\mid q(X))&=\mathsf H(M(X)\mid q(X)).
 \label{x_reciprocidad:v:ant:eq:ii-kb-aut-fibra}
\end{align}
\end{proposition}
\begin{proof}
The map \(d\) recovers \(X\) from the pair \((q(X),M(X))\),
so the first conditional entropy is zero. Since \(q\) is
a function of \(X\), the chain rule gives the second equality.
For the third, expand
\(\mathsf H(X,M(X)\mid q(X))\) in two ways. One gives
\(\mathsf H(X\mid q(X))\), since \(M\) is a function of \(X\);
the other gives
\(\mathsf H(M(X)\mid q(X))+\mathsf H(X\mid q(X),M(X))\).
\end{proof}

Euclidean division on the APP sheets displays the elementary
arithmetic mechanism. With \(i,j\in\{1,\ldots,9\}\), one preserves
\[
 i+j=9q^+_{ij}+R^+_{ij},\qquad
 ij=9q^\times_{ij}+R^\times_{ij},\qquad
 1\leq R^\sigma_{ij}\leq9\quad(\sigma\in\{+,\times\}).
\]
The residue--quotient pair recovers the integer value of the operation.
Summing the differences over the eighty-one positions gives
\[
 2025-810=54+9\cdot129.
\]
Quotient memory restores the contribution not published by
residual reduction. Recoverability of the complete state additionally
requires preserving the relevant position and route coordinates;
the residue--quotient pair identifies the result of the operation here.

\section{Reversible update, tritic decision and erasure}

Let \(U:\Omega\to\Omega\) be the TPK update on a domain where
enriched transport has an inverse. The transported law is
\(p'_y=p_{U^{-1}y}\). The bijection preserves the multiset of
probabilities and therefore
\begin{equation}
 \mathsf H(U(X))=\mathsf H(X),\qquad
 \mathsf H(X\mid U(X))=0.
 \label{x_reciprocidad:v:ant:eq:ii-kb-aut-reversible}
\end{equation}
If the observed output is \(q(U(X))\), the unpublished information
is quantified by \(\mathsf H(X\mid q(U(X)))\). The complete state,
its publication and the resetting of a register are thus three distinct maps.

For the three oriented TRIT states, let
\(p=(p_{-1},p_0,p_{+1})\) be a normalized distribution. Its information is
\begin{equation}
 I_{\mathrm{TRIT}}(p)=-\sum_{a=-1}^{1}p_a\log p_a,
 \qquad 0\leq I_{\mathrm{TRIT}}(p)\leq\log3.
 \label{x_reciprocidad:v:ant:eq:ii-kb-aut-trit}
\end{equation}
The upper bound corresponds to \(p_a=1/3\). Indeed,
\[
 D(p\Vert(1/3,1/3,1/3))
 =\sum_ap_a\log(3p_a)=\log3-I_{\mathrm{TRIT}}(p)\geq0;
\]
the inequality follows from the convexity of \(t\log t\), and
equality characterizes the uniform distribution. The lower bound
corresponds to a distribution concentrated on one state.

The ideal Landauer realization uses a memory with energetically
degenerate logical states, an isothermal environment at temperature \(\Theta\)
and a balance that preserves the output of the logical map and accounts for
auxiliary registers. In that model, the decrease in logical entropy is
\(k_B[\mathsf H(X)-\mathsf H(F(X))]
=k_B\mathsf H(X\mid F(X))\) for a deterministic map \(F\).
The entropic balance requires transferring at least that amount to the environment;
its energy expression is \(Q\geq k_B\Theta\mathsf H(X\mid F(X))\).
The distinction between reversible logical evolution and physical erasure
corresponds to the Landauer--Bennett realization \cite{x_reciprocidad:bennett2003}.
The invertible update
\eqref{x_reciprocidad:v:ant:eq:ii-kb-aut-reversible} has zero logical contribution;
resetting a uniform triple requires transferring \(\log3\) nats
to the environment and satisfies \(Q\geq k_B\Theta\log3\).
The control energy and dissipation of a specific physical
realization belong to its additional thermodynamic balance.

\section{Action, frequency, and decision energy}

Let \(\mathcal L_{\rm act}\) and \(\mathcal L_T\) be the dimensional
lines of action and duration; the energy line is
\(\mathcal L_E=\mathcal L_{\rm act}\otimes\mathcal L_T^*\).
The preceding development provides the positive sections
\(\hbar_a\in\mathcal L_{\rm act}\), with
\(a\in\{\mathrm{pre},\mathrm{ret}\}\), and the elementary duration
\(t_0\in\mathcal L_T\). The complete calendar provides
\[
 T_{108}=108t_0,\qquad
 \omega_{108}=\frac{2\pi}{T_{108}},\qquad
 h_a=2\pi\hbar_a.
\]
The dimensional contraction between action and frequency defines
\begin{equation}
 \epsilon_a:=\hbar_a\omega_{108}
 =\frac{h_a}{108t_0}
 =\frac{\hbar_a}{t_P},\qquad
 t_P=\frac{54}{\pi}t_0.
 \label{x_reciprocidad:v:ant:eq:ii-kb-aut-energia}
\end{equation}
The final expression uses the same reduced circular period as the
gravitational section. The three expressions correspond to a single
energy section, with an orientation \(a\) maintained in all
the factors. The normalized information of a uniform decision
produces the energy section per nat
\begin{equation}
 \varepsilon_a:=\frac{\epsilon_a}{\log3}\in\mathcal L_E^+.
 \label{x_reciprocidad:v:ant:eq:ii-kb-aut-energia-nat}
\end{equation}
The previously constructed action determines the energy scale once
the duration is fixed; the cardinality of the TRIT alternatives determines its
informational normalization. The value of \(\log3\) enters through the
count and measure declared in \eqref{x_reciprocidad:v:ant:eq:ii-kb-aut-trit}.

\section{Dimensional definition of Boltzmann and selection of coordinates}

Let \(\mathcal L_\Theta\) be the oriented temperature line.
The dimensional entropy space is
\(\mathcal L_{\rm ent}=\mathcal L_E\otimes\mathcal L_\Theta^*\).
A positive linear map
\(k_B:\mathcal L_\Theta\to\mathcal L_E\) is, equivalently,
a positive element of \(\mathcal L_{\rm ent}\). Thus the same
object can act on a temperature or multiply dimensionless
information.

\begin{definition}[Thermal normalization of the tritic decision]
A thermal pair compatible with the oriented section \(a\) consists of
a positive transducer \(k_B\) and a positive section
\(\Theta_{{\rm clk},a}\in\mathcal L_\Theta\) such that
\begin{equation}
 k_B(\Theta_{{\rm clk},a})=\varepsilon_a,
 \qquad
 k_B\Theta_{{\rm clk},a}\log3=\epsilon_a.
 \label{x_reciprocidad:v:ant:eq:ii-kb-aut-par-termico}
\end{equation}
The Boltzmann constant serves as the energy--temperature
transducer of tritic information in this realization.
\end{definition}

\begin{theorem}[Uniqueness relative to the thermal section and dimensional covariance]
Given a positive section \(\Theta_{{\rm clk},a}\), there exists a
unique positive linear map satisfying
\eqref{x_reciprocidad:v:ant:eq:ii-kb-aut-par-termico}. For positive bases \(u_E,u_\Theta\),
if
\[
 \epsilon_a=\epsilon_a^{\rm num}u_E,\quad
 \Theta_{{\rm clk},a}=\theta_a u_\Theta,\quad
 k_B(u_\Theta)=b\,u_E,
\]
its coordinate is
\begin{equation}
 b=\frac{\epsilon_a^{\rm num}}{\theta_a\log3}.
 \label{x_reciprocidad:v:ant:eq:ii-kb-aut-coordenada}
\end{equation}
Under \(u'_E=s_Eu_E\), \(u'_\Theta=s_\Theta u_\Theta\), with
\(s_E,s_\Theta>0\), one has
\[
 b'=\frac{s_\Theta}{s_E}b,\qquad
 \theta'_a=\frac{\theta_a}{s_\Theta},\qquad
 (\epsilon_a^{\rm num})'=\frac{\epsilon_a^{\rm num}}{s_E}.
\]
\end{theorem}
\begin{proof}
A nonzero section generates a line. Prescribing its image determines
the unique linear map; the positivity of both sections proves
its positivity. Substituting the coordinates into
\eqref{x_reciprocidad:v:ant:eq:ii-kb-aut-par-termico} gives
\(b\theta_a\log3=\epsilon_a^{\rm num}\). The three transformation
laws are obtained by expressing the same sections and the same
map in the new bases.
\end{proof}

The same Boltzmann constant must act in both orientations.
Since \(t_P\) is common, simultaneous compatibility is equivalent to
\begin{equation}
 \frac{\Theta_{{\rm clk},{\rm pre}}}{\Theta_{{\rm clk},{\rm ret}}}
 =\frac{\epsilon_{\rm pre}}{\epsilon_{\rm ret}}
 =\frac{\hbar_{\rm pre}}{\hbar_{\rm ret}}.
 \label{x_reciprocidad:v:ant:eq:ii-kb-aut-orientaciones}
\end{equation}
Indeed, an isomorphism of lines preserves quotients of sections.
Conversely, if the preceding equality holds, the isomorphism
determined in one orientation also satisfies the normalization of
the other. Thus, return changes the thermal section of the cycle in the
same proportion as the energy, while the transducer remains
common to both readings.

The preceding equation makes explicit the choice required to publish an
individual value. If the thermal section has not yet been selected
by an independent reader, for each \(\lambda>0\) the pair
\[
 (k_B,\Theta_{{\rm clk},a})
 \longmapsto(\lambda k_B,\lambda^{-1}\Theta_{{\rm clk},a})
\]
preserves \(\epsilon_a\). Consequently, presenting an
internally selected individual coordinate requires specifying the reader of
\(\Theta_{{\rm clk},a}\), the bases and their dimensional transport.
The action and information equations remain determined before
that specification. The subsequent identification of the bases with
joule and kelvin belongs to the metrological publication of the result.

\section{Dimensional entropy and information of the reduced state}

For a distribution of admitted routes \(p\), or for a reduced
state with density matrix \(\rho\), let
\[
 s(p)=-\sum_xp_x\log p_x,\qquad
 s(\rho)=-\operatorname{Tr}(\rho\log\rho).
\]
Their dimensional publications are
\begin{equation}
 S(p)=k_Bs(p),\qquad S(\rho)=k_Bs(\rho)
 \quad\text{in }\mathcal L_{\rm ent}.
 \label{x_reciprocidad:v:ant:eq:ii-kb-aut-entropia}
\end{equation}
Evaluation at \(\Theta\) produces an energy
\(\Theta S=k_B\Theta\,s\). For the uniform triple at the cycle
temperature, \eqref{x_reciprocidad:v:ant:eq:ii-kb-aut-par-termico} gives exactly
\[
 S_{\rm TRIT}=k_B\log3,\qquad
 \Theta_{{\rm clk},a}S_{\rm TRIT}=\epsilon_a.
\]
If \(\rho_A\) is the reduction of the bipartite purification
constructed in the areal section, \(s(\rho_A)\) measures the entanglement
of that bipartition. Unit transport preserves its
eigenvalues, its orientation and the sectoral occupations. The areal
scale \(\gamma a_0\ell_P^2\) converts occupations into area; the
transducer \(k_B\) converts information into dimensional entropy.
Both readings use the same register through distinct operators.

\section{Spectral normalization of thermal histories}

Let \(\mathcal G\) be an already constructed finite graph of admissible
extensions, with adjacency matrix \(A\). Each edge \(e\) preserves
its route coordinates and receives a dimensionless action increment
\(a_*(e)>0\) from the corresponding reader. Having chosen the energy section
\(\epsilon_*\), set \(t=\beta\epsilon_*\), where
\(\beta=(k_B\Theta)^{-1}\). The thermal operator is
\begin{equation}
 (L_tf)(y)=\sum_{e:x\to y}\exp[-t a_*(e)]f(x).
 \label{x_reciprocidad:v:ant:eq:ii-kb-aut-operador}
\end{equation}
All its exponents are dimensionless. The condition
\(\rho(L_t)=1\) normalizes the weighted growth of histories.

% REMATE_LOCAL_X_TERMICA
\Needspace{35\baselineskip}
\begin{proposition}[Existence and uniqueness of normalized pressure]
Suppose that \(A\) is irreducible, \(\rho(A)>1\), and that
\(0<m\leq a_*(e)\leq M\) for all edges. There exists a unique
\(t_*>0\) such that \(\rho(L_{t_*})=1\), and
\begin{equation}
 \frac{\log\rho(A)}{M}\leq t_*\leq
 \frac{\log\rho(A)}{m}.
 \label{x_reciprocidad:v:ant:eq:ii-kb-aut-presion}
\end{equation}
\end{proposition}
\begin{proof}
Entrywise comparison gives
\(e^{-tM}A\leq L_t\leq e^{-tm}A\) for \(t\geq0\).
Monotonicity of the spectral radius of nonnegative matrices implies
the corresponding bounds for \(\rho(L_t)\). The spectral radius
is continuous, equals \(\rho(A)>1\) at zero and tends to zero at infinity.
For \(s>0\), moreover,
\(L_{t+s}\leq e^{-sm}L_t\), so that
\(\rho(L_{t+s})\leq e^{-sm}\rho(L_t)<\rho(L_t)\).
Continuity and strict decrease prove existence and
uniqueness. Substituting the value one into the initial bounds proves
\eqref{x_reciprocidad:v:ant:eq:ii-kb-aut-presion}.
\end{proof}

The local realization of three equally weighted alternatives, represented
by a vertex with three loops and \(a_*(e)=1\), satisfies
\(\rho(L_t)=3e^{-t}\). Its normalization gives \(t_*=\log3\).
With \(\epsilon_*=\epsilon_a\), we obtain
\(\beta_{\rm clk}=\log3/\epsilon_a\), precisely the relation
\eqref{x_reciprocidad:v:ant:eq:ii-kb-aut-par-termico}. This calculation identifies the
local tritic confluence. For a graph of histories with additional
restrictions, confluence with the same clock requires checking
\[
 \rho\!\left(L_{(\log3/\epsilon_a)\epsilon_*}\right)=1
\]
with its effective increments. The local normalization and the pressure
of the complete system thus retain their respective domains.

The resulting relation brings together structure, value and meaning. The register
determines the recoverable information; probabilities over
continuations determine its entropic reading; action and period
determine \(\epsilon_a\); the energy--temperature transducer
publishes dimensional entropy. Area and entanglement are
related through the reduced state and its boundary operators,
keeping visible which datum produces each quantity.

% END OWNER


% BEGIN OWNER /Users/ruben/Documents/New project/output/EXTREMA_MEDIA_RAZON_RECIPROCIDAD_ES_EN_20260918/ARTICULO/ES/source/antecedentes_vii/28_portador_tensorial.tex
% FORMALIZACION_REUNIDA: representación pentacomponente de II,
% con antecedente Paley y demostración directa de los momentos del marco.
% El retorno transporta la carta; no se identifica un ciclo de orden doce con A5.
\chapter{Oriented incidence and gravitational tensor representation}
\label{x_reciprocidad:vii:sec:pentacomponente}
The exceptional publication of the state preserves the incidence accompanying
its arithmetic reading. The ternary matrix \(A_W\) of
\eqref{x_reciprocidad:exc:aw} provides a precise material antecedent here. Its
balanced lift \(C=\operatorname{bal}(A_W)\), with representatives
\(0,1,-1\), satisfies
\[
 C=C^{\mathsf T},\qquad C^2=5I_6,\qquad \operatorname{Tr}C=0.
\]
The six positions have the order \((\infty,0,1,2,3,4)\).
This incidence first produces a three-dimensional space; its quadratic
reading then produces the five-component space on which
the gravitational coupling acts.

\section{Signed covering and calendar transport}
We define
\[
 E_C=\frac12(I_6+C/\sqrt5),\quad
 \mathcal H_C=\operatorname{im}E_C,\quad
 v_{u,\sigma}=\sigma\sqrt2E_Ce_u,\quad \sigma\in\{1,-1\}.
\]
\begin{proposition}[Twelve oriented positions]
\label{x_reciprocidad:vii:prop:doce-orientadas}
The vectors \(v_{u,\sigma}\) are twelve distinct unit vectors
in a three-dimensional space. Their incidence is determined by
\(\langle v_{u,\sigma},v_{v,\tau}\rangle=1/\sqrt5\).
Orientation reversal acts by antipodality.
\end{proposition}
\begin{proof}
The quadratic identity for \(C\) implies \(E_C^2=E_C=E_C^*\);
its trace is three. Since the diagonal of \(C\) is zero,
\[
 \langle v_{u,\sigma},v_{v,\tau}\rangle
 =\begin{cases}\sigma\tau,&u=v,\\
 \sigma\tau C_{uv}/\sqrt5,&u\ne v.
 \end{cases}
\]
This proves normalization, distinction of the twelve vectors and incidence.
Changing \(\sigma\) to \(-\sigma\) changes the sign of the vector.
\end{proof}

An explicit Euclidean expression is obtained by setting
\(\nu=\sqrt{1+\varphi^2}\) and
\[
 B=\frac1\nu
 \begin{pmatrix}
 0&-1&-\varphi&0&\varphi&1\\
 -1&-\varphi&0&1&0&-\varphi\\
 -\varphi&0&-1&-\varphi&-1&0
 \end{pmatrix}.
\]
The relation \(\varphi^2=\varphi+1\), for the coordinate already generated,
gives \(B^*B=2E_C\).
Consequently \(B/\sqrt2|_{\mathcal H_C}\) is an isometry and takes
\(v_{u,\sigma}\) to \(\sigma Be_u\). The image is
\[
 \mathcal V=\frac1\nu\bigl(
 \{0\}\times\{\pm1\}\times\{\pm\varphi\}\ \cup\
 \{\pm1\}\times\{\pm\varphi\}\times\{0\}\ \cup\
 \{\pm\varphi\}\times\{0\}\times\{\pm1\}\bigr).
\]
The internal coordinate, the incidence and its
twelve geometric representatives are thus related.

The calendar coordinate is written \((b,r)\in
\mathbb Z/12\mathbb Z\times\{0,\ldots,8\}\). A route of length \(n\)
induces
\[
 k(r,n)=\left\lfloor\frac{r+n}{9}\right\rfloor,\qquad
 (b,r)\longmapsto
 \bigl(b+k(r,n)\bmod12,\ (r+n)\bmod9\bigr).
\]
The explicit alignment of its positions \(p=b+1\) with the signed chart is
\[
\begin{array}{c|rrrrrrrrrrrr}
p&1&2&3&4&5&6&7&8&9&10&11&12\\ \hline
u&\infty&0&0&\infty&1&3&3&1&2&4&4&2\\
\sigma&+&-&+&-&-&+&-&+&-&+&-&+
\end{array}
\]
This table fixes a trivialization of the carrier. Each change of position
jointly transports the vectors and their incidence matrix.

\begin{proposition}[Covariance of the frame under return]
If \(S e_p=e_{p+1\bmod12}\), the matrix representation of an
operator \(P\) of the carrier is transported as \(P_k=S^kPS^{-k}\).
This transport preserves inner products, incidences and route
composition. After twelve returns the chart returns, preserving the memory
accumulated in the source state.
\end{proposition}
\begin{proof}
\(S\) is orthogonal and \(S^{12}=I\); the metric properties
and periodicity of the chart follow. For composition, with
\(r'=(r+n)\bmod9\), Euclidean division gives
\(k(r,n+m)=k(r,n)+k(r',m)\). The exponents of \(S\) are added
exactly with the carry. The operation acts on the calendar
coordinate; the route's memory updates remain
present in its enriched domain.
\end{proof}

\section{Quadratic reading and five-component space}
Let \(\mathbb R^{\mathcal V}\) be the space of the twelve coordinates,
\(\mathsf A e_v=e_{-v}\) its antipodal involution and
\(\mathsf J=\mathbf1\mathbf1^{\mathsf T}\). We define
\[
 P_5=\frac12(I+\mathsf A)-\frac1{12}\mathsf J,\qquad
 Q_v=vv^{\mathsf T}-\frac13I_3,\qquad
 \Gamma_0x=\sum_{v\in\mathcal V}x_vQ_v.
\]
\begin{theorem}[Five-component projector and tensor isometry]
\label{x_reciprocidad:vii:thm:portador-pentacomponente}
\(P_5\) is the rank-five orthogonal projector onto the space of
coordinates even under antipodality and of zero sum. Moreover,
\[
 \Gamma_0^*\Gamma_0=\frac85P_5,\qquad
 \Gamma_0\Gamma_0^*=\frac85I_{\operatorname{Sym}^2_0(\mathbb R^3)}.
\]
Therefore
\[
 \Gamma_5=\sqrt{\frac58}\,\Gamma_0|_{\operatorname{im}P_5}
\]
is an isometry onto the space of traceless symmetric tensors.
\end{theorem}
\begin{proof}
The even space has one coordinate for each antipodal pair, so
its dimension is six. The constant vector belongs to it; subtracting its
rank-one projector gives \(P_5\), of rank five.

Since
\(\langle Q_v,Q_w\rangle_F=(v^{\mathsf T}w)^2-1/3\), the entries
of \(\Gamma_0^*\Gamma_0\) equal \(2/3\) for \(w=\pm v\) and
\(-2/15\) in the other ten positions of each row. These are
exactly the entries of \((8/5)P_5\).

The identity can also be checked in tensor space without
resorting to a classification of representations. The list of vertices gives
\[
 \sum_vv_i^4=\frac{12}{5},\qquad
 \sum_vv_i^2v_j^2=\frac45\quad(i\ne j),
\]
and moments with any odd power vanish by sign
changes. Indeed,
\((1+\varphi^2)^2=5\varphi^2\) and \(1+\varphi^4=3\varphi^2\).
For a symmetric tensor \(Q\), these sums give
\[
 \sum_v(v^{\mathsf T}Qv)\,vv^{\mathsf T}
 =\frac85Q+\frac45(\operatorname{tr}Q)I.
\]
With \(\operatorname{tr}Q=0\), and using
\(\sum_vv^{\mathsf T}Qv=4\operatorname{tr}Q=0\), we obtain
\(\Gamma_0\Gamma_0^*Q=(8/5)Q\). The two identities prove
the normalized isometry and its surjectivity.
\end{proof}

The gravitational section in \(\ref{x_reciprocidad:v:ant:sec:gravity}\) determines the
operator \(G_aP_5\). Its spectrum has \(G_a\) with multiplicity
five and zero with multiplicity seven. Orientation is preserved in
the carrier and its transport; \(G_a\) fixes the common coupling.

\section{Transverse projection and helicities}
For \(n\in S^2\), set \(\Pi=I-nn^{\mathsf T}\) and
\[
 \Lambda_nQ=\Pi Q\Pi-\frac12\operatorname{tr}(\Pi Q\Pi)\Pi.
\]
\begin{proposition}[Rank-two transverse reduction]
\(\Lambda_n\) is the orthogonal projector onto
\(\{Q\in\operatorname{Sym}^2_0(\mathbb R^3):Qn=0\}\), of dimension two.
\end{proposition}
\begin{proof}
\(\Pi n=0\) gives transversality and \(\operatorname{tr}\Pi=2\) gives
zero trace. A transverse traceless tensor remains fixed.
The formula is self-adjoint for the Frobenius inner product.
In a frame \((e_1,e_2,n)\), the image has orthonormal basis
\[
 E_+=\frac{e_1e_1^{\mathsf T}-e_2e_2^{\mathsf T}}{\sqrt2},\qquad
 E_\times=\frac{e_1e_2^{\mathsf T}+e_2e_1^{\mathsf T}}{\sqrt2}.
\]
\end{proof}
The other three vectors of an orthonormal basis of the carrier are
\[
 E_1=\frac{e_1n^{\mathsf T}+ne_1^{\mathsf T}}{\sqrt2},\quad
 E_2=\frac{e_2n^{\mathsf T}+ne_2^{\mathsf T}}{\sqrt2},\quad
 E_0=\frac{2nn^{\mathsf T}-e_1e_1^{\mathsf T}-e_2e_2^{\mathsf T}}{\sqrt6}.
\]
When the transverse frame is rotated by an angle \(\theta\), the pairs
\((E_+,E_\times)\) and \((E_1,E_2)\) rotate by \(2\theta\) and
\(\theta\), respectively, while \(E_0\) remains fixed. The
complexification thus has weights \(2,-2,1,-1,0\).
\(\Lambda_n\) preserves the first two and annihilates the rest.
The chain is established
\[
 \mathbb R^{12}\xrightarrow{P_5}\operatorname{im}P_5
 \xrightarrow{\Gamma_5}\operatorname{Sym}^2_0(\mathbb R^3)
 \xrightarrow{\Lambda_n}\mathcal H_{\rm TT}(n),
 \qquad 12\longrightarrow5\longrightarrow5\longrightarrow2.
\]
The five-component representation and its reduction are algebraic
results preceding the choice of field equations.
In the massless gravitational realization, the gauge constraints
select the transverse image. This distinction preserves both the
five structural coordinates and the two radiative helicities.


% END OWNER


% BEGIN OWNER /Users/ruben/Documents/New project/output/EXTREMA_MEDIA_RAZON_RECIPROCIDAD_ES_EN_20260918/ARTICULO/ES/source/antecedentes_vii/29_retorno_y_memoria_traslacional.tex
% Artículo VII, recuperación aditiva, revisión 04, 10-09-2026.
% RESULTADO_RECUPERADO: S15, 11c_gravedad_holonomia_rev6.tex:16--126.
% El cociclo canónico c27=(1,18) se conserva como antecedente del cómputo
% de rutas; este fragmento prueba su composición y representación.
% La memoria traslacional finita y la amplitud cosmológica s0 tienen
% tipos distintos. El cociclo no evalúa por sí solo una densidad local.

\chapter{Phase return and representation of accumulated memory}
\label{x_reciprocidad:vii:sec:retorno-memoria-gravedad}

The chronology of enriched routes distinguishes position return,
orientation return and phase return. The observable publication
identifies certain arrival states, while the cumulative register
preserves the sequence traversed. The geometric representation must
transport both data jointly.

\section{Inversion cocycle and change of sheet}
\label{x_reciprocidad:vii:subsec:cociclo-memoria}

Let \(F_{\rm obs}\) be the chronological publication of TPK transport.
Its canonical blocks have the hierarchy
\[
\begin{array}{c|l}
27&\text{positional return with orientation reversal},\\
54&\text{positional and orientational return},\\
108&\text{complete return of the observable phase}.
\end{array}
\]
The register counts orientation inversions \(g(a)\) and
effective changes of sheet \(s(a)\). For a route,
\begin{equation}
 c(\gamma)=\sum_{a\subset\gamma}(g(a),s(a)),\qquad
 c(\gamma_2\gamma_1)=c(\gamma_2)+c(\gamma_1).
 \label{x_reciprocidad:vii:eq:cociclo-ruta}
\end{equation}
The canonical computation of the block of length \(27\) publishes
\(c_{27}=(1,18)\). Composition preserves those increments in
the four blocks that complete the observable period.
On forward-oriented routes the register is cumulative;
its extension to the route groupoid uses \(\mathbb Z^2\) and
\(c(\bar\gamma)=-c(\gamma)\).

\begin{proposition}[Memory of the complete observable turn]
\label{x_reciprocidad:vii:prop:memoria-vuelta-completa}
Writing \(\mathcal K_{27}=F_{\rm obs}^{27}\), the lift
to the register, for \(x\in\mathcal O_{108}\) and \(z\in\mathbb Z^2\), satisfies
\[
 \widetilde{\mathcal K}_{27}(x,z)
   =(\mathcal K_{27}x,z+(1,18)),\qquad
 \widetilde{\mathcal K}_{27}^{\,4}(x,z)
   =(x,z+(4,72)).
\]
After \(N\) complete turns, the increment is \(N(4,72)\).
\end{proposition}

\begin{proof}
The observable projection of \(\mathcal K_{27}\) has order four.
The preceding canonical computation fixes \(c_{27}=(1,18)\) as a constant
over that observable orbit. The additivity of
\eqref{x_reciprocidad:vii:eq:cociclo-ruta} sums the four increments and gives
\(c_{108}=4(1,18)=(4,72)\).
Repetition applies the same concatenation law and produces
\(Nc_{108}\). The return of the first coordinate preserves the
increment of the second.
\end{proof}

\section{Affine representation of the cocycle}
\label{x_reciprocidad:vii:subsec:representacion-afin-cociclo}

The discrete realization uses an element \(R_4\) of order four in
\(\operatorname{Spin}^+(1,3)\), a unit temporal vector \(u\)
fixed by its vector action and a unit of length \(\ell_0>0\).
The scale \(\ell_0\) retains its provenance in the dimensional section.
We define
\begin{equation}
 \rho_{\rm C}^{\rm disc}(\gamma)=
 \bigl(R_4^{c_g(\gamma)},\,\ell_0c_s(\gamma)u\bigr)
 \in\operatorname{Spin}^+(1,3)\ltimes\mathbb R^{1,3}.
 \label{x_reciprocidad:vii:eq:representacion-discreta-memoria}
\end{equation}
The power acts on the vector through the associated vector
representation of the spin group.

\begin{theorem}[Composition and translational amplitude of return]
\label{x_reciprocidad:vii:thm:traslacion-memoria}
The map \(\rho_{\rm C}^{\rm disc}\) preserves the identity,
concatenation and inversion. At complete returns,
\begin{equation}
 \rho_{\rm C}^{\rm disc}(\gamma_{108})=(I,72\ell_0u),\qquad
 \rho_{\rm C}^{\rm disc}(\gamma_{108}^{\,N})=(I,72N\ell_0u).
 \label{x_reciprocidad:vii:eq:traslacion-retorno-108}
\end{equation}
Inversion changes the sign of the translation.
\end{theorem}

\begin{proof}
The semidirect product is
\((R,a)(S,b)=(RS,a+Rb)\). Since \(R_4u=u\),
\[
\begin{split}
 \rho_{\rm C}^{\rm disc}(\gamma_2)
 \rho_{\rm C}^{\rm disc}(\gamma_1)
 &=
 \bigl(R_4^{c_g(\gamma_2)+c_g(\gamma_1)},
 \ell_0(c_s(\gamma_2)+c_s(\gamma_1))u\bigr)\\
 &=\rho_{\rm C}^{\rm disc}(\gamma_2\gamma_1).
\end{split}
\]
The identity cocycle is zero and that of the inverse is its opposite.
For the complete turn, substitute \(c_{108}=(4,72)\) and
\(R_4^4=I\). Concatenating \(N\) turns keeps the rotational part
equal to the identity and sums their translations.
\end{proof}

The amplitude \(72\ell_0\) is a length associated with a finite route.
A realization through a co-tetrad and a connection transports it to its
affine holonomy. The local contribution to the gravitational equations
is subsequently obtained through the curvature of that connection, the
variation of the matter action and the normalization of its density.
The representation thus preserves the exact cumulative content that
must enter that realization.

% END OWNER


% BEGIN OWNER /Users/ruben/Documents/New project/output/EXTREMA_MEDIA_RAZON_RECIPROCIDAD_ES_EN_20260918/ARTICULO/ES/source/antecedentes_vii/30_pantalla_y_respuesta.tex
% Artículo VII. Incorporación de resultados preexistentes, 10-09-2026.
% Propietarios íntegros preservados en fuentes_conservadas:
% 17_pantalla_cubica_soldadura.tex (P), 18_clausura_energetica.tex (E).
% P: incidencia, cociente, soldadura celular y refinamiento.
% E: respuesta, unicidad, extensión mínima y expectativa tracial.
% Dependencias anteriores de VII: núcleo APP--TRIT--TPK; realización
% cúbica de la ruta, cociclo de transporte y compresión q_vis=5/9.
% Este fragmento no constituye por sí solo el artículo autónomo.
% La identidad de pantalla se distingue de la contracción material de espín.

\chapter{Boundary soldering and energy response}
\label{x_reciprocidad:vii:sec:pantalla-respuesta}

The cubic realization of transport preserves six oriented incidences.
Its geometric quotient, the Lorentzian form and the energy response
are obtained from the same incidence matrix. Transport proceeds from
the enriched APP--TRIT--TPK route: the face, sign, carry and
memory accompany each edge. This section develops the
linear operation acting on that boundary publication.

\section{Cubic incidence and the rank-four quotient}
\label{x_reciprocidad:vii:sec:cociente-pantalla}

Let \(H_F=\mathbb R^F\), with
\(F=\{(i,\epsilon):1\leq i\leq3,\ \epsilon\in\{-1,1\}\}\),
its Euclidean inner product, and its basis \(e_{i,\epsilon}\).
Face opposition is the orthogonal involution
\(Re_{i,\epsilon}=e_{i,-\epsilon}\). Define
\[
 u=\sum_{i,\epsilon}e_{i,\epsilon},\qquad
 P_u=\frac{uu^{\mathsf T}}6,\qquad P_-=\frac{I-R}{2},\qquad
 P_0=\frac{I+R}{2}-P_u,\qquad P_\square=P_u+P_-.
\]
The vertices are \(V=\{-1,1\}^3\). The incidence matrix
\(M:\mathbb R^V\longrightarrow H_F\) is determined by
\[
 M_{(i,\epsilon),\nu}=\mathbf1_{\{\nu_i=\epsilon\}}.
\]

\begin{proposition}[Incidence decomposition]
\label{x_reciprocidad:vii:prop:incidencia}
We have
\[
 MM^{\mathsf T}=12P_u+4P_-,
 \qquad
 \ker M^{\mathsf T}=P_0H_F.
\]
Consequently, the quotient
\(Q=H_F/P_0H_F\), orthogonally identified with
\(P_\square H_F\), has dimension four.
\end{proposition}

\begin{proof}
A face contains four vertices. Two opposite faces have no common
vertices and two faces on different axes have two common vertices.
Writing \(\mathbf J=uu^{\mathsf T}\), we obtain
\[
 MM^{\mathsf T}=4I+2(\mathbf J-I-R)
 =2(I-R)+2\mathbf J=4P_-+12P_u.
\]
The three projectors \(P_u,P_-,P_0\) are orthogonal and sum to \(I\);
their ranks are one, three and two, respectively. The identity
\(\|M^{\mathsf T}x\|^2=\langle x,MM^{\mathsf T}x\rangle\)
identifies the stated kernel. The complementary rank is four.
\end{proof}

On \(Q\), the form
\[
 B=P_--P_u
\]
has signature \((3,1)\). The time orientation is fixed by \(u\).
With
\[
 t=\frac{u}{\sqrt6},\qquad
 d_i=\frac{e_{i,+}-e_{i,-}}{\sqrt2},\qquad
 W=\sqrt6P_u+\sqrt2P_-,
\]
we have \(We_{i,\epsilon}=t+\epsilon d_i\). These six vectors
are null for \(B\), whereas
\[
 3t+\nu_1d_1+\nu_2d_2+\nu_3d_3
\]
has squared norm \(-6\). The same incidence satisfies
\(MM^{\mathsf T}=2W^*W\) on \(Q\). The proper rotations of the cube
commute with the projectors, preserve both forms and fix \(t\).
These identities determine the metric of the cubic publication.

\section{Cellular soldering and memory transport}
\label{x_reciprocidad:vii:sec:soldadura-celular}

For each oriented edge \(a\), the route publication provides
the face \(\lambda(\underline a)\), the sign
\(\sigma(a)\), the source and target fibres and the transport
\(U(a):Q_{s(a)}\longrightarrow Q_{t(a)}\).
Inversion satisfies \(U(\bar a)=U(a)^{-1}\).
The boundary half-edges retain their formal inversion and their incidence
at the boundary. With \(\ell_m=\ell_0\,9^{-m}>0\), the soldering is
\begin{equation}
 \beta_m(a)=\ell_m\sigma_m(a)
              n_{\lambda(\underline a)}^{s(a)},\qquad
 n_{i,\epsilon}=t+\epsilon d_i.
 \label{x_reciprocidad:vii:eq:soldadura-arista}
\end{equation}
The face and sign are transported together with the route memory.
The inversion convention requires
\begin{equation}
 \beta_m(\bar a)=-U_m(a)\beta_m(a).
 \label{x_reciprocidad:vii:eq:inversion-soldadura}
\end{equation}
Applying this relation twice recovers \(\beta_m(a)\);
the carry register separately preserves the ordered composition.

\begin{proposition}[Naturality of soldering under refinement]
\label{x_reciprocidad:vii:prop:refinamiento-soldadura}
Let \(a=a_1\cdots a_9\) be a compatible refinement of an edge.
Suppose that the nine child incidences, transported to the parent
fibre, preserve the face and sign, and that
\(\ell_{m+1}=\ell_m/9\). Let
\(\mathcal U_j:Q_{m,s(a)}\longrightarrow Q_{m+1,s(a_j)}\)
be the accumulated transport that includes the refinement identification
and the path to the origin of child \(j\). Then
\[
 \beta_m(a)=
 \sum_{j=1}^{9}\mathcal U_j^{-1}\beta_{m+1}(a_j).
\]
The equality is compatible with inversion and with concatenation
of the memory register.
\end{proposition}

\begin{proof}
Each summand, expressed in the fibre at the origin of the parent, is
\((\ell_m/9)\sigma_m(a)n_{\lambda(\underline a)}\).
The sum of nine terms reproduces the parent soldering.
Inversion reverses the order of the transports and replaces each one
with its inverse; the relation \eqref{x_reciprocidad:vii:eq:inversion-soldadura} provides
the overall sign. Memory is concatenated in that same order, so
refinement preserves both the vector publication and its
enriched register.
\end{proof}

\section{Positive response, uniqueness, and minimal extension}
\label{x_reciprocidad:vii:sec:respuesta-positiva}

The linear screen soldering is
\(\bar\beta_m=\ell_mW|_Q\). It is invertible, with
\[
 \bar\beta_m^{-1}
 =\ell_m^{-1}\left(\frac1{\sqrt6}P_u+\frac1{\sqrt2}P_-\right).
\]
The eigenvalue \(5/9\) of the antisymmetric sector of the normalized block
\(B_c/9\), constructed in Section~\ref{x_reciprocidad:vii:sec:bloque-compresion},
determines the compression
coefficient \(q_{\mathrm{vis}}=5/9\). We reserve this notation to
distinguish it from the Euclidean quotient \(q_9(n)\) of the arithmetic kernel.
Its realization on the screen uses an isometry
\(V_9:Q\longrightarrow\mathcal K\). The visible contribution
\(C_9=q_{\mathrm{vis}}V_9\) and its complement
\(D_9=\sqrt{1-q_{\mathrm{vis}}^2}\,V_9\) satisfy
\[
 C_9^*C_9+D_9^*D_9=I_Q,\qquad D_9^*D_9=\frac{56}{81}I_Q.
\]
The coefficient \(56\) is \(9^2-5^2\); it belongs to the quadratic defect
of that compression reader.

\begin{theorem}[Energy response determined by soldering]
\label{x_reciprocidad:vii:thm:respuesta-unica}
There is a unique operator \(\Lambda_m:Q\longrightarrow Q\) such that
\(\bar\beta_m^*\Lambda_m\bar\beta_m=D_9^*D_9\). It is positive definite and
\begin{equation}
 \Lambda_m=\frac{28}{243\ell_m^2}(P_u+3P_-)
 =\frac{112}{81\ell_m^2}
       (MM^{\mathsf T}|_Q)^{-1}.
 \label{x_reciprocidad:vii:eq:lambda-pantalla}
\end{equation}
On \(H_F\), the corresponding identity has right-hand side
\((56/81)P_\square\).
\end{theorem}

\begin{proof}
Congruence by the invertible soldering necessarily determines
\[
 \Lambda_m=\bar\beta_m^{-*}\frac{56}{81}I_Q\bar\beta_m^{-1}.
\]
Substituting the inverse, the coefficient of \(P_u\) is
\(56/(486\ell_m^2)=28/(243\ell_m^2)\); that of \(P_-\) is
\(28/(81\ell_m^2)\). Both are positive. By
Proposition \ref{x_reciprocidad:vii:prop:incidencia},
\((MM^{\mathsf T}|_Q)^{-1}=P_u/12+P_-/4\),
which gives the second expression. Extension to six faces annihilates
\(P_0H_F\); the identity therefore extends through \(P_\square\).
\end{proof}

\begin{proposition}[Characterization by minimal extension]
\label{x_reciprocidad:vii:prop:minima-extension}
For every \(b\in Q\),
\[
 \langle b,\Lambda_m b\rangle
 =\frac{112}{81\ell_m^2}
   \min_{Mx=b}\|x\|^2.
\]
The minimizing extension is unique:
\(x_{\min}=M^{\mathsf T}(MM^{\mathsf T}|_Q)^{-1}b\).
\end{proposition}

\begin{proof}
The stated vector satisfies \(Mx_{\min}=b\) and belongs to
\(\operatorname{im}M^{\mathsf T}=(\ker M)^\perp\).
Every other solution is \(x_{\min}+z\), with \(z\in\ker M\); by
orthogonality, its squared norm is \(\|x_{\min}\|^2+\|z\|^2\).
Moreover,
\[
 \|x_{\min}\|^2
 =\langle b,(MM^{\mathsf T}|_Q)^{-1}b\rangle.
\]
The formula for \(\Lambda_m\) completes the proof.
\end{proof}

The operator \(MM^{\mathsf T}|_Q\) realizes the
Dirichlet--to--Neumann response of this incidence formulation; its inverse
is the Neumann--to--Dirichlet response. The preceding energy corresponds
to the latter, with the normalization of the soldering and the compression
defect. The Schur complement of
\(\left(\begin{smallmatrix}I&M^{\mathsf T}\\M&0\end{smallmatrix}\right)\)
is \(-MM^{\mathsf T}\), which offers the same variational elimination.

\section{Preservation of density by refinement}
\label{x_reciprocidad:vii:sec:densidad-refinamiento}

Consider the sector of cubic transports, orthogonal for the energy
inner product and compatible with the form \(B\). We fix a
common label space \(Q^{\mathrm{lab}}\), identified
isometrically with the parent's screen. The solderings have
types
\(\bar\beta_m:Q^{\mathrm{lab}}\longrightarrow Q_m\) and
\(\bar\beta_{m+1,j}:Q^{\mathrm{lab}}\longrightarrow Q_{m+1,j}\);
the operators \(\Lambda\) act in their respective target fibers.
If
\(U_j:Q_j\longrightarrow Q\) transports the child fiber to the parent's,
then
\[
 \bar\beta_{m+1,j}=\frac19U_j^{-1}\bar\beta_m,\qquad
 \Lambda_{m+1,j}=81U_j^{-1}\Lambda_mU_j.
\]
The factor \(81\) comes from \(\ell_{m+1}^{-2}=81\ell_m^{-2}\).

\begin{theorem}[Invariance of the normalized density]
\label{x_reciprocidad:vii:thm:densidad}
Each child preserves the parent's quadratic response:
\[
 \bar\beta_{m+1,j}^*\Lambda_{m+1,j}\bar\beta_{m+1,j}
 =\frac{56}{81}I_{Q^{\mathrm{lab}}}.
\]
For \(r\) refinements, the joint response, expressed in the parent's
fibers, is
\[
 E_{m+r}=\frac{56}{81}I_{Q^{\mathrm{lab}}}\otimes I_{9^r}.
\]
The expectation with normalized trace
\(\tau_{9^r}(A)=9^{-r}\operatorname{Tr}(A)\) satisfies
\[
 (\operatorname{id}\otimes\tau_{9^r})(E_{m+r})
 =\frac{56}{81}I_{Q^{\mathrm{lab}}}
\]
at every finite depth.
\end{theorem}

\begin{proof}
The congruence of a child contains the factors \(1/9,81,1/9\);
their product is one. Orthogonality cancels the intermediate
transports. The direct sum of nine equal responses identifies
the first refinement with \((56/81)I_Q\otimes I_9\).
Iteration produces the tensor with \(I_{9^r}\).
Finally \(\tau_{9^r}(I_{9^r})=1\). The expectations compose
through the identity
\(\tau_{9^{r+s}}=\tau_{9^r}\otimes\tau_{9^s}\).
\end{proof}

The unital inclusions \(A\mapsto A\otimes I_9\) are compatible with
that family of traces. In the inductive limit of the matrix algebras,
the response is the central element
\((56/81)I_Q\otimes\mathbf1\). The unnormalized extensive sum produces,
in contrast, \(9^r(56/81)I_Q\). This distinction specifies which quantity
is preserved: the quadratic screen density.

\section{Variational realization of the response}
\label{x_reciprocidad:vii:sec:realizacion-respuesta}

The material realization uses an isometry
\(J_{\mathrm{scr}}:Q\longrightarrow\mathcal H_\partial\)
compatible with the spin action and with refinement.
The boundary co-tetrad is \(e=J_{\mathrm{scr}}\bar\beta_m\).
Let \(\Lambda_{\mathrm{CE}}\) be the boundary response obtained
by second variation of the material and geometric action, with its
boundary conditions and elimination of the interior variables.
The compatibility defect is
\begin{equation}
 \mathcal R_{\mathrm{grav}}
 =\bar\beta_m^*
   (J_{\mathrm{scr}}^*\Lambda_{\mathrm{CE}}J_{\mathrm{scr}}-\Lambda_m)
   \bar\beta_m
 =e^*\Lambda_{\mathrm{CE}}e-D_9^*D_9.
 \label{x_reciprocidad:vii:eq:defecto-respuesta}
\end{equation}
By the invertibility of \(\bar\beta_m\),
\[
 \mathcal R_{\mathrm{grav}}=0
 \quad\Longleftrightarrow\quad
 J_{\mathrm{scr}}^*\Lambda_{\mathrm{CE}}J_{\mathrm{scr}}=\Lambda_m.
\]
The isometry \(J_{\mathrm{scr}}\) transports the boundary states;
the variational spin current is a degree-three form with values
in bivectors. Each enters with its own domain in the
realization of the response.

The construction determines a positive energy form, its minimizing
interior extension and the density preserved by refinement.
The amplitude of a concrete material contribution is obtained by evaluating
the current of that realization in its variational quadratic form.
This differentiates the operator fixing the response, the current
on which it acts and the gravitational observable published by the evaluation.

% END OWNER


% BEGIN OWNER /Users/ruben/Documents/New project/output/EXTREMA_MEDIA_RAZON_RECIPROCIDAD_ES_EN_20260918/ARTICULO/ES/source/antecedentes_vii/30d_realizacion_geometrica.tex
% RESULTADO_RECUPERADO: S15, 11c_gravedad_holonomia_rev6.tex:128--212;
% S01, 02_dinamica_tres_hojas.tex:155--219; pantalla y respuesta de30.
% Pruebas de curvatura, covariancia y Bianchi reunidas para autonomía.
% La realización suave conserva sus condiciones de holonomía explícitas.

\chapter{Geometric realization of transport and soldering}
\label{x_reciprocidad:vii:sec:realizacion-cartan}

The discrete representation of return and screen soldering
determine two aspects of transport: the accumulation of memory along
routes and the comparison of their fibres. Their geometric realization
preserves both operations. This section fixes the target
space, the dimensional normalization of the connection and the identities
that its variation will use.

\section{Screen, co-tetrad and connection}

The screen quotient \(Q_\square\), its Lorentzian form, and the
soldering \(\beta_m\) come from
\ref{x_reciprocidad:vii:sec:pantalla-respuesta}. In each cell, a screen realization
\(J_{{\rm scr},m}\) transports the form and soldering to the
corresponding geometric fibre. Its cellular image is
\[
 e_m=J_{{\rm scr},m}\beta_m.
\]
The scale \(\ell_m\) remains within the soldering. Changes
of frame act jointly on \(J_{{\rm scr},m}\), the connection
and the transport of the cells.

A smooth realization of those data on an oriented manifold
\(\mathcal U\) of dimension four consists of a nondegenerate
co-tetrad \(e=(e^I)\), a Lorentz connection
\(\omega=(\omega^I{}_J)\) and a realization of routes
\(\mathfrak R_{\rm C}\) satisfying
\begin{equation}
 \operatorname{Hol}_{\mathcal A_{\ell_0}}
       (\mathfrak R_{\rm C}(\gamma))
 =\rho_{\rm C}(\operatorname{Hol}_{\rm TPK}(\gamma)).
 \label{x_reciprocidad:vii:eq:compatibilidad-holonomia-geometrica}
\end{equation}
In this matrix equality, the representation is normalized by
\[
 N_{\ell_0}(R,a)=(\Lambda(R),a/\ell_0),\qquad
 \rho_{\rm C}=N_{\ell_0}\circ\rho_{\rm C}^{\rm disc},
\]
where \(\Lambda:\operatorname{Spin}^+(3,1)\to SO^+(3,1)\)
is the vector representation. The map \(N_{\ell_0}\)
preserves the semidirect product: dividing the translation
\(a+\Lambda(R)b\) by \(\ell_0\) is equivalent to composing the
two normalized translations. The complete turn therefore publishes
\(72u\) in the normalized connection and \(72\ell_0u\) in the
dimensional chart. The spin lift separately preserves
\(R_4^{c_g(\gamma)}\); the vector representation has the
central kernel \(\{1,-1\}\), so that lift is retained
when using spinor fields.
Identity, concatenation, inversion and subdivision are maintained in the
corresponding realization. The screen pairing
and the holonomy condition specify the data of the
realization used in the field equations.

With \(\eta=\operatorname{diag}(-1,1,1,1)\), define
\(g=\eta_{IJ}e^I\otimes e^J\). For a positive elementary scale
\(\ell_0\), constant in the chart, the affine connection is
\begin{equation}
 \mathcal A_{\ell_0}
 =\begin{pmatrix}\omega&\ell_0^{-1}e\\0&0\end{pmatrix}.
 \label{x_reciprocidad:vii:eq:conexion-afin-realizada}
\end{equation}
The co-tetrad has the dimension of length; the factor
\(\ell_0^{-1}\) makes its translational component dimensionally homogeneous.

\begin{proposition}[Curvature and torsion of the realized connection]
\label{x_reciprocidad:vii:prop:curvatura-afin}
The curvature of \eqref{x_reciprocidad:vii:eq:conexion-afin-realizada} is
\begin{equation}
 \mathcal F_{\ell_0}
 =\mathrm d\mathcal A_{\ell_0}
   +\mathcal A_{\ell_0}\wedge\mathcal A_{\ell_0}
 =\begin{pmatrix}F_\omega&\ell_0^{-1}T\\0&0\end{pmatrix},
 \quad F_\omega=\mathrm d\omega+\omega\wedge\omega,
 \quad T=\mathrm de+\omega\wedge e.
 \label{x_reciprocidad:vii:eq:curvatura-afin-realizada}
\end{equation}
\end{proposition}
\begin{proof}
The product of matrices of forms produces
\(\omega\wedge\omega\) in the diagonal block and
\(\ell_0^{-1}\omega\wedge e\) in the translational block.
The exterior derivative contributes
\(\mathrm d\omega\) and \(\ell_0^{-1}\mathrm de\), since
\(\ell_0\) is constant. Their sum proves the equality.
\end{proof}

Curvature describes the frame-transport defect and torsion
that of the co-tetrad. Both belong to the same realized connection;
their components retain distinct tensorial domains.

\begin{proposition}[Covariance and Bianchi identities]
\label{x_reciprocidad:vii:prop:bianchi-cartan}
Under a change of frame \(g_L:\mathcal U\to SO^+(3,1)\),
\[
 e'=g_Le,\qquad
 \omega'=g_L\omega g_L^{-1}-\mathrm dg_L\,g_L^{-1},
\]
we have \(T'=g_LT\) and \(F_{\omega'}=g_LF_\omega g_L^{-1}\).
Furthermore,
\begin{equation}
 D_\omega T=F_\omega\wedge e,\qquad D_\omega F_\omega=0.
 \label{x_reciprocidad:vii:eq:bianchi-realizacion}
\end{equation}
\end{proposition}
\begin{proof}
In \(\mathrm d(g_Le)+\omega'\wedge g_Le\), the terms
containing \(\mathrm dg_L\wedge e\) cancel.
The same substitution into the curvature gives its transformation by
conjugation. For the first Bianchi identity, expand
\[
 D_\omega(\mathrm de+\omega\wedge e)
 =(\mathrm d\omega+\omega\wedge\omega)\wedge e.
\]
In \(\mathrm dF_\omega+\omega\wedge F_\omega-F_\omega\wedge\omega\),
the terms containing \(\mathrm d\omega\) and the cubic
products cancel in pairs. This proves the second identity.
\end{proof}

\section{The three quadratic realizations of curvature}

The TRIT sheet \(\tau\in\{+1,0,-1\}\) is represented in
the Cartan algebra by generators \(J_{ab}=-J_{ba}\)
and \(P_a^{(\tau)}\), with relations
\begin{align*}
 [J_{ab},J_{cd}]&=
 \eta_{bc}J_{ad}-\eta_{ac}J_{bd}
 -\eta_{bd}J_{ac}+\eta_{ad}J_{bc},\\
 [J_{ab},P_c^{(\tau)}]&=
 \eta_{bc}P_a^{(\tau)}-\eta_{ac}P_b^{(\tau)},\\
 [P_a^{(\tau)},P_b^{(\tau)}]&=-\tau J_{ab}.
\end{align*}
For constant \(\ell>0\), write
\[
 \mathcal A_\tau^{\rm Cart}
 =\tfrac12\omega^{ab}J_{ab}+\ell^{-1}e^aP_a^{(\tau)}.
\]

\begin{proposition}[Curvature of the tritified realization]
\label{x_reciprocidad:vii:prop:curvatura-tritificada}
With the preceding relations,
\begin{equation}
 \mathcal F_\tau^{\rm Cart}
 =\tfrac12(F_\omega^{ab}-\tau\ell^{-2}e^a\wedge e^b)J_{ab}
     +\ell^{-1}T^aP_a^{(\tau)}.
 \label{x_reciprocidad:vii:eq:curvatura-tritificada}
\end{equation}
\end{proposition}
\begin{proof}
Expand \(\mathrm d\mathcal A+\tfrac12[\mathcal A,\mathcal A]\).
The bracket of the Lorentz sector produces \(F_\omega\), the
mixed bracket produces \(D_\omega e\), and the translational
bracket contributes
\(-\tfrac12\tau\ell^{-2}e^a\wedge e^bJ_{ab}\).
\end{proof}

On the central sheet the affine connection is recovered; the extreme sheets
preserve the two signs of the constant-curvature term. The
representation thus retains the tritic regime of the route before
applying the variational equations. The holonomy equality
\eqref{x_reciprocidad:vii:eq:compatibilidad-holonomia-geometrica} corresponds to the
central sheet. On each extreme sheet, the representation
\(\rho_{{\rm C},\tau}\) in its Cartan group and its respective holonomy
condition are used. The current and the metric tensor use
the fields and representation of the same realization.

% END OWNER


% BEGIN OWNER /Users/ruben/Documents/New project/output/EXTREMA_MEDIA_RAZON_RECIPROCIDAD_ES_EN_20260918/ARTICULO/ES/source/antecedentes_vii/30b_variacion_energia_memoria.tex
% Artículo VII. Desarrollo aditivo, revisión 04, 10-09-2026.
% ARQUITECTURA_AUTORAL_PREEXISTENTE: transporte enriquecido, energía
% D_9^* W D_9 y compatibilidad por refinamiento, propietario S01,
% 02_dinamica_tres_hojas.tex, líneas 221--250.
% FORMALIZACION_NUEVA / CERTIFICADO_NUEVO: diferencial explícito de ese
% funcional, representante antihermítico, composición e inversión ponderada.
% Las variaciones pertenecen al espacio de realizaciones unitarias del
% transporte. No se presume que toda variación continua sea una ruta TPK.
% Su realización como corriente material de espín se trata por separado.
% \mathsf W_m es el peso energético; W de la sección anterior es soldadura.

\chapter{Variational response of memory energy}
\label{x_reciprocidad:vii:sec:variacion-memoria}

The transport of a route identifies the fibers at its endpoints and preserves
the orientation and memory entering that identification. The associated
energy compares the final field with the transported initial field. Its
differential quantifies how that comparison responds to a variation of
transport. This section obtains the response directly from the discrete
functional and develops its composition, covariance and inversion.

\section{Discrete functional and space of variations}
\label{x_reciprocidad:vii:subsec:funcional-memoria}

At a finite level \(m\), let \(\mathcal V_m\) be the vertices of the
route publication and \(\mathcal E_m^+\) an orientation of each edge.
Each vertex carries a finite-dimensional Hermitian fiber \(H_v\).
We define
\[
 \mathcal H_m=\bigoplus_{v\in\mathcal V_m}H_v,\qquad
 \mathcal K_m=\bigoplus_{a\in\mathcal E_m^+}H_{t(a)}.
\]
The Hermitian inner product is antilinear in the first variable.
For unitary \(U_a:H_{s(a)}\to H_{t(a)}\), the difference operator is
\begin{equation}
 (D_mf)_a=f_{t(a)}-U_af_{s(a)},\qquad
 E_m(f,U)=\langle D_mf,\mathsf W_mD_mf\rangle,\qquad
 \mathsf W_m=\mathsf W_m^*>0.
 \label{x_reciprocidad:vii:eq:energia-memoria-discreta}
\end{equation}
This is the energy of the term \(D_m^*\mathsf W_mD_m\) in the discrete generator.
The weight \(\mathsf W_m\) may couple different edges. The soldering operator \(W\)
of Section~\ref{x_reciprocidad:vii:sec:pantalla-respuesta} and this energy
weight have different domains and roles.

The transport \(U\) is a publication of the enriched state. The field
\(f\) and the weight \(\mathsf W_m\) remain explicit as arguments
of the functional; each material application preserves the rules that
determine them. Differentiable families of the unitary realizations
of transport are considered,
\[
 \dot U_a=A_aU_a,\qquad A_a^*=-A_a.
\]
The variations used physically form the subspace preserving
the conditions of the corresponding realization. The differential is
calculated first in the unitary tangent space and then restricted
to that subspace.

\begin{proposition}[Transport-response covector]
\label{x_reciprocidad:vii:prop:covector-memoria}
Write
\[
 v_a=U_af_{s(a)},\qquad r=D_mf,\qquad q=\mathsf W_mr.
\]
At fixed field and weight, the energy differential is
\begin{equation}
 \mathfrak j_m(A)=-2\operatorname{Re}
       \sum_{a\in\mathcal E_m^+}\langle q_a,A_av_a\rangle .
 \label{x_reciprocidad:vii:eq:covector-memoria}
\end{equation}
Its representative for the real Hilbert--Schmidt inner product is
\begin{equation}
 \mathcal J_a=v_aq_a^*-q_av_a^*,\qquad
 \mathcal J_a^*=-\mathcal J_a,\qquad
 \mathfrak j_m(A)=
 \sum_a\operatorname{Re}\operatorname{Tr}(\mathcal J_a^*A_a).
 \label{x_reciprocidad:vii:eq:representante-corriente-discreta}
\end{equation}
For simultaneous variations of field, transport and weight we have
\begin{equation}
 \dot E_m=
 2\operatorname{Re}\sum_a
 \left\langle q_a,\dot f_{t(a)}-U_a\dot f_{s(a)}-A_av_a\right\rangle
 +\langle r,\dot{\mathsf W}_mr\rangle .
 \label{x_reciprocidad:vii:eq:variacion-completa-memoria}
\end{equation}
\end{proposition}

\begin{proof}
The first-order update, written in the algebra of dual
numbers with \(\epsilon^2=0\), gives
\[
 r+\epsilon\dot r,\qquad
 \dot r_a=\dot f_{t(a)}-U_a\dot f_{s(a)}-A_av_a.
\]
The coefficient of \(\epsilon\) in the energy is
\[
 \langle\dot r,\mathsf W_mr\rangle+
 \langle r,\mathsf W_m\dot r\rangle+
 \langle r,\dot{\mathsf W}_mr\rangle.
\]
Self-adjointness of the weight yields
\eqref{x_reciprocidad:vii:eq:variacion-completa-memoria}. Fixing field and weight leaves
\eqref{x_reciprocidad:vii:eq:covector-memoria}. Moreover,
\[
 \operatorname{Tr}(\mathcal J_a^*A_a)
 =v_a^*A_aq_a-q_a^*A_av_a.
\]
Anti-Hermiticity of \(A_a\) implies
\(v_a^*A_aq_a=-\overline{q_a^*A_av_a}\), giving
the stated representative.
\end{proof}

The response is thus evaluated: on each edge the field is transported,
its difference from the final field is calculated and the energy weight is applied.
The antisymmetrized product of those two vectors determines
\(\mathcal J_a\in\mathfrak u(H_{t(a)})\). This operator represents,
through the real Hilbert--Schmidt inner product, the response covector
on the unitary variations of the target fiber.

\section{Composition and transport of the response}
\label{x_reciprocidad:vii:subsec:composicion-respuesta}

\begin{proposition}[Composition rule along a route]
\label{x_reciprocidad:vii:prop:composicion-respuesta}
For a route formed by \(N\) transports, let
\(U_\gamma=U_N\cdots U_1\) and
\(B_k=U_N\cdots U_{k+1}\), with \(B_N=I\).
If \(\dot U_k=A_kU_k\), then
\begin{equation}
 \dot U_\gamma U_\gamma^{-1}
 =\sum_{k=1}^N B_kA_kB_k^{-1}.
 \label{x_reciprocidad:vii:eq:variacion-transporte-compuesto}
\end{equation}
Thus a representative \(\mathcal J_\gamma\) in the final fibre
induces in factor \(k\) the contribution
\(B_k^*\mathcal J_\gamma B_k\). When a factor belongs to several
routes in the functional, its contributions are summed with the incidences
used by that functional.
\end{proposition}

\begin{proof}
The product rule gives
\[
 \dot U_\gamma=\sum_k
 U_N\cdots U_{k+1}A_kU_k\cdots U_1.
\]
Multiplication by \(U_\gamma^{-1}\) eliminates the factors to the right
of \(A_k\) and yields \eqref{x_reciprocidad:vii:eq:variacion-transporte-compuesto}.
By unitarity and cyclicity of the trace,
\[
 \operatorname{Re}\operatorname{Tr}
   (\mathcal J_\gamma^*B_kA_kB_k^{-1})
 =\operatorname{Re}\operatorname{Tr}
   ((B_k^*\mathcal J_\gamma B_k)^*A_k).
\]
Linearity of the differential proves the sum of contributions.
\end{proof}

This formula preserves the position of each operation within the route.
Transport following an incidence determines how
its contribution is represented in the final fiber; the local response is recovered
by transporting it to the fiber of that incidence.

\begin{proposition}[Covariance under unitary frame changes]
\label{x_reciprocidad:vii:prop:covariancia-corriente-memoria}
Let \(g_v\) be unitary frame changes independent of the variation
parameter, and \(G_{\mathcal E}=\bigoplus_a g_{t(a)}\).
Under the transformations
\[
 f'_v=g_vf_v,\quad U'_a=g_{t(a)}U_ag_{s(a)}^{-1},\quad
 \mathsf W'_m=G_{\mathcal E}\mathsf W_mG_{\mathcal E}^{-1},\quad
 A'_a=g_{t(a)}A_ag_{t(a)}^{-1},
\]
we have
\begin{equation}
 \mathcal J'_a=g_{t(a)}\mathcal J_ag_{t(a)}^{-1},
 \qquad \mathfrak j'_m(A')=\mathfrak j_m(A).
 \label{x_reciprocidad:vii:eq:covariancia-respuesta}
\end{equation}
\end{proposition}

\begin{proof}
We obtain \(r'=G_{\mathcal E}r\), \(q'=G_{\mathcal E}q\) and
\(v'_a=g_{t(a)}v_a\). Substitution into
\eqref{x_reciprocidad:vii:eq:representante-corriente-discreta} proves the first
identity. Invariance of the Hermitian inner product, or cyclicity of
the trace, proves the second.
\end{proof}

When frame changes depend on the parameter, their derivatives
also act on the field and weight. The corresponding covariant
expression is then the complete differential
\eqref{x_reciprocidad:vii:eq:variacion-completa-memoria}.

\section{Oriented inversion and transport of the weight}
\label{x_reciprocidad:vii:subsec:inversion-respuesta}

Inversion can be seen explicitly in the case of a local energy
per edge, \(\mathsf W_m=\bigoplus_a W_a\).
For \(a:s\to t\), its inverse carries
\[
 U_{\bar a}=U_a^{-1},\qquad
 r_{\bar a}=-U_a^{-1}r_a,\qquad
 W_{\bar a}=U_a^{-1}W_aU_a.
\]
These rules preserve the edge energy exactly.

\begin{proposition}[Differential compatibility with inversion]
\label{x_reciprocidad:vii:prop:inversion-respuesta}
If \(W_a\) and the fields at both endpoints remain fixed, then
\[
 A_{\bar a}=-U_a^{-1}A_aU_a,\qquad
 \dot W_{\bar a}=U_a^{-1}[W_a,A_a]U_a.
\]
The complete variation of the reversed edge's energy coincides with
the original. In particular, if \([W_a,A_a]=0\), the partial transport
covector satisfies
\begin{equation}
 \mathfrak j_{\bar a}(U_a^{-1}A_aU_a)
 =-\mathfrak j_a(A_a).
 \label{x_reciprocidad:vii:eq:imparidad-transporte}
\end{equation}
\end{proposition}

\begin{proof}
The derivative of \(U_a^{-1}\) is \(-U_a^{-1}A_a\), from which
the two formulas result. Writing \(v_a=U_af_s\) and \(f_t=r_a+v_a\),
the part due exclusively to reversed transport equals
\[
 \mathfrak j_{\bar a}(A_{\bar a})
 =-2\operatorname{Re}\langle r_a,W_aA_af_t\rangle
 =\mathfrak j_a(A_a)
   -2\operatorname{Re}\langle r_a,W_aA_ar_a\rangle.
\]
Variation of the weight contributes
\[
 \langle r_{\bar a},\dot W_{\bar a}r_{\bar a}\rangle
 =\langle r_a,[W_a,A_a]r_a\rangle
 =2\operatorname{Re}\langle r_a,W_aA_ar_a\rangle.
\]
The two terms compensate each other. If the commutator is zero, that contribution
disappears; linearity in \(A_a\) gives
\eqref{x_reciprocidad:vii:eq:imparidad-transporte}.
\end{proof}

Energy is summed over one orientation per edge. A sum over both
orientations carries the factor \(1/2\). This convention preserves the same
functional and hence the same response.

\section{Compatibility with refinement}
\label{x_reciprocidad:vii:subsec:refinamiento-respuesta}

\begin{theorem}[Transport of the differential by refinement]
\label{x_reciprocidad:vii:thm:refinamiento-respuesta}
Let \(J_m:\mathcal H_m\to\mathcal H_{m+1}\) and
\(K_m:\mathcal K_m\to\mathcal K_{m+1}\) be fixed maps.
For a family \(C^1\), suppose the refinement identities
\begin{equation}
 D_{m+1}(t)J_m=K_mD_m(t),\qquad
 K_m^*\mathsf W_{m+1}(t)K_m=\mathsf W_m(t).
 \label{x_reciprocidad:vii:eq:familia-refinamiento-memoria}
\end{equation}
Then \(E_{m+1}(J_mf)=E_m(f)\).
The transport covectors agree on variations linked by
the first identity. When the weights vary, their contributions
to the differential also agree.
\end{theorem}

\begin{proof}
By substitution,
\[
 \langle D_{m+1}J_mf,\mathsf W_{m+1}D_{m+1}J_mf\rangle
 =\langle D_mf,K_m^*\mathsf W_{m+1}K_mD_mf\rangle=E_m(f).
\]
Differentiating the first identity of
\eqref{x_reciprocidad:vii:eq:familia-refinamiento-memoria} gives
\(\dot D_{m+1}J_m=K_m\dot D_m\). Consequently,
\[
 2\operatorname{Re}
 \langle D_{m+1}J_mf,\mathsf W_{m+1}\dot D_{m+1}J_mf\rangle
 =2\operatorname{Re}\langle D_mf,\mathsf W_m\dot D_mf\rangle.
\]
The derivative of the second identity is
\(K_m^*\dot{\mathsf W}_{m+1}K_m=\dot{\mathsf W}_m\), proving
agreement of the weight-variation terms.
\end{proof}

The statement uses compatibility of the transport family.
If \(J_m\) depends on the parameter, the chain rule also
preserves the term
\[
 2\operatorname{Re}
 \langle D_{m+1}J_mf,\mathsf W_{m+1}D_{m+1}\dot J_mf\rangle.
\]
The resulting naturality jointly transports the functional, its
incidences and its differential.

\section{Material realization of the response}
\label{x_reciprocidad:vii:subsec:realizacion-respuesta-discreta}

The result of this section is the explicit covector
\(\mathfrak j_m\), with its representative \(\mathcal J_a\).
The material realization composes that covector with the variation of
transport induced by the connection and with the normalization of the action.
If \(\mathcal L_m\) denotes that differential map, the covector on
connection variations is \(\mathcal L_m^*\mathfrak j_m\).
The composition rule
\eqref{x_reciprocidad:vii:eq:variacion-transporte-compuesto} evaluates its part
corresponding to finite products of transports.

In section~\ref{x_reciprocidad:vii:sec:corriente-respuesta-cartan}, the current
\(\sigma_{IJ}\) is a three-form defined by variation of the material
action. Identification between the two realizations preserves the field,
representation, weight, dimensional normalization and passage to the
limit used by that action. These data allow evaluation of the
contraction entering the torsional density. The present
development provides the discrete differential composed in that calculation.

% END OWNER


% BEGIN OWNER /Users/ruben/Documents/New project/output/EXTREMA_MEDIA_RAZON_RECIPROCIDAD_ES_EN_20260918/ARTICULO/ES/source/antecedentes_vii/30c_composicion_corriente_conexion.tex
% Artículo VII, incorporación aditiva, revisión 05.
% Arquitectura autoral: energía de memoria APP--TRIT--TPK, S01.
% Antecedente recuperado: diferencial y refinamiento de 30b.
% Formalización reunida: nota COMPOSICION_VARIACIONAL_TRIT_CONEXION.md.
% Se conserva íntegra 30b; su representante antihermítico es la
% restricción unitaria del diferencial que aquí se desarrolla.
% Esta sección evalúa la corriente de una acción y una realización
% declaradas. No selecciona retrospectivamente su campo ni su peso.

\chapter{Realization of the current in connection coordinates}
\label{x_reciprocidad:vii:sec:composicion-corriente-conexion}

Memory energy assigns each route a response preserving its
incidences and the order of transports. Inserting it into a geometric
action requires its components with respect to connection
variations. This step is obtained by differentiating the transport that already realizes
the route: the connection current is the composition of that differential
with the energy covector of the preceding section.

The elliptic, parabolic and hyperbolic regimes of TRIT require
preserving the complete differential before restricting it to a
particular representation. Unitary rotations constitute a
domain of that construction. The following extension also preserves
the dilation and transport directions admitted by the geometric
realization used.

\section{Differential on invertible transports}

The fibers, fields and weight of
\eqref{x_reciprocidad:vii:eq:energia-memoria-discreta} are retained. We now consider invertible
transports and differentiable families with
\(\dot U_a=A_aU_a\), where \(A_a\) belongs to the tangent space
of the chosen representation. We write
\[
 v_a=U_af_{s(a)},\qquad r_a=f_{t(a)}-v_a,\qquad
 q=\mathsf W_mr,\qquad E_m=\langle r,\mathsf W_mr\rangle.
\]

\begin{proposition}[Complete representative of the differential]
\label{x_reciprocidad:vii:prop:diferencial-invertible}
At fixed field and weight, the following holds
\begin{equation}
 \mathrm d_UE_m(A)=-2\operatorname{Re}\sum_a
       \langle q_a,A_av_a\rangle
 =\sum_a\operatorname{Re}\operatorname{Tr}(\mathcal G_a^*A_a),
 \qquad \mathcal G_a=-2q_av_a^*.
 \label{x_reciprocidad:vii:eq:gradiente-transporte-completo}
\end{equation}
The anti-Hermitian and Hermitian projections of this representative are
\begin{equation}
 \frac{\mathcal G_a-\mathcal G_a^*}{2}
   =v_aq_a^*-q_av_a^*=\mathcal J_a,
 \qquad
 \frac{\mathcal G_a+\mathcal G_a^*}{2}
   =-(q_av_a^*+v_aq_a^*).
 \label{x_reciprocidad:vii:eq:proyecciones-gradiente-memoria}
\end{equation}
Restriction to anti-Hermitian generators recovers exactly
\eqref{x_reciprocidad:vii:eq:representante-corriente-discreta}.
\end{proposition}

\begin{proof}
The identity \(\dot r_a=-A_av_a\) and self-adjointness of the weight give
\(\dot E_m=2\operatorname{Re}\langle q,\dot r\rangle\).
Moreover,
\[
 \operatorname{Tr}(\mathcal G_a^*A_a)
 =-2\operatorname{Tr}(v_aq_a^*A_a)=-2q_a^*A_av_a.
\]
The two projections are obtained by taking the adjoint. For the real
Hilbert--Schmidt inner product, the Hermitian and anti-Hermitian subspaces are
orthogonal; therefore the Hermitian part has zero contraction with
an anti-Hermitian generator.
\end{proof}

The formula allows weights coupling edges: their action is calculated in
\(q=\mathsf W_mr\) before separating the contributions of each
incidence. If field or weight vary, the complete expression
\eqref{x_reciprocidad:vii:eq:variacion-completa-memoria} is retained, also valid for the
invertible transports considered here.

\section{Quadratic representations of TRIT}

In a two-dimensional matrix realization of
\(\mathbb A_\tau\), one may take
\[
 J_{+}=\begin{pmatrix}0&1\\-1&0\end{pmatrix},\qquad
 J_{-}=\begin{pmatrix}0&1\\1&0\end{pmatrix},\qquad
 J_{0}=\begin{pmatrix}0&1\\0&0\end{pmatrix}.
\]
Direct multiplication gives
\[
 J_+^2=-I,\qquad J_-^2=I,\qquad J_0^2=0.
\]
Their exponentials realize, respectively, the elliptic,
hyperbolic and parabolic regimes. The first is orthogonal for the fixed
positive inner product. The second preserves the form
\(B=\operatorname{diag}(-1,1)\), since
\(J_-^{\mathsf T}B+BJ_-=0\). The third represents the algebra
of dual numbers. These representations fix examples of the quadratic
types; their use on a screen of another dimension preserves the
corresponding representation map.

A calculation distinguishes the components of the differential. On an edge,
take \(U=I\), \(f_s=(1,1)^{\mathsf T}\),
\(f_t=2f_s\) and \(\mathsf W=I\). Then \(q=v=f_s\),
\(\mathcal J=0\), whereas
\begin{equation}
 \mathrm dE(J_-)=-4,\qquad \mathrm dE(J_0)=-2.
 \label{x_reciprocidad:vii:eq:control-regimenes-corriente}
\end{equation}
The complete representative preserves both responses. The distinction
identifies the space of variations before performing the material
contraction.

\section{Composition of routes and oriented inversion}

\begin{proposition}[Contragredient transport of the differential]
\label{x_reciprocidad:vii:prop:pullback-invertible}
For \(U_\gamma=U_N\cdots U_1\), let
\(B_k=U_N\cdots U_{k+1}\) and \(B_N=I\). We have
\[
 \dot U_\gamma U_\gamma^{-1}
   =\sum_kB_kA_kB_k^{-1}.
\]
The representative \(\mathcal G_\gamma\) in the final fiber induces
at incidence \(k\)
\begin{equation}
 \mathcal G_k=B_k^*\mathcal G_\gamma(B_k^{-1})^*.
 \label{x_reciprocidad:vii:eq:transporte-contragrediente}
\end{equation}
\end{proposition}

\begin{proof}
The product rule, followed by multiplication by
\(U_\gamma^{-1}\), proves the first equality. Cyclicity of the
trace gives
\[
 \operatorname{Re}\operatorname{Tr}
  (\mathcal G_\gamma^*B_kA_kB_k^{-1})
 =\operatorname{Re}\operatorname{Tr}
  ((B_k^*\mathcal G_\gamma(B_k^{-1})^*)^*A_k).
\]
For unitary \(B_k\), the conjugation of
\ref{x_reciprocidad:vii:prop:composicion-respuesta} is recovered.
\end{proof}

\begin{proposition}[Inversion by energy congruence]
\label{x_reciprocidad:vii:prop:inversion-congruencia}
For a local energy per edge, inversion is represented by
\begin{equation}
 U_{\bar a}=U_a^{-1},\qquad
 r_{\bar a}=-U_a^{-1}r_a,\qquad
 W_{\bar a}=U_a^*W_aU_a.
 \label{x_reciprocidad:vii:eq:inversion-peso-general}
\end{equation}
It preserves energy. If \(W_a\) remains fixed,
\[
 A_{\bar a}=-U_a^{-1}A_aU_a,\qquad
 \dot W_{\bar a}=U_a^*(A_a^*W_a+W_aA_a)U_a.
\]
The complete variation of the inverted energy coincides with the original.
\end{proposition}

\begin{proof}
Substitution into \(r_{\bar a}^*W_{\bar a}r_{\bar a}\)
produces \(r_a^*W_ar_a\). The derivative of the inverse is
\(\dot U_a^{-1}=-U_a^{-1}A_a\); that of the weight is obtained by
the product rule. Equality of energies holds for the entire
family, so their derivatives coincide, including the term
\(\langle r_{\bar a},\dot W_{\bar a}r_{\bar a}\rangle\).
\end{proof}

Congruence preserves positivity for every invertible transport.
In the unitary representation it reduces to the inversion rule of
\ref{x_reciprocidad:vii:prop:inversion-respuesta}. The sum uses one orientation
per edge, or half the sum over both orientations.

\section{Evaluation of current components}

Let \(\theta=(\theta_b)\) be a real variation of the connection
coordinates of the realization used. The effective differential of
transport defines the operators
\begin{equation}
 A_a(\theta)=\dot U_aU_a^{-1}
     =\sum_bB_{ab}\theta_b.
 \label{x_reciprocidad:vii:eq:diferencial-conexion-transporte}
\end{equation}
For a finite composition, their contributions are obtained through
\eqref{x_reciprocidad:vii:eq:transporte-contragrediente}. The operators
\(B_{ab}\) are the derivatives of the connection representation that
realizes the route.

\begin{theorem}[Connection current of the memory action]
\label{x_reciprocidad:vii:thm:corriente-conexion-evaluada}
For a material term \(S_{\rm mem}=\mathfrak a_m E_m\),
with \(\mathfrak a_m\) its normalization on the action line,
at fixed field, weight and normalization we have
\begin{equation}
 \delta S_{\rm mem}=-\frac12\sum_b\theta_b s_b,\qquad
 s_b=4\mathfrak a_m\operatorname{Re}
        \sum_a\langle q_a,B_{ab}v_a\rangle.
 \label{x_reciprocidad:vii:eq:componentes-corriente-conexion}
\end{equation}
If the cellular pairing between the connection and current bases
is an invertible matrix \(M\), defined by
\(\delta S_{\rm mem}=-\tfrac12\theta^{\mathsf T}M\sigma\),
the current coordinates are
\begin{equation}
 \sigma=M^{-1}s.
 \label{x_reciprocidad:vii:eq:coordenadas-corriente-material}
\end{equation}
\end{theorem}

\begin{proof}
Substitute \eqref{x_reciprocidad:vii:eq:diferencial-conexion-transporte} into
\eqref{x_reciprocidad:vii:eq:gradiente-transporte-completo}, multiply by
\(\mathfrak a_m\), and collect the coefficients of
\(\theta_b\). The variational convention \(-1/2\) fixes the
factor \(4\). Comparing both expressions for every variation
\(\theta\) gives \(M\sigma=s\); invertibility of the pairing
proves existence and uniqueness of these coordinates.
\end{proof}

The formula explicitly realizes the composition mentioned in
\ref{x_reciprocidad:vii:subsec:realizacion-respuesta-discreta}. To evaluate each
component, the field is transported, the residual is calculated, the
weight is applied and the result is contracted with the connection variation of the same realization.
The pairing \(M\) preserves the orientation and volume of cells.
A diagonal pairing gives the corresponding division by their
volumes; other bases preserve the complete matrix.

If the connection also enters \(f\), \(\mathsf W_m\)
or \(\mathfrak a_m\), we apply
\eqref{x_reciprocidad:vii:eq:variacion-completa-memoria} and add
\(E_m\delta\mathfrak a_m\). Contributions from all
material terms are summed before composing the Cartan response.

\begin{corollary}[Compatibility of the current with refinement]
\label{x_reciprocidad:vii:cor:refinamiento-corriente-conexion}
For the compatible families of
\ref{x_reciprocidad:vii:thm:refinamiento-respuesta}, let \(P_m\) be the fixed map
taking connection variations from level \(m\) to level
\(m+1\). If the memory actions coincide on those refined fields
and transports, their covectors satisfy
\[
 s_m=P_m^{\mathsf T}s_{m+1},\qquad
 M_m\sigma_m=P_m^{\mathsf T}M_{m+1}\sigma_{m+1}.
\]
\end{corollary}

\begin{proof}
Differentiating the equality of actions in the direction
\(P_m\theta\) gives
\(-\tfrac12\theta^{\mathsf T}s_m
 =-\tfrac12\theta^{\mathsf T}P_m^{\mathsf T}s_{m+1}\).
The arbitrariness of \(\theta\) and the cellular representations
of both covectors prove the two identities. Composition of
these maps iterates the equality to any finite number of refinements.
\end{proof}

\section{Insertion into the Cartan equation}

Section~\ref{x_reciprocidad:vii:sec:corriente-respuesta-cartan} normalizes the
material three-form through
\(\delta_\omega S_m=-\tfrac12\int\delta\omega^{IJ}\wedge
\sigma_{IJ}\). The preceding theorem provides its cellular
components for the declared memory action and pairing.
The Cartan response uses that current together with the co-tetrad,
the Barbero--Immirzi parameter and the gravitational coupling of the
same realization.

The dependence of matter on contorsion determines how this equation is
solved. For affine matter, the current is independent of contorsion,
and the linear inversion and quadratic elimination proved in the next
section apply. For a holonomy action, the differential is evaluated at
the current connection, and that dependence is retained when solving
the equation. For example, an edge with
\(f_s=f_t=1\), unit weight, and \(U(t)=e^{it}\) has
\[
 E(t)=2-2\cos t,\qquad E'(t)=2\sin t.
\]
Here the current depends on \(t\). The complete formula retains
that dependence instead of replacing it by its value at a particular
connection.

Finally, the gravitational density is obtained from the metric
variation of the effective action, with its field and normalization.
The representation of the current, the contortion it produces and the
resulting density are the successive operations of that composition.
The subsequent dilution law transports the amplitude thus evaluated in
the corresponding material section.

% END OWNER


% BEGIN OWNER /Users/ruben/Documents/New project/output/EXTREMA_MEDIA_RAZON_RECIPROCIDAD_ES_EN_20260918/ARTICULO/ES/source/antecedentes_vii/30e_corriente_extension_minima.tex
% Artículo VII. Incorporación aditiva, 10-09-2026.
% RESULTADO_RECUPERADO: fuentes_conservadas/18_clausura_energetica.tex,
% líneas 357--485: Gramiano, peso, extensión mínima y eliminación de Schur;
% líneas 149--180 y 235--295: normalización y unicidad de la respuesta.
% Residencia anterior en VII: 30_pantalla_y_respuesta.tex,
% vii:prop:incidencia, vii:eq:lambda-pantalla, vii:prop:minima-extension.
% FORMALIZACION_NUEVA / CERTIFICADO_NUEVO: diferencial del mínimo con
% incidencia, frontera y normalización variables; incidencia transportada
% por bloques y composición con coordenadas de conexión.
% La familia de rutas que realiza cada incidencia se mantiene explícita.
% No se atribuyen al propietario histórico ni esa familia diferenciable
% completa ni la evaluación de una amplitud material de espín.
% Dependencias: 30, 30c y la convención variacional de 31. No modifica
% esos fragmentos ni constituye por sí solo una realización material.

\chapter{Differential of the minimal boundary extension}
\label{x_reciprocidad:vii:sec:corriente-extension-minima}

The cubic incidence simultaneously determines the minimal-extension field and the weight of its response. For a boundary datum, both are calculated through the inverse of the incidence Gramian. We now develop the differential of this same operation, preserving the variations of the boundary, transport and normalization scale. In this way, the field involved in the response is selected by the screen variational problem.

\section{Minimizing field and response weight}
\label{x_reciprocidad:vii:subsec:extension-minima-parametrizada}

Let \(H\) and \(Q\) be finite-dimensional Hermitian spaces, with
fixed inner products antilinear in the first variable. The real case is
obtained by replacing adjoints with transposes. On an open set of
real parameters \(\Omega\), consider families of class \(C^1\)
\[
 \mathsf C(\theta):H\longrightarrow Q,\qquad
 b(\theta)\in Q,\qquad \chi(\theta)>0,
 \qquad \operatorname{Ran}\mathsf C(\theta)=Q.
\]
Surjectivity is equivalent to full row rank in orthonormal bases. We define the Gramian, the conjugate potential and the extension
\begin{equation}
 \mathsf L=\mathsf C\mathsf C^*,\qquad
 y=\mathsf L^{-1}b,\qquad f_{\min}=\mathsf C^*y,
 \qquad \Lambda=\chi\mathsf L^{-1}.
 \label{x_reciprocidad:vii:eq:gramiano-extension-minima}
\end{equation}
The incidence \(\mathsf C\) and the cellular pairing \(M\) of \eqref{x_reciprocidad:vii:eq:coordenadas-corriente-material} are different operators.

\begin{proposition}[Variational selection of the extension]
\label{x_reciprocidad:vii:prop:minimo-parametrizado}
The extension \(f_{\min}\) is the unique norm minimizer among the fields \(f\in H\) satisfying \(\mathsf C f=b\). Its energy is
\begin{equation}
 E(\theta)=\chi\min_{\mathsf C f=b}\|f\|^2
 =\chi\langle b,\mathsf L^{-1}b\rangle
 =\langle b,\Lambda b\rangle.
 \label{x_reciprocidad:vii:eq:energia-extension-minima}
\end{equation}
The field \(f_{\min}\), potential \(y\), and energy \(E\)
depend \(C^1\)-smoothly on the parameters.
\end{proposition}

\begin{proof}
Surjectivity of \(\mathsf C\) makes \(\mathsf C^*\) injective,
so \(\langle z,\mathsf Lz\rangle=\|\mathsf C^*z\|^2>0\)
for \(z\ne0\). Thus \(\mathsf L\) is invertible and
\(\mathsf C f_{\min}=\mathsf L y=b\). Every solution is
\(f_{\min}+z\), with \(z\in\ker\mathsf C\). Since
\(f_{\min}\in\operatorname{Ran}\mathsf C^*=(\ker\mathsf C)^\perp\),
\[
 \|f_{\min}+z\|^2=\|f_{\min}\|^2+\|z\|^2,
 \qquad
 \|f_{\min}\|^2=\langle y,\mathsf L y\rangle
 =\langle b,\mathsf L^{-1}b\rangle.
\]
These identities prove the minimum, its uniqueness and its value. Inversion is a differentiable operation on the open set of invertible matrices; the expressions in \eqref{x_reciprocidad:vii:eq:gramiano-extension-minima} provide the asserted regularity.
\end{proof}

The cubic realization of \ref{x_reciprocidad:vii:prop:incidencia} corresponds to \(H=\mathbb R^{\{-1,1\}^3}\), the screen quotient \(Q=P_\square\mathbb R^6\) and the incidence \(\mathsf C=M:H\to Q\). Its Gramian and normalization are
\begin{equation}
 \mathsf L_\square=12P_u+4P_-,\qquad
 \chi_m=\frac{112}{81\ell_m^2},\qquad
 \Lambda_m=\frac{28}{243\ell_m^2}(P_u+3P_-).
 \label{x_reciprocidad:vii:eq:normalizacion-minimo-cubico}
\end{equation}
The factor \(112/81\) comes from the soldering and the quadratic compression defect of \eqref{x_reciprocidad:vii:eq:lambda-pantalla}. If \(b\) has length dimension, so do \(y\) and \(f_{\min}\), the incidence is dimensionless and \(E\) is dimensionless. Conversion to action uses the normalization in the action line specified below.

\Needspace{15\baselineskip}
\section{Full differential of the minimum}
\label{x_reciprocidad:vii:subsec:diferencial-minimo}

\begin{theorem}[Response to the variation of incidence and boundary]
\label{x_reciprocidad:vii:thm:diferencial-minimo}
For every real tangent direction \(\delta\theta\), one has
\begin{equation}
 \delta E
 =\delta\chi\,\langle b,y\rangle
   +2\chi\operatorname{Re}
       \langle y,\delta b-(\delta\mathsf C)f_{\min}\rangle.
 \label{x_reciprocidad:vii:eq:diferencial-minimo-completo}
\end{equation}
If \(\chi=112/(81\ell_m^2)\), with variable \(\ell_m>0\),
\begin{equation}
 \delta E=-2E\frac{\delta\ell_m}{\ell_m}
 +2\chi_m\operatorname{Re}
       \langle y,\delta b-(\delta\mathsf C)f_{\min}\rangle.
 \label{x_reciprocidad:vii:eq:diferencial-minimo-escala}
\end{equation}
\end{theorem}

\begin{proof}
Differentiating \(\mathsf L\mathsf L^{-1}=I_Q\) gives
\[
 \delta\mathsf L^{-1}
 =-\mathsf L^{-1}(\delta\mathsf L)\mathsf L^{-1},\qquad
 \delta\mathsf L=(\delta\mathsf C)\mathsf C^*
                      +\mathsf C(\delta\mathsf C)^*.
\]
Applying the product rule to \eqref{x_reciprocidad:vii:eq:energia-extension-minima} and using the self-adjointness of \(\mathsf L^{-1}\) yields
\[
 \delta E=\delta\chi\,\langle b,y\rangle
 +2\chi\operatorname{Re}\langle y,\delta b\rangle
 -\chi\langle y,(\delta\mathsf L)y\rangle.
\]
The two terms of the last product are conjugate to each other and \(\mathsf C^*y=f_{\min}\); therefore,
\[
 \langle y,(\delta\mathsf L)y\rangle
 =2\operatorname{Re}\langle y,(\delta\mathsf C)f_{\min}\rangle.
\]
One obtains \eqref{x_reciprocidad:vii:eq:diferencial-minimo-completo}. Finally, \(\delta\chi_m=-2\chi_m\delta\ell_m/\ell_m\) produces \eqref{x_reciprocidad:vii:eq:diferencial-minimo-escala}.
\end{proof}

The formula incorporates the variation of the minimizing field through the equation that determines it. It is also checked directly: \(\mathsf C\delta f_{\min}=\delta b-(\delta\mathsf C)f_{\min}\) and \(f_{\min}=\mathsf C^*y\) imply
\[
 \delta\|f_{\min}\|^2
 =2\operatorname{Re}\langle f_{\min},\delta f_{\min}\rangle
 =2\operatorname{Re}
     \langle y,\delta b-(\delta\mathsf C)f_{\min}\rangle.
\]
For example, if \(\mathsf C\) is fixed and \(b=\ell_m b_0\), with \(b_0\) fixed, the two terms of \eqref{x_reciprocidad:vii:eq:diferencial-minimo-escala} cancel and the energy remains constant. This check jointly preserves the boundary dimension and its normalization.

\section{Fixed quotient and frame changes}
\label{x_reciprocidad:vii:subsec:marcos-minimo}

The preceding derivatives are taken in fixed spaces \(H,Q\) and Hermitian
inner products. In the six-face presentation, the inverse always
acts on \(Q=P_\square H_F\). For a family of quotients
\(Q(\theta)\), an isometric trivialization
\(V_Q(\theta):Q\to Q(\theta)\) is fixed first, and likewise for the domain.
The derivatives of those maps are included when differentiating the expressions
transported to the fixed spaces. The formulas apply on each
open set where rank is preserved. Changes in the energy inner product
additionally require differentiating that product; a Lorentzian metric
and the positive inner product used for minimization retain their distinct
types here.

\begin{proposition}[Covariance of the minimum]
\label{x_reciprocidad:vii:prop:covariancia-minimo}
Let \(g_H,g_Q\) be unitary. Under
\[
 \mathsf C'=g_Q\mathsf C g_H^*,\qquad b'=g_Qb,\qquad \chi'=\chi,
\]
\(y'=g_Qy\), \(f'_{\min}=g_Hf_{\min}\) and \(E'=E\) hold. When the frame changes depend on the parameter, the full differential preserves this equality. In particular, a pure frame variation has \(\delta E=0\).
\end{proposition}

\begin{proof}
The transformed Gramian is \(\mathsf L'=g_Q\mathsf Lg_Q^*\), from which the transformations of \(y,f_{\min}\) and equality of energies are obtained. For a pure frame variation, at a point with identity frames, let \(A_Q^*=-A_Q\) and \(A_H^*=-A_H\) be their generators. Then
\[
 \delta\mathsf C=A_Q\mathsf C-\mathsf C A_H,\qquad
 \delta b=A_Qb,\qquad
 \delta b-(\delta\mathsf C)f_{\min}=\mathsf C A_Hf_{\min}.
\]
Its contribution to \eqref{x_reciprocidad:vii:eq:diferencial-minimo-completo} is \(2\chi\operatorname{Re}\langle f_{\min},A_Hf_{\min}\rangle=0\). The proof thus includes the terms due to the moving frames.
\end{proof}

\section{Transported incidence and connection components}
\label{x_reciprocidad:vii:subsec:incidencia-conexion-minima}

The connection dependence is specified before performing the contraction. An explicit blockwise realization of the incidence is obtained as follows. Let \(\mathcal K\) be a common Hermitian fibre, \(F=\{(i,\epsilon):1\leq i\leq3,\ \epsilon=\pm1\}\) and \(V=\{-1,1\}^3\). For each incidence \(\nu_i=\epsilon\), the realization used declares a route \(\gamma_{i,\epsilon,\nu}\) from the vertex fibre to the face fibre, with transport \(U_{i,\epsilon,\nu}\). The trivializations identify these fibres with \(\mathcal K\). We define
\begin{equation}
 \begin{split}
 (\mathsf B_Uf)_{i,\epsilon}
   &=\sum_{\nu_i=\epsilon}U_{i,\epsilon,\nu}f_\nu,\\
 \mathsf C_U&=(q_\square\otimes I_{\mathcal K})\mathsf B_U:
   \bigoplus_{\nu\in V}\mathcal K\longrightarrow Q\otimes\mathcal K,
 \qquad q_\square=P_\square:H_F\to Q.
 \end{split}
 \label{x_reciprocidad:vii:eq:incidencia-transportada-minima}
\end{equation}
This is an explicit variational continuation of the cubic incidence in the declared route representation. For identity transports, \(\mathsf C_U=M\otimes I_{\mathcal K}\) and one recovers \eqref{x_reciprocidad:vii:eq:normalizacion-minimo-cubico}, tensored with the fibre. The Gramian remains positive in a neighbourhood of that realization, since positive definiteness is an open condition. On each full-rank domain, the weight and field are fixed by \eqref{x_reciprocidad:vii:eq:gramiano-extension-minima}. The routes realizing the incidences, their orientations and their memories are data of this realization of APP--TRIT--TPK transport and accompany the construction.

If \(\omega=(\omega_j)\) are real connection coordinates, the derivative of \(\mathsf C_U\) is calculated by substituting in each block
\begin{equation}
 \partial_j U_\gamma
 =\sum_{k=1}^{N}
    U_N\cdots U_{k+1}(\partial_jU_k)U_{k-1}\cdots U_1,
 \qquad U_\gamma=U_N\cdots U_1.
 \label{x_reciprocidad:vii:eq:derivada-incidencia-rutas}
\end{equation}
Empty products are identities. The formula follows from the product rule and preserves the position of each operation in the route. Thus, \(\mathsf C_j:=\partial_j\mathsf C_U\) is obtained from a finite sum of known matrices in the chosen realization, followed by the projection \(q_\square\otimes I\). If a moving quotient is used, its derivative is incorporated according to \ref{x_reciprocidad:vii:subsec:marcos-minimo}. The fixed integer cubic matrix corresponds to the case \(\mathsf C_j=0\); the response to the connection follows from the transport dependence actually declared in \eqref{x_reciprocidad:vii:eq:incidencia-transportada-minima}.

% REMATE_LOCAL_X_EXTENSION_MINIMA
\Needspace{34\baselineskip}
\begin{theorem}[Variational components of the minimal extension]
\label{x_reciprocidad:vii:thm:componentes-minimo-conexion}
For the screen action \(S_{\min}=\mathfrak a_m E\), where \(\mathfrak a_m\) is its normalization in the action line, let us write
\[
 b_j=\partial_jb,\qquad \chi_j=\partial_j\chi,
 \qquad a_j=\partial_j\mathfrak a_m.
\]
With the convention \(\delta S_{\min}=-\tfrac12\sum_j s_j\delta\omega_j\),
the components are
\begin{equation}
 \begin{split}
 s_j={}&4\mathfrak a_m\chi\operatorname{Re}
             \langle y,\mathsf C_jf_{\min}-b_j\rangle\\
      &-2\mathfrak a_m\chi_j\langle b,y\rangle-2a_jE.
 \end{split}
 \label{x_reciprocidad:vii:eq:componentes-minimo-conexion}
\end{equation}
For the cubic normalization, the second term equals \(4\mathfrak a_m E\,\partial_j\ell_m/\ell_m\). If the cellular pairing of this same action has invertible matrix \(\mathsf M_{\rm cel}\), its current coordinates are \(\sigma=\mathsf M_{\rm cel}^{-1}s\).
\end{theorem}

\begin{proof}
The product rule gives \(\partial_jS_{\min}=a_jE+\mathfrak a_m\partial_jE\). Substituting \eqref{x_reciprocidad:vii:eq:diferencial-minimo-completo} and multiplying by \(-2\) proves \eqref{x_reciprocidad:vii:eq:componentes-minimo-conexion}. The derivative of \(\chi_m\) produces the scale expression. Finally, comparing \(\delta S_{\min}=-\tfrac12\delta\omega^{\mathsf T}s\) with \(-\tfrac12\delta\omega^{\mathsf T}\mathsf M_{\rm cel}\sigma\) for every variation gives \(\mathsf M_{\rm cel}\sigma=s\), and its inverse determines the coordinates.
\end{proof}

The construction recovers a field and a weight jointly selected
by the incidence, and evaluates their response to variations of the same
realization. The boundary datum \(b\), the route family, the scale
\(\ell_m\), the normalization \(\mathfrak a_m\) and the cellular
pairing retain their publication rules. Interpretation as a
material current uses those objects and the complete action. The
Cartan response and its subsequent metric density are composed with
it in the corresponding matter domain, preserving any
dependence of the current on the connection.

% END OWNER

\begin{HMTCompactSection}

% BEGIN OWNER /Users/ruben/Documents/New project/output/EXTREMA_MEDIA_RAZON_RECIPROCIDAD_ES_EN_20260918/ARTICULO/ES/source/antecedentes_vii/31_corriente_y_respuesta_cartan.tex
% Artículo VII. Fragmento de cuerpo: corriente y respuesta Cartan--Holst.
% RESULTADO_RECUPERADO: acción, corriente variacional, inversa de Holst,
% torsión algebraica, respuesta cuadrática y normalización métrica.
% Fuentes leídas íntegramente:
% [II10] output/ARTICULO_II_REV10_EDICION_INTEGRADA_20260910/manuscrito/
%        sections/10_gravedad.tex, líneas 109--155.
% [G11c] integral activo 2249, fuente/manuscrito/sections/md/
%        11c_gravedad_holonomia_rev6.tex, líneas 214--334.
% [B028] integral activo, fuente/manuscrito/integracion_83/parte_iv/body/
%        B028_ch68_realizacion_cartan_holst_barbero_7df2afd9f675_11_gravedad_cartan_holst.tex,
%        líneas 76--218; copia idéntica conservada en II10/fuentes/propietarios_II/G03/.
%        SHA-256 e1e1f3325399b3ef09c331ed8268e3883f7bcc561cd7a560cf80d1be7eb1e91c.
% CERTIFICADO_NUEVO de resultados recuperados: inversión explícita del mapa
% torsión -> D(e wedge e), prueba de unicidad de contorsión y eliminación
% homogénea del término cuadrático. Las dos inversas se comprobaron sobre
% las 24 bases correspondientes mediante aritmética racional exacta.
% Precisión de dominio: la respuesta lineal en corriente y su eliminación
% cuadrática se enuncian para materia afín en la contorsión. Se conserva
% la dependencia del coeficiente respecto de la acción material concreta.
%
% DEPENDENCIAS DE INTEGRACIÓN (no constituyen remisiones bibliográficas):
% 1. vii:sec:realizacion-cartan: cotetrada, conexión y realización del
%    transporte enriquecido; incluir su prueba en el cuerpo precedente.
% 2. vii:sec:acoplamientos: genealogía de gamma_HMT, G_int, c_int,
%    homogeneidad dimensional y elección de la carta, con pruebas.
% 3. Para evaluar una especie material: acción S_m concreta, representación
%    de sus campos y cálculo de su corriente sigma. Este fragmento recibe
%    esa acción declarada y demuestra su respuesta variacional.
% 4. La identificación de una corriente de pantalla con sigma, su
%    contracción física y su normalización material tienen residencia propia.
%    Aquí no se introduce ni se selecciona una amplitud s_0.
% 5. El balance métrico utiliza la covariancia de la acción material y
%    sus ecuaciones de movimiento; no supone conservación separada del
%    tensor de espín y del tensor material de Levi--Civita.

\chapter{Spin current and algebraic Cartan--Holst response}
\label{x_reciprocidad:vii:sec:corriente-respuesta-cartan}

The realization of the enriched transport in Section~\ref{x_reciprocidad:vii:sec:realizacion-cartan} provides the cotetrad and the connection; the dimensional sections and the incidence parameter of Sections~\ref{x_reciprocidad:v:ant:sec:barbero} and~\ref{x_reciprocidad:v:ant:sec:gravity} fix their couplings. Their composition determines a first-order variational problem. The material current occupies a precise position in it: it is the derivative of the action with respect to the connection and determines the torsional response with the same normalization.

\section{Geometric domain and current normalization}
\label{x_reciprocidad:vii:subsec:dominio-corriente-cartan}

We work on a smooth oriented four-dimensional manifold \(M\), with internal Lorentzian bundle \(E\), metric \(\eta=\operatorname{diag}(-1,1,1,1)\) and non-degenerate cotetrad \(e\in\Omega^1(M;E)\). In a local frame we write \(e^I\), with \(I=0,1,2,3\). The metric connection \(\omega^{IJ}=-\omega^{JI}\) acts through \(D_\omega\). When the fields \(\Psi\) are spinorial, a spin structure and the corresponding material representation are preserved. We define
\begin{equation}
 \Sigma^{IJ}=e^I\wedge e^J,\qquad
 F^{IJ}=\mathrm d\omega^{IJ}+\omega^I{}_{K}\wedge\omega^{KJ},
 \qquad T^I=D_\omega e^I.
 \label{x_reciprocidad:vii:eq:campos-cartan}
\end{equation}
The operator \(\star\) acts on the internal bivector indices,
\begin{equation}
 (\star X)^{IJ}=\tfrac12\epsilon^{IJ}{}_{KL}X^{KL},
 \qquad \star^2=-I,\qquad
 \mathsf P_\gamma=\star+\gamma^{-1}I.
 \label{x_reciprocidad:vii:eq:operador-holst}
\end{equation}
This distinguishes this duality from the exterior duality of forms on \(M\). Internal indices are raised and contracted with \(\eta\).

We fix a chart with \(G_{\rm int}>0\), \(c_{\rm int}>0\),
\(\gamma=\gamma_{\rm HMT}\in\mathbb R\setminus\{0\}\) and
\(\Lambda_{\rm int}\) constant on it. We abbreviate
\(\kappa=8\pi G_{\rm int}/c_{\rm int}^4\) and \(c=c_{\rm int}\).
These sections remain fixed during variation. The action is
\begin{equation}
 \begin{split}
 S[e,\omega,\Psi]={}&
 \frac1{2\kappa c}\int_M\Sigma_{IJ}\wedge
                         (\mathsf P_\gamma F)^{IJ}
 -\frac{\Lambda_{\rm int}}{\kappa c}
                         \int_M\operatorname{vol}_e
 +S_{\rm m}[e,\omega,\Psi].
 \end{split}
 \label{x_reciprocidad:vii:eq:accion-cartan-holst}
\end{equation}
The volume and orientation convention is the one that gives \(\Sigma_{IJ}\wedge(\star F)^{IJ}=R\operatorname{vol}_e\) in the Levi--Civita sector. The prior dimensional type makes each summand an action section.

For the variation of the connection, variations with compact support in the interior are used, or \(\delta\omega|_{\partial M}=0\) is fixed. A boundary completion whose variational problem cancels the same boundary term may also be used. At fixed cotetrad and material fields, the spin current is the degree-three form, valued in dual internal bivectors, determined by
\begin{equation}
 \delta_\omega S_{\rm m}
 =-\tfrac12\int_M\delta\omega^{IJ}\wedge\sigma_{IJ}.
 \label{x_reciprocidad:vii:eq:corriente-variacional}
\end{equation}
Pairing with an arbitrary compactly supported variation fixes \(\sigma\) uniquely. This definition preserves the action unit of the current and the normalization factor involved in its geometric response.

\section{Connection variation and Holst inversion}
\label{x_reciprocidad:vii:subsec:variacion-holst}

\begin{theorem}[Variational Cartan--Holst equation]
\label{x_reciprocidad:vii:thm:ecuacion-cartan-holst}
On the preceding domain, stationarity with respect to the connection is equivalent to
\begin{equation}
 D_\omega(\mathsf P_\gamma\Sigma)^{IJ}
       =\kappa c\,\sigma^{IJ}.
 \label{x_reciprocidad:vii:eq:ecuacion-cartan-holst}
\end{equation}
\end{theorem}
\begin{proof}
The variation of the curvature is \(\delta F=D_\omega\delta\omega\). The internal duality is parallel and self-adjoint for the bivector contraction; since \(\gamma\) is constant, \(D_\omega\mathsf P_\gamma=\mathsf P_\gamma D_\omega\). Setting \(B_{IJ}=(\mathsf P_\gamma\Sigma)_{IJ}\) yields
\[
 B_{IJ}\wedge D_\omega\delta\omega^{IJ}
 =\mathrm d(B_{IJ}\wedge\delta\omega^{IJ})
   -D_\omega B_{IJ}\wedge\delta\omega^{IJ}.
\]
The last form changes sign when its factors of degrees three and one are interchanged. The boundary conditions eliminate the integral of the exact term. The cosmological term is independent of \(\omega\). Therefore,
\[
 \delta_\omega S=
 \frac1{2\kappa c}\int_M\delta\omega^{IJ}\wedge
       \bigl[D_\omega B_{IJ}-\kappa c\,\sigma_{IJ}\bigr].
\]
The arbitrariness of the variation proves the equation, including its
sign and coefficient.
\end{proof}

\begin{lemma}[Real inverse of the Holst operator]
\label{x_reciprocidad:vii:lem:inversa-holst}
For \(\gamma\in\mathbb R\setminus\{0\}\),
\begin{equation}
 \mathsf P_\gamma^{-1}
 =\frac{\gamma^2}{1+\gamma^2}
                      (-\star+\gamma^{-1}I).
 \label{x_reciprocidad:vii:eq:inversa-holst}
\end{equation}
\end{lemma}
\begin{proof}
Both factors are polynomials in \(\star\) and commute. Their
product before normalization is
\[
 (\star+\gamma^{-1}I)(-\star+\gamma^{-1}I)
 =-\star^2+\gamma^{-2}I=(1+\gamma^{-2})I.
\]
The denominator \(1+\gamma^2\) is positive on the fixed real domain. Multiplication by \(\gamma^2/(1+\gamma^2)\) gives the identity.
\end{proof}

The current thus determines the degree-three bivector-valued form
\begin{equation}
 Y^{IJ}:=\kappa c\,
   \frac{\gamma^2}{1+\gamma^2}
   (-\star+\gamma^{-1}I)\sigma^{IJ},
 \qquad D_\omega\Sigma^{IJ}=Y^{IJ}.
 \label{x_reciprocidad:vii:eq:corriente-dualizada}
\end{equation}

\section{Pointwise reconstruction of torsion and contorsion}
\label{x_reciprocidad:vii:subsec:reconstruccion-torsion}

The Leibniz rule gives
\begin{equation}
 D_\omega\Sigma^{IJ}
 =T^I\wedge e^J-e^I\wedge T^J
 =:\mathsf L_e(T)^{IJ}.
 \label{x_reciprocidad:vii:eq:mapa-torsion}
\end{equation}
At each point, \(\mathsf L_e\) acts from \(\Lambda^2T^*M\otimes E\) to \(\Lambda^3T^*M\otimes\Lambda^2E\). Both spaces have dimension twenty-four. Let \(E_I\) be the dual frame of the cotetrad, so that \(\iota_{E_I}e^J=\delta_I^J\).

\begin{theorem}[Algebraic inversion of the torsional equation]
\label{x_reciprocidad:vii:thm:inversion-torsion}
For a nondegenerate co-tetrad, \(\mathsf L_e\) is an isomorphism.
Its inverse on the form \(Y\) is written
\begin{equation}
 B^I=\sum_J\iota_{E_J}Y^{IJ},\qquad
 b=\tfrac14\sum_I\iota_{E_I}B^I,\qquad
 T^I=B^I-e^I\wedge b.
 \label{x_reciprocidad:vii:eq:torsion-explicita}
\end{equation}
Equation~\eqref{x_reciprocidad:vii:eq:ecuacion-cartan-holst} therefore determines
torsion pointwise from \(e\) and the current appearing in it.
\end{theorem}
\begin{proof}
First let \(Y=\mathsf L_e(T)\) and set \(t=\sum_J\iota_{E_J}T^J\). The contraction of \eqref{x_reciprocidad:vii:eq:mapa-torsion}, using that \(T^I\) has degree two and that the cotetrad has four elements, gives
\[
 \sum_J\iota_{E_J}(T^I\wedge e^J)=2T^I,
 \qquad
 \sum_J\iota_{E_J}(e^I\wedge T^J)=T^I-e^I\wedge t.
\]
Hence \(B^I=T^I+e^I\wedge t\). A second contraction produces
\[
 \sum_I\iota_{E_I}B^I=t+3t=4t.
\]
One recovers \(t=b\) and then \(T^I=B^I-e^I\wedge b\). In particular, the kernel of \(\mathsf L_e\) is zero. Equality of dimensions proves its surjectivity and makes the formula valid for every form \(Y\) in the codomain.
\end{proof}

The connection decomposes uniquely as
\begin{equation}
 \omega^{IJ}=\mathring\omega^{IJ}(e)+K^{IJ},\qquad
 K^{IJ}=-K^{JI},\qquad T^I=K^I{}_{J}\wedge e^J,
 \label{x_reciprocidad:vii:eq:contorsion}
\end{equation}
where \(\mathring\omega(e)\) is the Levi--Civita connection and \(K\) the contorsion. To make the inverse explicit, we write \(K_{IJ}=K_{IJK}e^K\) and \(T_I=\tfrac12T_{IJK}e^J\wedge e^K\). Then
\begin{equation}
 T_{IJK}=K_{IKJ}-K_{IJK},\qquad
 K_{IJK}=\tfrac12(T_{KIJ}-T_{IJK}-T_{JKI}).
 \label{x_reciprocidad:vii:eq:contorsion-componentes}
\end{equation}
The second equality is obtained by combining the first cyclically and using \(K_{IJK}=-K_{JIK}\); substituting it recovers the components of \(T\). Torsion and contorsion are therefore two equivalent presentations of the algebraic response at fixed cotetrad.

\begin{corollary}[Current-free sector and minimal linear response]
\label{x_reciprocidad:vii:cor:desacoplamiento-torsional}
If \(\sigma=0\), one obtains \(T=0\) and \(\omega=\mathring\omega(e)\). For a material coupling whose current \(\sigma[e,\Psi]\) is independent of \(K\), equations~\eqref{x_reciprocidad:vii:eq:corriente-dualizada}, \eqref{x_reciprocidad:vii:eq:torsion-explicita} and \eqref{x_reciprocidad:vii:eq:contorsion-componentes} provide a unique response \(K_*(e,\Psi)\), linear in \(\sigma\).
\end{corollary}
\begin{proof}
The three reconstruction maps are linear with \(e\), \(\gamma\) and \(\kappa c\) fixed. The zero current produces \(Y=0\), then \(T=0\) and finally \(K=0\).
\end{proof}

In the minimal action, the connection equation contains this pointwise response. Its evolution is obtained by evolving \(e\) and \(\Psi\): \(K_*=K_*(e,\Psi)\) changes with them. If a material action preserves additional algebraic dependence on \(K\), the Cartan equation retains its variational form and must be solved with that current; the preceding linear response corresponds to the specified minimal domain.

\section{Elimination of contorsion and quadratic contribution}
\label{x_reciprocidad:vii:subsec:respuesta-cuadratica}

Now consider the minimal coupling affine in \(K\), expressed by
\begin{equation}
 S_{\rm m}[e,\mathring\omega+K,\Psi]
 =S_{\rm m}[e,\mathring\omega,\Psi]
   -\tfrac12\int_M K^{IJ}\wedge\sigma_{IJ}[e,\Psi].
 \label{x_reciprocidad:vii:eq:materia-afin-contorsion}
\end{equation}
This condition identifies the material action to which the quadratic elimination applies; the current remains the one defined by \eqref{x_reciprocidad:vii:eq:corriente-variacional}.

The curvature expansion is
\begin{equation}
 F(\mathring\omega+K)
 =F(\mathring\omega)+D_{\mathring\omega}K+K\wedge K,
 \qquad (K\wedge K)^{IJ}=K^I{}_{L}\wedge K^{LJ}.
 \label{x_reciprocidad:vii:eq:curvatura-contorsion}
\end{equation}
Since \(D_{\mathring\omega}\Sigma=0\), the linear term is
the exterior derivative of \(\Sigma_{IJ}\wedge(\mathsf P_\gamma K)^{IJ}\).
The corresponding boundary contribution remains explicit:
\begin{equation}
 S_\partial[e,K]=\frac1{2\kappa c}
     \int_{\partial M}\Sigma_{IJ}\wedge(\mathsf P_\gamma K)^{IJ}.
 \label{x_reciprocidad:vii:eq:borde-contorsion}
\end{equation}
On an interior domain, or with the compatible boundary completion, the action density that depends on the contorsion is
\begin{equation}
 \mathcal Q_{e,\gamma}(K)-\tfrac12K^{IJ}\wedge\sigma_{IJ},
 \qquad
 \mathcal Q_{e,\gamma}(K)
  =\frac1{2\kappa c}\Sigma_{IJ}\wedge
                         \mathsf P_\gamma(K\wedge K)^{IJ}.
 \label{x_reciprocidad:vii:eq:forma-cuadratica-contorsion}
\end{equation}
The density \(\mathcal Q_{e,\gamma}\) is a degree-four differential form, homogeneous of degree two in \(K\).

\begin{theorem}[Effective response quadratic in the current]
\label{x_reciprocidad:vii:thm:respuesta-cuadratica}
Under \eqref{x_reciprocidad:vii:eq:materia-afin-contorsion}, substituting the algebraic solution \(K_*\) produces the effective spin density
\begin{equation}
 \mathcal L_{\rm spin}^{\rm eff}
 =-\tfrac14 K_*^{IJ}\wedge\sigma_{IJ}
 =-\mathcal Q_{e,\gamma}(K_*).
 \label{x_reciprocidad:vii:eq:accion-efectiva-spin}
\end{equation}
At fixed cotetrad and couplings, it is homogeneous of degree two in the current. The effective action also preserves \(S_\partial[e,K_*]\), or its declared boundary completion.
\end{theorem}
\begin{proof}
By Corollary~\ref{x_reciprocidad:vii:cor:desacoplamiento-torsional}, \(K_*\) is linear in \(\sigma\). The connection equation is precisely the stationarity of \eqref{x_reciprocidad:vii:eq:forma-cuadratica-contorsion}. Since it is algebraic, the radial variation \(\delta K=K_*\) is applied pointwise; it may also be multiplied by a compactly supported function. Euler's identity for a quadratic form gives
\[
 2\mathcal Q_{e,\gamma}(K_*)
             -\tfrac12K_*^{IJ}\wedge\sigma_{IJ}=0.
\]
Substitution into the action density proves \eqref{x_reciprocidad:vii:eq:accion-efectiva-spin}. If \(\sigma\mapsto\lambda\sigma\), then \(K_*\mapsto\lambda K_*\) and the effective density is multiplied by \(\lambda^2\). The expansion \eqref{x_reciprocidad:vii:eq:curvatura-contorsion} separately identifies the boundary term accompanying this elimination.
\end{proof}

The specific current, its Lorentzian contraction, the duality and the conventions of the material action determine the sign and coefficient of this contribution. The quadratic identity provides the composition rule that allows them to be evaluated once these data have been specified.

\section{Effective metric source and covariant balance}
\label{x_reciprocidad:vii:subsec:fuente-metrica-efectiva}

The normalization of the metric source is fixed in the same action line. After eliminating \(K\), we define
\begin{align}
 \delta_g S_{\rm m}[e,\mathring\omega,\Psi]
 &=-\tfrac12\int_M\operatorname{vol}_e\,
                   \mathcal T^{S,\rm LC}_{\mu\nu}\,\delta g^{\mu\nu},
 \\
 \delta_g S_{\rm spin}^{\rm eff}
 &=-\tfrac12\int_M\operatorname{vol}_e\,
                   \mathcal U^{S,\rm spin}_{\mu\nu}\,\delta g^{\mu\nu}.
 \label{x_reciprocidad:vii:eq:fuentes-normalizadas-accion}
\end{align}
Here interior variations are considered, with boundary terms treated according to the fixed variational problem. For the affine coupling, the second source is the metric variation of \eqref{x_reciprocidad:vii:eq:accion-efectiva-spin}. The variation includes the dependence of the material current on \(e\) and \(\Psi\).

\begin{proposition}[Metric equation and conservation of the total source]
\label{x_reciprocidad:vii:prop:fuente-metrica-balance}
With the preceding conventions, the metric equations are
\begin{equation}
 G_{\mu\nu}[g]+\Lambda_{\rm int}g_{\mu\nu}
 =\kappa c\bigl(\mathcal T^{S,\rm LC}_{\mu\nu}
                    +\mathcal U^{S,\rm spin}_{\mu\nu}\bigr)
 =\kappa\bigl(T^{\rm phys,LC}_{\mu\nu}
                    +U^{\rm phys,spin}_{\mu\nu}\bigr),
 \label{x_reciprocidad:vii:eq:ecuacion-metrica-efectiva}
\end{equation}
where \(T^{\rm phys,LC}=c\mathcal T^{S,\rm LC}\) and \(U^{\rm phys,spin}=c\mathcal U^{S,\rm spin}\) when the cotetrad takes values in the length line. Every solution satisfies
\begin{equation}
 \mathring\nabla^\mu
   (T^{\rm phys,LC}_{\mu\nu}+U^{\rm phys,spin}_{\mu\nu})=0.
 \label{x_reciprocidad:vii:eq:balance-metrico-total}
\end{equation}
\end{proposition}
\begin{proof}
Eliminating \(K\) preserves the equations of the remaining variables: in the chain rule, the term multiplying \(\delta K_*\) vanishes by the connection equation. In the Levi--Civita sector, the first Bianchi identity cancels the term \(\Sigma_{IJ}\wedge F^{IJ}\), and the gravitational part is written \((2\kappa c)^{-1}\int_M(R-2\Lambda_{\rm int})\operatorname{vol}_e\). Its interior variation is
\[
 \frac1{2\kappa c}\int_M\operatorname{vol}_e\,
        (G_{\mu\nu}+\Lambda_{\rm int}g_{\mu\nu})\delta g^{\mu\nu}.
\]
The sum with \eqref{x_reciprocidad:vii:eq:fuentes-normalizadas-accion} gives
\eqref{x_reciprocidad:vii:eq:ecuacion-metrica-efectiva}. The contracted Bianchi identity
for \(\mathring\nabla\), metric compatibility and
the constancy of \(\kappa\) and \(\Lambda_{\rm int}\) in the chart
produce \eqref{x_reciprocidad:vii:eq:balance-metrico-total}.
\end{proof}

Before eliminating the connection, the geometric identities retain
the form
\begin{equation}
 D_\omega T^I=F^I{}_{J}\wedge e^J,\qquad D_\omega F^{IJ}=0.
 \label{x_reciprocidad:vii:eq:bianchi-cartan}
\end{equation}
The first is \(D_\omega^2e=F\wedge e\); the second is obtained by expanding \(D_\omega(\mathrm d\omega+\omega\wedge\omega)\) and cancelling the terms. In a covariant material action, the total metric balance coincides with the Noether identity on the material equations. Geometry and matter thus preserve a joint source; the partition between its Levi--Civita contribution and its spin contribution is determined by the specialized action.

% END OWNER

\end{HMTCompactSection}

% BEGIN OWNER /Users/ruben/Documents/New project/output/EXTREMA_MEDIA_RAZON_RECIPROCIDAD_ES_EN_20260918/ARTICULO/ES/source/antecedentes_vii/32_eliminacion_torsional_no_lineal.tex
% FORMALIZACION_NUEVA / CERTIFICADO_NUEVO, 10-09-2026.
% Composición de las operaciones efectivas ya reunidas en 30e y 31.
% No modifica la selección de rutas, la frontera, los acoplamientos ni
% la normalización de la acción; no evalúa una amplitud cosmológica s_0.
% Propietarios y alcance: desarrollo/COMPOSICION_METRICA_REV06.md.
% Control independiente: pruebas/verificar_eliminacion_torsional_no_lineal.py.

\chapter{Torsional elimination and metric response of the minimal extension}
\label{x_reciprocidad:vii:sec:eliminacion-torsional-no-lineal}

The minimal extension provides a field, a response weight and a
current through a single variational operation. Its transported
incidence depends on the connection through the route products of
\eqref{x_reciprocidad:vii:eq:derivada-incidencia-rutas}. That dependence is preserved when
composing the current with the Cartan--Holst equation. Elimination
of the connection is therefore performed on the complete action of that
extension. The affine case of
\ref{x_reciprocidad:vii:thm:respuesta-cuadratica} is contained as an
exact specialization of the following construction.

\section{Variational composition at a cellular level}
\label{x_reciprocidad:vii:subsec:composicion-torsional-celular}

Let us fix a finite level, its routes, the transport representation and
the cellular pairing of
\ref{x_reciprocidad:vii:thm:componentes-minimo-conexion}. Let \(p\) be real coordinates
of the co-tetrad and the material fields in a declared local
trivialization, and \(k\in\mathbb R^d\) the coordinates of the
contortion \(K\). The family \(p\mapsto e(p)\) preserves
nondegeneracy of the co-tetrad. The connection acting on each route is
\begin{equation}
 \omega(p,k)=\mathring\omega(e(p))+K(k),\qquad
 F(p,k)=S_{\min}[e(p),\omega(p,k),\Psi(p)]
       =\mathfrak a(p,k)E(p,k).
 \label{x_reciprocidad:vii:eq:accion-minima-en-contorsion}
\end{equation}
The spaces and bases of this expression are transported to fixed
spaces before differentiation. We work on an open set where the incidence
\(\mathsf C(p,k)\) is surjective and the families are \(C^2\).
This additional regularity allows the Hessian of the action to be examined.
The couplings \(G_{\rm int},c_{\rm int},\gamma_{\rm HMT}\)
preserve the chart and the variation convention of
\ref{x_reciprocidad:vii:subsec:dominio-corriente-cartan}.

Let \(Q(p,k)\) be the integral, or the declared cellular
realization, of the quadratic form
\eqref{x_reciprocidad:vii:eq:forma-cuadratica-contorsion}. In real bases,
\begin{equation}
 Q(p,k)=\tfrac12 k^{\mathsf T}\mathsf A(p)k,
 \qquad \mathsf A(p)^{\mathsf T}=\mathsf A(p).
 \label{x_reciprocidad:vii:eq:cuadratica-celular-contorsion}
\end{equation}
The operator \(\mathsf A\) comes from the gravitational action and
its pairing, whereas the Gramian
\(\mathsf C\mathsf C^*\) comes from the extension problem.
Positivity of the latter guarantees existence of the minimizing
field; the gravitational form retains its own signature.
Boundary terms are retained according to
\eqref{x_reciprocidad:vii:eq:borde-contorsion}; the following interior formulas
apply with admissible variations that annul their contribution, or
with the corresponding boundary completion.

Define the response covector by
\begin{equation}
 d_kF(p,k)[h]= -\tfrac12 h^{\mathsf T}s(p,k),
 \qquad s(p,k)=\mathsf M_{\rm cel}(p,k)\sigma(p,k).
 \label{x_reciprocidad:vii:eq:covector-torsional-minimo}
\end{equation}
Here \(s\) and \(\sigma\) preserve the distinct types of
\eqref{x_reciprocidad:vii:eq:componentes-minimo-conexion}: the former contains the
coefficients of the differential and the latter the coordinates of the
current in the material pairing. If the basis in
\(k\) changes, that pairing is also transported. When differentiating
\(s=\mathsf M_{\rm cel}\sigma\), the derivatives of
both factors are included when the pairing depends on the parameter.

\begin{proposition}[Composed equation of extension and contorsion]
\label{x_reciprocidad:vii:prop:ecuacion-torsional-no-lineal}
Interior stationarity of the action with respect to contorsion is
\begin{equation}
 \mathcal R(p,k):=\mathsf A(p)k-\tfrac12s(p,k)=0.
 \label{x_reciprocidad:vii:eq:ecuacion-torsional-no-lineal}
\end{equation}
Each component of \(s\) is calculated from the same extension:
\begin{equation}
 \begin{split}
 s_j={}&4\mathfrak a\chi\operatorname{Re}
       \langle y,(\partial_{k_j}\mathsf C)f_{\min}
                         -\partial_{k_j}b\rangle\\
 &-2\mathfrak a(\partial_{k_j}\chi)\langle b,y\rangle
       -2(\partial_{k_j}\mathfrak a)E.
 \end{split}
 \label{x_reciprocidad:vii:eq:seleccion-variacional-corriente-k}
\end{equation}
When the cellular pairing realizes the connection equation of
\ref{x_reciprocidad:vii:thm:ecuacion-cartan-holst}, this equation is equivalent to composing
\(\sigma(p,k)\) with the Holst inverse, the inverse of
\(\mathsf L_e\) and the contortion reconstruction:
\begin{equation}
 K=\mathsf C_e^{\rm tor}\,\mathsf L_e^{-1}
           \bigl(\kappa c\,\mathsf P_\gamma^{-1}
             \sigma[e,\mathring\omega(e)+K,\Psi]\bigr).
 \label{x_reciprocidad:vii:eq:composicion-implicita-cartan}
\end{equation}
The symbol \(\mathsf C_e^{\rm tor}\) exclusively designates the
map \(T\mapsto K\) of
\eqref{x_reciprocidad:vii:eq:contorsion-componentes}, with fixed co-tetrad.
\end{proposition}

\begin{proof}
The differential of \(Q+F\), in the arbitrary direction \(h\), is
\[
 d_k(Q+F)[h]
   =h^{\mathsf T}\bigl(\mathsf A k-\tfrac12s\bigr).
\]
Its vanishing in every direction proves
\eqref{x_reciprocidad:vii:eq:ecuacion-torsional-no-lineal}.
Theorem~\ref{x_reciprocidad:vii:thm:componentes-minimo-conexion}, applied to
\(\partial_{k_j}\), gives
\eqref{x_reciprocidad:vii:eq:seleccion-variacional-corriente-k}. With \(p\) fixed,
\(\delta\omega=\delta K\). The Cartan equation is transformed
successively through \eqref{x_reciprocidad:vii:eq:inversa-holst},
\eqref{x_reciprocidad:vii:eq:torsion-explicita} and
\eqref{x_reciprocidad:vii:eq:contorsion-componentes}; each of these operations
is invertible on its domain. This proves equivalence with
\eqref{x_reciprocidad:vii:eq:composicion-implicita-cartan} when both expressions
use the same material and cellular realization.
\end{proof}

The composition preserves the connection on both sides. Routes
passing through several cells can couple the current
components between them. The pointwise reconstruction of \(T\) from
a given current and the joint solution of
\eqref{x_reciprocidad:vii:eq:ecuacion-torsional-no-lineal} are, in that case,
successive operations. This finite formulation keeps explicit
the route representation entering the functional.

\section{Local existence of a stationary branch}
\label{x_reciprocidad:vii:subsec:rama-torsional-local}

\begin{theorem}[Local continuation of the torsional solution]
\label{x_reciprocidad:vii:thm:rama-torsional-local}
Suppose \((p_0,k_0)\) belongs to the preceding open set,
\(\mathcal R(p_0,k_0)=0\), and the total Hessian
\begin{equation}
 \mathsf H_0
 =\mathsf A(p_0)+\partial_k^2F(p_0,k_0)
 =\mathsf A(p_0)-\tfrac12\partial_ks(p_0,k_0)
 \label{x_reciprocidad:vii:eq:hessiano-total-torsional}
\end{equation}
is invertible. There exist neighborhoods of \(p_0\) and \(k_0\) in which
a unique branch \(C^1\), \(k=k_*(p)\), satisfies
\eqref{x_reciprocidad:vii:eq:ecuacion-torsional-no-lineal} and passes through \(k_0\).
On that branch,
\begin{equation}
 \partial_{p_i}k_*
 =-\bigl[\mathsf A+\partial_k^2F\bigr]^{-1}
     \left[(\partial_{p_i}\mathsf A)k_*
                      +\partial_{p_i}\partial_kF\right].
 \label{x_reciprocidad:vii:eq:derivada-rama-torsional}
\end{equation}
\end{theorem}

\begin{proof}
Inversion of the Gramian is \(C^2\) on the full-rank open
set. Hence \(F\) is \(C^2\) and \(\mathcal R\) is
\(C^1\). Its derivative with respect to \(k\), at \((p_0,k_0)\), is
\(\mathsf H_0\). The implicit function theorem provides
the branch, its regularity and its uniqueness in a neighborhood of the point.
Existence can also be seen through the map
\[
 \mathcal T_p(k)=k-\mathsf H_0^{-1}\mathcal R(p,k).
\]
Its derivative with respect to \(k\) vanishes at \((p_0,k_0)\).
On a sufficiently small ball it has norm at most
\(q<1\). Shrinking the neighborhood of \(p_0\) yields
\(\|\mathcal T_p(k_0)-k_0\|<(1-q)r\), where \(r\) is the
radius of the ball. Thus, \(\mathcal T_p\) maps the closed ball
into itself and is contractive. Its fixed point is the local solution.
Differentiating \(\mathcal R(p,k_*(p))=0\) and solving for
\(\partial_{p_i}k_*\) gives
\eqref{x_reciprocidad:vii:eq:derivada-rama-torsional}.
\end{proof}

The hypothesis is verified on the composite action: the Gramian's
rank ensures the screen minimum and the total Hessian controls
the torsional branch. The asserted uniqueness is local around the
specified solution; other branches may coexist in the same
full-rank domain.

\section{Effective action and non-affine contribution}
\label{x_reciprocidad:vii:subsec:accion-efectiva-no-afin}

\begin{theorem}[Exact elimination identity]
\label{x_reciprocidad:vii:thm:eliminacion-torsional-no-lineal}
Let \(k_*(p)\) be a stationary branch. Suppose, moreover, that the
segment \(t\mapsto (p,tk_*)\), \(0\leq t\leq1\), remains
in the open set where \(F\) is defined. The action correction
relative to the Levi--Civita realization is
\begin{equation}
 \begin{split}
 \Delta S(p)
  &:=Q(p,k_*)+F(p,k_*)-F(p,0)\\
  &=\tfrac14 k_*^{\mathsf T}s(p,k_*)
     -\tfrac12\int_0^1
           k_*^{\mathsf T}s(p,tk_*)\,dt.
 \end{split}
 \label{x_reciprocidad:vii:eq:accion-efectiva-no-lineal}
\end{equation}
Equivalently,
\begin{equation}
 \Delta S=-\tfrac14 k_*^{\mathsf T}s(p,k_*)
  +\tfrac12\int_0^1 k_*^{\mathsf T}
       \bigl[s(p,k_*)-s(p,tk_*)\bigr]\,dt.
 \label{x_reciprocidad:vii:eq:correccion-no-afin-integrada}
\end{equation}
These identities refer to the integrated or cellular action.
In a realization with local density, they apply to the density
with its corresponding pairing of forms. The effective action
preserves the declared boundary term separately.
\end{theorem}

\begin{proof}
By the chain rule and
\eqref{x_reciprocidad:vii:eq:covector-torsional-minimo},
\[
 \frac{d}{dt}F(p,tk_*)
      =-\tfrac12 k_*^{\mathsf T}s(p,tk_*).
\]
Integrating between \(0\) and \(1\) gives the difference of material
actions. The stationarity equation, evaluated in the
direction \(k_*\), and the quadratic homogeneity of \(Q\) give
\[
 2Q(p,k_*)=\tfrac12 k_*^{\mathsf T}s(p,k_*).
\]
The sum of both identities proves
\eqref{x_reciprocidad:vii:eq:accion-efectiva-no-lineal}; adding and subtracting
\(\tfrac12k_*^{\mathsf T}s(p,k_*)\) provides
\eqref{x_reciprocidad:vii:eq:correccion-no-afin-integrada}.
\end{proof}

\begin{corollary}[Affine specialization]
\label{x_reciprocidad:vii:cor:recuperacion-eliminacion-afin}
If \(s(p,k)=s(p)\), then
\[
 \Delta S=-\tfrac14 k_*^{\mathsf T}s(p)=-Q(p,k_*).
\]
This is the integrated realization of
\eqref{x_reciprocidad:vii:eq:accion-efectiva-spin}.
\end{corollary}
\begin{proof}
The difference inside the integral of
\eqref{x_reciprocidad:vii:eq:correccion-no-afin-integrada} vanishes. The second
equality follows from the radial stationarity identity.
\end{proof}

\section{Metric differential of the eliminated action}
\label{x_reciprocidad:vii:subsec:diferencial-metrico-no-lineal}

\begin{theorem}[Response of the geometric variables]
\label{x_reciprocidad:vii:thm:envolvente-metrica}
On the branch of Theorem~\ref{x_reciprocidad:vii:thm:rama-torsional-local},
\begin{equation}
 \partial_{p_i}\Delta S
 =\tfrac12 k_*^{\mathsf T}(\partial_{p_i}\mathsf A)k_*
  +(\partial_{p_i}F)(p,k_*)-(\partial_{p_i}F)(p,0),
 \label{x_reciprocidad:vii:eq:envolvente-metrica-no-lineal}
\end{equation}
where the partial derivatives keep the coordinates \(k\) fixed.
In fixed Hermitian inner products, each material derivative is calculated by
\begin{equation}
 \begin{split}
 \partial_{p_i}F={}&(\partial_{p_i}\mathfrak a)E
       +\mathfrak a(\partial_{p_i}\chi)\langle b,y\rangle\\
   &+2\mathfrak a\chi\operatorname{Re}
       \langle y,\partial_{p_i}b
                    -(\partial_{p_i}\mathsf C)f_{\min}\rangle.
 \end{split}
 \label{x_reciprocidad:vii:eq:diferencial-geometrico-minimo}
\end{equation}
If \(p_i\) varies the co-tetrad, \(\partial_{p_i}\mathsf C\)
includes the variation of
\(\mathring\omega(e(p))+K(k)\), of the realized routes and of
the trivializations used, according to their declared dependencies.
\end{theorem}

\begin{proof}
The total derivative of \(Q(p,k_*)+F(p,k_*)\) contains, in
addition to its partial derivatives, the term
\[
 \bigl[\partial_kQ(p,k_*)+\partial_kF(p,k_*)\bigr]
                      \partial_{p_i}k_*.
\]
The bracket is zero by the connection equation. Subtracting the derivative
of \(F(p,0)\) proves
\eqref{x_reciprocidad:vii:eq:envolvente-metrica-no-lineal}.
The product rule and
\eqref{x_reciprocidad:vii:eq:diferencial-minimo-completo} prove
\eqref{x_reciprocidad:vii:eq:diferencial-geometrico-minimo}. The composition
\eqref{x_reciprocidad:vii:eq:accion-minima-en-contorsion} requires differentiating
the Levi--Civita connection when varying \(e\) with \(k\) fixed.
\end{proof}

When the positive extension inner product also varies, its
contribution is preserved explicitly. In coordinates, let
\(\mathsf h(p,k)>0\) be its Hermitian matrix and consider the minimum
of \(f^*\mathsf h f\) under \(\mathsf C f=b\). Then
\begin{equation}
 \mathsf L=\mathsf C\mathsf h^{-1}\mathsf C^*,\quad
 y=\mathsf L^{-1}b,\quad
 f_{\min}=\mathsf h^{-1}\mathsf C^*y,\quad
 E=\chi b^*y.
 \label{x_reciprocidad:vii:eq:minimo-producto-variable}
\end{equation}
Here \(y\) represents the constraint covector in fixed
bases. The positive matrix \(\mathsf h\) retains its energy
role, distinct from that of the Lorentzian metric.

\begin{lemma}[Variation of the extension inner product]
\label{x_reciprocidad:vii:lem:producto-variable-extension}
With the preceding conventions,
\begin{equation}
 \delta E=(\delta\chi)b^*y
   +2\chi\operatorname{Re}\bigl[y^*(\delta b-\delta\mathsf C f_{\min})\bigr]
   +\chi f_{\min}^*(\delta\mathsf h)f_{\min}.
 \label{x_reciprocidad:vii:eq:diferencial-producto-variable}
\end{equation}
\end{lemma}
\begin{proof}
The formula \(\delta\mathsf h^{-1}
=-\mathsf h^{-1}(\delta\mathsf h)\mathsf h^{-1}\) gives
\[
 \delta\mathsf L=(\delta\mathsf C)\mathsf h^{-1}\mathsf C^*
 +\mathsf C\mathsf h^{-1}(\delta\mathsf C)^*
 -\mathsf C\mathsf h^{-1}(\delta\mathsf h)
                              \mathsf h^{-1}\mathsf C^*.
\]
When differentiating \(E=\chi b^*\mathsf L^{-1}b\), the term
\(-\chi y^*(\delta\mathsf L)y\) produces the two terms
with \(\delta\mathsf C\) and the positive term with
\(\delta\mathsf h\). This proves the complete identity.
\end{proof}

In a covariant material realization whose variation admits the
metric pairing of
\eqref{x_reciprocidad:vii:eq:fuentes-normalizadas-accion}, the correction determines
\begin{equation}
 \delta_g\Delta S
 =-\tfrac12\int_M\operatorname{vol}_e\,
          \mathcal U^{S,\min}_{\mu\nu}\,\delta g^{\mu\nu},
 \qquad U^{\rm phys,\min}_{\mu\nu}
                      =c\mathcal U^{S,\min}_{\mu\nu}.
 \label{x_reciprocidad:vii:eq:fuente-metrica-minimo-no-lineal}
\end{equation}
Equations \eqref{x_reciprocidad:vii:eq:envolvente-metrica-no-lineal} and
\eqref{x_reciprocidad:vii:eq:diferencial-geometrico-minimo}, supplemented by
\eqref{x_reciprocidad:vii:eq:diferencial-producto-variable} where appropriate,
evaluate their coefficients in the declared metric variations.
The screen minimum, torsional solution and metric derivative
are three consecutive operations on the same action.
The reduced action preserves the possible couplings between
cells present in the original routes.

The homogeneous specialization of
\ref{x_reciprocidad:vii:sec:dilucion-rebote} additionally uses the realization of the
resulting tensor: its spatial symmetry, the sign of contraction
with the observer and the scaling law. When that evaluation gives
\(\rho_{\rm corr}(a)=-s_0a^{-6}\), the amplitude is read in the
reference state as
\(s_0=-a_{\rm ref}^{6}\rho_{\rm corr}(a_{\rm ref})\).
This reading preserves the material state and its normalization; the
elimination theorem provides the action to be contracted.

\section{Exact non-affine checking model}
\label{x_reciprocidad:vii:subsec:control-no-afin}

The following finite model checks the coefficients and the
non-affine dependence of the identities. Its parameters belong
to the algebraic example and keep the problem of material
realization of the routes separate.

\begin{proposition}[Exact elimination of a rank-one extension]
\label{x_reciprocidad:vii:prop:ejemplo-no-afin}
Let \(H=\mathbb R^2\), \(Q=\mathbb R\),
\(\mathsf C(k)=(1\ \ k)\), \(b=1\), \(\chi=1\),
\(\mathfrak a=\lambda>0\), and \(Q_{\rm ex}(k)=k^2/2\).
Then
\begin{equation}
 f_{\min}=\frac{(1,k)^{\mathsf T}}{1+k^2},\quad
 F(\lambda,k)=\frac{\lambda}{1+k^2},\quad
 s(\lambda,k)=\frac{4\lambda k}{(1+k^2)^2},\quad
 \mathcal R(\lambda,k)
    =k\left(1-\frac{2\lambda}{(1+k^2)^2}\right).
 \label{x_reciprocidad:vii:eq:ejemplo-minimo-no-afin}
\end{equation}
For \(\lambda=25/2\), the stationary solutions are
\(0,2,-2\). At \(k_*=2\),
\begin{equation}
 \mathsf H_*=\frac{16}{5},\quad s_*=4,\quad
 Q_{\rm ex}(k_*)=2,\quad F(\lambda,k_*)=\frac52,
 \quad \Delta S=-8.
 \label{x_reciprocidad:vii:eq:valores-ejemplo-no-afin}
\end{equation}
The affine expression \(-k_*s_*/4=-2\) differs from the exact
correction. On the positive branch,
\begin{equation}
 \partial_\lambda\Delta S=-\frac45,
 \qquad \partial_\lambda k_* =\frac1{20}
 \quad\text{in }\lambda=25/2.
 \label{x_reciprocidad:vii:eq:derivadas-ejemplo-no-afin}
\end{equation}
\end{proposition}

\begin{proof}
The Gramian is \(1+k^2>0\), so the minimum has the
indicated expressions. Differentiating \(F\) gives
\(s=-2\partial_kF\) and then \(\mathcal R=k-s/2\).
For \(\lambda=25/2\), the equation is
\(k[(1+k^2)^2-25]=0\); since \(1+k^2>0\), its roots
are precisely \(0,\pm2\). The total Hessian is
\[
 1-\frac{2\lambda(1-3k^2)}{(1+k^2)^3},
\]
which equals \(16/5\) at both nonzero roots. The correction is
\(2+5/2-25/2=-8\). In the radial identity,
\[
 \int_0^1 k_*s(\lambda,tk_*)\,dt
    =\int_0^1\frac{200t}{(1+4t^2)^2}\,dt=20,
\]
since an antiderivative is \(-25/(1+4t^2)\).
Thus \(\Delta S=8/4-20/2=-8\), including the non-affine term.
For \(\lambda>1/2\), the positive branch satisfies
\(k_*^2=\sqrt{2\lambda}-1\) and
\(\Delta S=\sqrt{2\lambda}-1/2-\lambda\).
Their derivatives give \eqref{x_reciprocidad:vii:eq:derivadas-ejemplo-no-afin}.
The envelope formula gives the same result:
\((\partial_\lambda F)(\lambda,k_*)
 -(\partial_\lambda F)(\lambda,0)=1/5-1=-4/5\).
\end{proof}

% END OWNER


% BEGIN OWNER /Users/ruben/Documents/New project/output/EXTREMA_MEDIA_RAZON_RECIPROCIDAD_ES_EN_20260918/ARTICULO/ES/source/antecedentes_vii/33_corriente_espinorial_y_densidad.tex
% Artículo VII. Incorporación aditiva autorizada, 10-09-2026.
% RESULTADO_RECUPERADO: representación real de Clifford y CAR:
% TEORIA_HOLOGRAFICA_INTEGRAL_MASA_HMT_MD_REV2_2026-08-20/fuente/modulos/
% 18_cuantica_estadistica_y_termodinamica.md, líneas 131--181 y 218--255;
% acción de primer orden: 14_topologia_exoticos_y_gravedad.md, 459--517;
% distinción corriente media/segundo momento y dilución:
% HMT_OBRAS/OBRA_II/chapters/11_cartan_bianchi_cosmologia.tex, 201--230.
% PDF REV2: b56bb38d10b9a4f4d0de47b7da8010fb31a6e8e8684a3bfd4237e37d7a3dddfa.
% FORMALIZACION_NUEVA: acción espinorial afín explícita, contracción
% normalizada mediante 31, observable cuártico CAR y funcional de amplitud.
% CERTIFICADO_NUEVO: cálculo racional de la forma cuadrática en las cuatro
% componentes de una trivector y evaluación de la realización CAR finita.
% Convención explicitada: las matrices g tienen firma (-+++); A_D=-i g^0
% y la acción cinética simétrica lleva hbar/2. La combinación equivalente
% i hbar Gamma^I usa Gamma^I=-i g^I y firma opuesta. No se mezclan ambas.
% Dependencias: cotetrada/conexión de 30d; constantes internas heredadas;
% corriente e inversas de 31. Estado material, celda y normal orden son
% datos de la realización declarada, no constantes universales extraídas.
% La acción afín de este fragmento y la acción de extensión mínima de 30e
% conservan sus problemas variacionales propios. No se identifican sus
% corrientes. No se cambia 50: su dominio de amplitud positiva se concreta.

\chapter{Spinorial current and torsional amplitude of the material state}
\label{x_reciprocidad:vii:sec:corriente-espinorial-densidad}

The preceding development and the following spinorial realization articulate
two variational problems with common geometric antecedents. The
minimal extension action \(S_{\min}=\mathfrak a_m E\), constructed in
section~\ref{x_reciprocidad:vii:sec:corriente-extension-minima}, incorporates the
dependence of the transported incidence and its minimum on the
connection; its torsional elimination preserves that complete dependence in
section~\ref{x_reciprocidad:vii:sec:eliminacion-torsional-no-lineal}. The spinorial
action \(S_D\) of \eqref{x_reciprocidad:vii:eq:accion-dirac-material} couples the
same connection to the fields of its representation and is affine in the
contortion when the co-tetrad and fields are fixed. Each current is obtained
by varying its own action. A minimizing preparation can determine
the initial material state of the spinorial realization; substitution
of connection-dependent fields additionally preserves the terms
of the composite differential. This distinction places the following amplitude
calculation in its specific action and variational domain.

The spinorial representation of transport and the algebraic response
of Cartan--Holst allow torsional intensity to be calculated as a
functional of the material state. Its evaluation preserves three data:
the field representation, the second moment of its current and
the volume relative to which density is defined. The coupling
coefficient results from the action and its inversion; the initial state
determines the material amplitude.

\section{Spinorial representation and affine material action}
\label{x_reciprocidad:vii:subsec:accion-espinorial-afin}

We preserve the co-tetrad, orientation and couplings of
section~\ref{x_reciprocidad:vii:sec:corriente-respuesta-cartan}. The spinorial
realization uses the four-dimensional complex module obtained by
complexifying the real matrices
\begin{equation}
 \begin{gathered}
 J=\begin{pmatrix}0&-1\\1&0\end{pmatrix},\qquad
 R=\begin{pmatrix}0&1\\1&0\end{pmatrix},\qquad
 Z=\begin{pmatrix}1&0\\0&-1\end{pmatrix},\\
 g^0=J\otimes I_2,\qquad g^1=R\otimes I_2,\qquad
 g^2=Z\otimes R,\qquad g^3=Z\otimes Z.
 \end{gathered}
 \label{x_reciprocidad:vii:eq:clifford-material-explicito}
\end{equation}
Multiplication gives
\(\{g^I,g^J\}=2\eta^{IJ}I_4\), with
\(\eta=\operatorname{diag}(-1,1,1,1)\). The matrices are
distinguished from the scalar parameter \(\gamma\) of the Holst
operator. Define
\begin{equation}
 g_I=\eta_{IJ}g^J,\quad
 \mathbb B_{IJ}=\tfrac14[g_I,g_J],\quad
 A_D=-i g^0,\quad \bar\psi=\psi^\dagger A_D,\quad
 \Omega=\tfrac12\omega^{IJ}\mathbb B_{IJ}.
 \label{x_reciprocidad:vii:eq:adjunto-conexion-espinorial}
\end{equation}
The Clifford product satisfies
\[
 [\mathbb B_{IJ},g^K]=\delta_J^K g_I-\delta_I^K g_J.
\]
Thus, the same metric connection acts on the fields through
\(D\psi=\mathrm d\psi+\Omega\psi\) and
\(D\bar\psi=\mathrm d\bar\psi-\bar\psi\Omega\).
Transport along a route is obtained by integrating this connection in
its spinorial representation; it preserves the connection and the route of the
preceding geometric realization.

Let \(\eta_I=\iota_{E_I}\operatorname{vol}_e\), so that
\(e^J\wedge\eta_I=\delta_I^J\operatorname{vol}_e\). For a
field of dimensional type \([\psi]=L^{-3/2}\), we take the minimal
material action
\begin{equation}
 S_D=\int_M\left\{
 \frac{\hbar}{2}
 \bigl[\bar\psi g^I D\psi-(D\bar\psi)g^I\psi\bigr]\wedge\eta_I
 -mc\,\bar\psi\psi\operatorname{vol}_e\right\}.
 \label{x_reciprocidad:vii:eq:accion-dirac-material}
\end{equation}
Here \(\hbar=\hbar_{\rm int}\), \(c=c_{\rm int}\) and the mass
\(m\) is the reading of the material species in the chosen chart.
These values are outputs preceding this realization. The matrix
\(A_D\) is Hermitian and \(A_Dg^I\) is anti-Hermitian. Therefore the
symmetric kinetic term has the real convention corresponding to
\eqref{x_reciprocidad:vii:eq:clifford-material-explicito}. In this signature, the torsion-free
equation has principal term \(\hbar g^I D_I\); the equivalent
notation \(i\hbar\Gamma^I D_I\), with \(\Gamma^I=-i g^I\),
uses the Clifford relation of opposite signature. This specification jointly
fixes the adjoint, signature and imaginary factor of the action.

The material fields and co-tetrad remain fixed in the variation
of the connection. A preparation coming from the minimal extension of
section~\ref{x_reciprocidad:vii:sec:corriente-extension-minima} can determine
the initial datum of the field or state. Each action preserves its
own variation: if a connection-dependent minimizing field
is substituted into an action, its composite differential is also applied,
or the envelope theorem under its stationarity hypotheses.
Action~\eqref{x_reciprocidad:vii:eq:accion-dirac-material} here realizes a
coupling affine in contortion.

\begin{proposition}[Spinorial current with variational normalization]
\label{x_reciprocidad:vii:prop:corriente-dirac-explicita}
Let
\(B^{IJK}=\bar\psi g^{[I}g^Jg^{K]}\psi\), where brackets
denote normalized antisymmetrization. The current defined in
\eqref{x_reciprocidad:vii:eq:corriente-variacional} by action
\eqref{x_reciprocidad:vii:eq:accion-dirac-material} is
\begin{equation}
 \sigma_{JK}
 =-\frac{\hbar}{2}\,
   \bar\psi\{g^I,\mathbb B_{JK}\}\psi\,\eta_I
 =-\frac{\hbar}{2}\,B^I{}_{JK}\eta_I.
 \label{x_reciprocidad:vii:eq:corriente-dirac-trivector}
\end{equation}
It is independent of contortion when \(e\) and \(\psi\) are
fixed. The inversion of section~\ref{x_reciprocidad:vii:subsec:reconstruccion-torsion}
therefore applies to this current.
\end{proposition}
\begin{proof}
The variation of \(\Omega\) is
\(\delta\Omega=\frac12\delta\omega^{JK}\mathbb B_{JK}\).
The two kinetic terms give, respectively, the product with
\(g^I\) on the left and on the right. Consequently,
\[
 \delta_\omega S_D
 =\frac{\hbar}{4}\int_M\delta\omega^{JK}\wedge\eta_I\,
                 \bar\psi\{g^I,\mathbb B_{JK}\}\psi.
\]
Comparison with \eqref{x_reciprocidad:vii:eq:corriente-variacional} fixes the
factor and sign of \eqref{x_reciprocidad:vii:eq:corriente-dirac-trivector}.
Expanding the anticommutator and using the Clifford relation
cancels the degree-one contractions, leaving
\(\{g^I,\mathbb B_{JK}\}=g^{[I}g_Jg_{K]}\).
The action is linear in \(\Omega\), proving the last
statement.
\end{proof}

\section{Effective contraction and second-moment coefficient}
\label{x_reciprocidad:vii:subsec:contraccion-espinorial}

The Lorentzian quadratic form of the trivector is written
\begin{equation}
 q(B)=\frac1{3!}B^{IJK}B_{IJK}
 =-(B^{012})^2-(B^{013})^2-(B^{023})^2+(B^{123})^2.
 \label{x_reciprocidad:vii:eq:cuadratica-trivector}
\end{equation}
This contraction preserves the signature of the spinorial module. Its second
moment depends on the state and on the realization of field
products.

\begin{theorem}[Effective interaction of the spinorial current]
\label{x_reciprocidad:vii:thm:coeficiente-torsional-espinorial}
With material action~\eqref{x_reciprocidad:vii:eq:accion-dirac-material}, algebraic
elimination of contortion gives
\begin{equation}
 \mathcal L_{\rm spin}^{\rm eff}
 =C_\gamma q(B)\operatorname{vol}_e,\qquad
 C_\gamma=\frac{3\kappa c\hbar^2}{16}
                      \frac{\gamma^2}{1+\gamma^2}.
 \label{x_reciprocidad:vii:eq:coeficiente-cuadratico-espinorial}
\end{equation}
The coefficient is obtained from the current and the Holst inverse
with the normalizations of the preceding action.
\end{theorem}
\begin{proof}
By \eqref{x_reciprocidad:vii:eq:accion-efectiva-spin}, it suffices to evaluate
\(-\frac14K_*^{IJ}\wedge\sigma_{IJ}\), with
\[
 Y^{IJ}=\kappa c\frac{\gamma^2}{1+\gamma^2}
           (-\star+\gamma^{-1}I)\sigma^{IJ},\qquad
 K_*=\mathsf C_e(Y).
\]
Here \(\mathsf C_e\) is the explicit composition of
\eqref{x_reciprocidad:vii:eq:torsion-explicita} with
\eqref{x_reciprocidad:vii:eq:contorsion-componentes}. Linearity of those two
inverses reduces the calculation to the parts \(-\star\) and \(I\).
To give the complete contraction, we order the trivector
components as \((012,013,023,123)\), extract the factors
\(\kappa c\hbar^2\gamma^2/(1+\gamma^2)\), and substitute
\(\sigma_{JK}=-\frac12B^I{}_{JK}\eta_I\). The two matrices
of the remaining quadratic form are
\begin{equation}
 \mathcal B_{-\star}
 =\frac3{16}\operatorname{diag}(-1,-1,-1,1),\qquad
 \mathcal B_I=0_{4\times4}.
 \label{x_reciprocidad:vii:eq:matrices-contraccion-espinorial}
\end{equation}
These matrices are obtained using only
\(e^I\wedge\eta_J=\delta_J^I\operatorname{vol}_e\): for
each component \(B^{abc}\), we form \(\sigma\), calculate
\(Y\), then
\(T^I=\sum_J\iota_{E_J}Y^{IJ}
 -\frac14e^I\wedge\sum_{L,J}\iota_{E_L}\iota_{E_J}Y^{LJ}\),
and finally
\(K_{IJK}=(T_{KIJ}-T_{IJK}-T_{JKI})/2\).
The contraction gives the four diagonal entries shown; the
six polarizations between distinct components are zero for both
operators. This calculates all coefficients of the two
quadratic forms. The part \(\gamma^{-1}I\) vanishes and the part
\(-\star\) produces exactly \eqref{x_reciprocidad:vii:eq:cuadratica-trivector},
with the coefficient of \eqref{x_reciprocidad:vii:eq:coeficiente-cuadratico-espinorial}.
\end{proof}

\section{Finite fermionic realization and second moment of the state}
\label{x_reciprocidad:vii:subsec:segundo-momento-car}

We fix a periodic comoving cell, an orthonormal frame adapted to its
observer and the homogeneous spatial mode of the four spinorial
components. This finite realization has state space
\(\mathcal F=\bigoplus_{n=0}^4\Lambda^n\mathbb C^4\).
Let \(a_r,a_r^\dagger\), \(r=1,\ldots,4\) be its canonical
operators, with
\(\{a_r,a_s^\dagger\}=\delta_{rs}\) and
\(\{a_r,a_s\}=0\). The occupation vacuum fixes normal
ordering. For a matrix \(M\), we write
\[
 \mathrm d\Gamma(M)=\sum_{r,s}a_r^\dagger M_{rs}a_s,
 \qquad \widehat N=\mathrm d\Gamma(I_4).
\]
The letter \(\Gamma\) in this notation denotes second
quantization; nonadic holonomy retains its own notation.

The four matrices of the bilinears are calculated directly:
\begin{equation}
 \begin{array}{c|cccc}
 abc&012&013&023&123\\ \hline
 M_{abc}=A_Dg^ag^bg^c
       &iJ\otimes R&iJ\otimes Z&iI_2\otimes J&iZ\otimes J
 \end{array}
 \label{x_reciprocidad:vii:eq:matrices-bilineales-espin}
\end{equation}
All are Hermitian, have zero trace and satisfy
\(M_{abc}^2=I_4\). The quartic observable corresponding to
\eqref{x_reciprocidad:vii:eq:cuadratica-trivector}, in this declared normal ordering,
is
\begin{equation}
 \begin{split}
 \widehat{\mathscr Q}_{\rm sp}
 &=\sum_{a<b<c}\eta_{aa}\eta_{bb}\eta_{cc}
       :\!\mathrm d\Gamma(M_{abc})^2\!:\\
 &=\sum_{a<b<c}\eta_{aa}\eta_{bb}\eta_{cc}
       \mathrm d\Gamma(M_{abc})^2+2\widehat N.
 \end{split}
 \label{x_reciprocidad:vii:eq:observable-cuadratico-espin}
\end{equation}
Normal ordering is part of the specified finite realization;
its definition makes explicit the subtraction of single-occupation
contractions.

\begin{proposition}[Evaluation by the canonical relations]
\label{x_reciprocidad:vii:prop:operador-espin-car}
The operator \(\widehat{\mathscr Q}_{\rm sp}\) is self-adjoint,
preserves each occupation sector and has restrictions
\begin{equation}
 \widehat{\mathscr Q}_{\rm sp}|_{\Lambda^0\oplus\Lambda^1}=0,
 \qquad
 \widehat{\mathscr Q}_{\rm sp}|_{\Lambda^3}=4I_4,
 \qquad
 \widehat{\mathscr Q}_{\rm sp}|_{\Lambda^4}=8I_1.
 \label{x_reciprocidad:vii:eq:sectores-espin-car}
\end{equation}
In the basis of \(\Lambda^2\mathbb C^4\) whose occupation
subsets are \((12,13,23,14,24,34)\), its matrix is
\begin{equation}
 \begin{pmatrix}
 0&0&0&0&0&-4\\
 0&2&0&0&6&0\\
 0&0&2&-6&0&0\\
 0&0&-6&2&0&0\\
 0&6&0&0&2&0\\
 -4&0&0&0&0&0
 \end{pmatrix}.
 \label{x_reciprocidad:vii:eq:sector-dos-espin-car}
\end{equation}
\end{proposition}
\begin{proof}
The relation \(a_s a_t^\dagger=\delta_{st}-a_t^\dagger a_s\)
applied to \(\mathrm d\Gamma(M)^2\) separates its
single-occupation contraction and gives
\[
 :\!\mathrm d\Gamma(M)^2\!:
 =\mathrm d\Gamma(M)^2-\mathrm d\Gamma(M^2).
\]
Since the first three signs of
\eqref{x_reciprocidad:vii:eq:matrices-bilineales-espin} are negative and the last
positive, the sum of contractions is \(-2\widehat N\).
This proves \eqref{x_reciprocidad:vii:eq:observable-cuadratico-espin}.
Self-adjointness and preservation of occupation follow from
those same expressions.

On \(\Lambda^0\) all terms vanish. On \(\Lambda^1\),
\(\mathrm d\Gamma(M)=M\), and the squares sum to \(-2I_4\),
which cancels with \(2\widehat N\). On \(\Lambda^4\),
\(\mathrm d\Gamma(M)=\operatorname{tr}M=0\), leaving eight.
Identification by complement between \(\Lambda^3\mathbb C^4\)
and the dual of \(\mathbb C^4\) takes
\(\mathrm d\Gamma(M)\) to
\((\operatorname{tr}M)I-M^{\mathsf T}=-M^{\mathsf T}\).
Their squares again sum to \(-2I_4\), while
\(2\widehat N=6I_4\); therefore the restriction is \(4I_4\).
Finally, applying the four explicit matrices to the
six exterior products of the indicated basis produces
\eqref{x_reciprocidad:vii:eq:sector-dos-espin-car}. The operator is thus evaluated
in the sixteen dimensions of \(\mathcal F\).
\end{proof}

For a state operator \(\varrho\geq0\),
\(\operatorname{Tr}\varrho=1\), define
\(q_\varrho=\operatorname{Tr}(\varrho\widehat{\mathscr Q}_{\rm sp})\).
Its dependence on second moments is explicit:
\begin{equation}
 q_\varrho=\sum_{a<b<c}\eta_{aa}\eta_{bb}\eta_{cc}
 \left\{
  \langle\widehat B_{abc}\rangle_\varrho^2
  +\operatorname{Var}_\varrho(\widehat B_{abc})
  -\langle\widehat N\rangle_\varrho
 \right\},\qquad
 \widehat B_{abc}=\mathrm d\Gamma(M_{abc}).
 \label{x_reciprocidad:vii:eq:covarianza-espin-material}
\end{equation}
Covariance and normal subtraction are preserved even when
the mean currents vanish. In particular, if \(P_3\) is the
projection onto \(\Lambda^3\mathbb C^4\),
\begin{equation}
 \varrho_3=\tfrac14P_3,\qquad
 \langle\widehat B_{abc}\rangle_{\varrho_3}=0,\qquad
 q_{\varrho_3}=4>0.
 \label{x_reciprocidad:vii:eq:testigo-positivo-espin}
\end{equation}
It is a positive state of trace one and constitutes an explicit witness
to the positive domain of the functional. The two-occupation sector
also contains negative values; for example, the vector
\((|12\rangle+|34\rangle)/\sqrt2\) has value \(-4\).
The sign of the intensity thus comes from the contraction and the
state, and remains separate from the selection of cosmological conditions.

\section{Density, pressure and preservation of the comoving amplitude}
\label{x_reciprocidad:vii:subsec:amplitud-espinorial-comovil}

In the spatially homogeneous realization, fix
\[
 e^0=N(t)c\,\mathrm dt,\qquad e^j=a(t)\,\mathrm dx^j,
 \qquad \mathcal V(t)=a(t)^3V_c,
\]
with lapse \(N(t)>0\), dimensionless scale factor \(a(t)>0\)
and comoving volume \(V_c>0\). The frame and periodic mode of
subsection~\ref{x_reciprocidad:vii:subsec:segundo-momento-car} remain declared.
The normalization of the homogeneous field is
\(\widehat\psi=\mathcal V^{-1/2}(a_1,a_2,a_3,a_4)^{\mathsf T}\).
Each bilinear carries a factor \(\mathcal V^{-1}\), and its normally
ordered product carries \(\mathcal V^{-2}\). The symmetric kinetic
temporal form preserves this canonical normalization: terms
arising from differentiating \(\mathcal V^{-1/2}\) cancel between
its two summands.

\begin{theorem}[Torsional amplitude as a functional of the state]
\label{x_reciprocidad:vii:thm:amplitud-material-explicita}
Keeping the canonical state and the realization of the product fixed
during metric variations, the effective spin term has
reduced action
\begin{equation}
 S_{\rm spin}^{\rm hom}
 =\int\!\mathrm dt\,N(t)
       \frac{cC_\gamma}{a(t)^3V_c}\,q_\varrho.
 \label{x_reciprocidad:vii:eq:accion-homogenea-espin}
\end{equation}
Its energy density and isotropic pressure are
\begin{equation}
 \rho_{\rm spin}=p_{\rm spin}
 =-\frac{cC_\gamma}{V_c^2}\,q_\varrho a^{-6}.
 \label{x_reciprocidad:vii:eq:densidad-presion-espin}
\end{equation}
In any evolution preserving \(q_\varrho\), the corresponding
constant amplitude is
\begin{equation}
 \boxed{\displaystyle
 s_0[\varrho,V_c]
 =\frac{3\kappa c^2\hbar^2}{16}
    \frac{\gamma^2}{1+\gamma^2}
    \frac{\operatorname{Tr}
       (\varrho\widehat{\mathscr Q}_{\rm sp})}{V_c^2}.}
 \label{x_reciprocidad:vii:eq:amplitud-s0-espinorial}
\end{equation}
\end{theorem}
\begin{proof}
Spatial integration of
\eqref{x_reciprocidad:vii:eq:coeficiente-cuadratico-espinorial} multiplies the
factor \(\mathcal V^{-2}\) by
\(Nc\mathcal V\,\mathrm dt\), giving
\eqref{x_reciprocidad:vii:eq:accion-homogenea-espin}. In canonical variables,
lapse variation defines energy through
\(\delta_N S_{\rm m}=-\int E\,\delta N\,\mathrm dt\).
Therefore,
\[
 E_{\rm spin}=-\frac{cC_\gamma q_\varrho}{a^3V_c},
 \qquad \rho_{\rm spin}=\frac{E_{\rm spin}}{a^3V_c}.
\]
The isotropic spatial variation satisfies
\(\delta_a S_{\rm m}=\int3Na^2V_c p\,\delta a\,\mathrm dt\).
Applied to \eqref{x_reciprocidad:vii:eq:accion-homogenea-espin}, it gives
\[
 3Na^2V_c p_{\rm spin}
 =-\frac{3NcC_\gamma q_\varrho}{a^4V_c}.
\]
This simultaneously proves the sign, pressure and density of
\eqref{x_reciprocidad:vii:eq:densidad-presion-espin}. Since
\(q_\varrho\) is preserved, the coefficient of \(-a^{-6}\) is constant and
we obtain \eqref{x_reciprocidad:vii:eq:amplitud-s0-espinorial}. Dimensionality
is consistent: \([cC_\gamma]=\text{energy}\,L^3\),
so \(s_0\) has the dimension of energy density.
\end{proof}

\begin{corollary}[Preserved domain of positive amplitude]
\label{x_reciprocidad:vii:cor:dominio-positivo-conservado}
Let there be a unitary evolution of \(\mathcal F\) preserving
\(\Lambda^3\mathbb C^4\). For every initial state of trace one
supported on that sector, \(q_{\varrho(t)}=4\) holds, and
\begin{equation}
 s_0=\frac{3\kappa c^2\hbar^2}{4V_c^2}
                   \frac{\gamma^2}{1+\gamma^2}>0.
 \label{x_reciprocidad:vii:eq:amplitud-positiva-sector-tres}
\end{equation}
\end{corollary}
\begin{proof}
Invariance of the sector preserves the support of the state; on it,
\eqref{x_reciprocidad:vii:eq:sectores-espin-car} gives
\(\widehat{\mathscr Q}_{\rm sp}=4I\). Preservation of the
trace then fixes its expectation at four, independently of
the internal evolution of its coherences. Positivity of the
couplings, \(V_c>0\) and \(\gamma\ne0\) proves the formula
and its sign.
\end{proof}

For general states, preservation of
\(\langle\widehat N\rangle\) determines dilution of the occupation
density, \(n(t)=\langle\widehat N\rangle/(a^3V_c)\).
Preservation of the quartic moment is checked separately.
If \(i\hbar\dot\varrho=[H(t),\varrho]\), then
\begin{equation}
 \frac{\mathrm d q_\varrho}{\mathrm dt}
 =\frac{i}{\hbar}\operatorname{Tr}
       \bigl(\varrho[H(t),\widehat{\mathscr Q}_{\rm sp}]\bigr).
 \label{x_reciprocidad:vii:eq:conservacion-momento-cuartico}
\end{equation}
Commutation, invariance of a sector where the observable is
scalar, or a conservation proved for the specific state
provide the relevant conditions. If the moment evolves,
\eqref{x_reciprocidad:vii:eq:densidad-presion-espin} preserves its instantaneous
evaluation; its balance includes
\[
 \dot\rho_{\rm spin}+6H\rho_{\rm spin}
 =-\frac{cC_\gamma}{V_c^2}\,a^{-6}\dot q_\varrho,
 \qquad H=\dot a/a
\]
in proper time. This term records exchange with the
remaining material degrees of freedom.

The composition is determined: spinorial representation of
transport, material action, variational current, Cartan--Holst
inversion, second moment and normalized density. The domain
\(s_0>0\) of section~\ref{x_reciprocidad:vii:sec:dilucion-rebote} thus has
an evaluable functional and explicit witness states. Its use in
cosmological dynamics additionally preserves the material,
geometric and conservation conditions stated in that section.

% END OWNER


% BEGIN OWNER /Users/ruben/Documents/New project/output/EXTREMA_MEDIA_RAZON_RECIPROCIDAD_ES_EN_20260918/ARTICULO/ES/source/antecedentes_vii/50_dilucion_y_rebote.tex
% Artículo VII: pruebas preexistentes incorporadas al cuerpo.
% Propietario íntegro: fuentes_conservadas/17_rebote_einstein_cartan.tex.
% Incluye la ley de escala y las pruebas de rebote y persistencia.
% La amplitud s_0 es la salida de la realización material antecedente;
% este fragmento no certifica por sí solo la evaluación de esa realización.

\chapter{Memory dilution and persistence of the bounce}
\label{x_reciprocidad:vii:sec:dilucion-rebote}

The evolution of memory and spatial expansion are two readings of
the same sequence of states. The refinement index can be eliminated
between them; that elimination determines the dilution exponent transmitted
by the cosmological realization to its torsional contribution.
The initial amplitude depends on evaluation of the concrete material
current in its variational response. The scaling law and bounce
theorem that follow preserve that dependence.

Theorem~\ref{x_reciprocidad:vii:thm:amplitud-material-explicita} determines that
amplitude as a functional of the material state and comoving volume.
Corollary~\ref{x_reciprocidad:vii:cor:dominio-positivo-conservado} provides
an explicit domain of states with positive, conserved amplitude.
In that domain, evaluation of \(s_0\) precedes the cosmological
integration developed below. The ordinary density
\(\rho_0>0\) and its barotropic law are additionally specified in
section~\ref{x_reciprocidad:vii:sec:umbral-rebote}: the witness to positivity of the
amplitude does not in itself evaluate that material density.

\section{The \texorpdfstring{\(a^{-6}\)}{a to the power -6} dilution law}
\label{x_reciprocidad:vii:sec:dilucion}

\begin{proposition}[Elimination of the memory index]
\label{x_reciprocidad:vii:prop:dilucion}
Let \(a_{\mathrm{ref}}>0\), \(\varphi>1\), and let the route readers be
\[
 \frac{a_n}{a_{\mathrm{ref}}}=\varphi^{n/3},\qquad
 M_n=\varphi^{-2n}.
\]
Then
\[
 M_n=\left(\frac{a_n}{a_{\mathrm{ref}}}\right)^{-6}
\]
for every index of the sequence. If the realization of the quadratic
density is \(s_n^2=s_*^2M_n\), with \(s_*^2>0\), we obtain
\[
 s_n^2=s_*^2\left(\frac{a_n}{a_{\mathrm{ref}}}\right)^{-6}.
\]
\end{proposition}

\begin{proof}
Raising the first equality to the power \(-6\) gives
\(\varphi^{-2n}\), which is exactly the second reader. Multiplying
by the initial amplitude reproduces the second statement. The exponent
is determined by the quotient of the exponents of the two
readers; the amplitude is the evaluation of the reference state.
\end{proof}

In the continuous homogeneous realization, this constitutive law is preserved:
the torsional density is written \(s_0a^{-6}\), where the normalization
of the factor \(a\), the material coefficient and the reference amplitude
have been absorbed into \(s_0\). The scaling law and evaluation of
\(s_0\) are successive operations of the construction.

\section{Homogeneous threshold and transversality}
\label{x_reciprocidad:vii:sec:umbral-rebote}

Consider the homogeneous, spatially flat chart, in units \(c=1\),
with constant \(G>0\), positive densities \(\rho_0,s_0\) and constant
barotropic parameter \(w<1\). The specified realization satisfies
\(\Lambda_{\mathrm{eff}}=0\) and the equations
\begin{align}
 \rho_{\mathrm m}(a)&=\rho_0a^{-3(1+w)},&
 \rho_{\mathrm s}(a)&=s_0a^{-6},\\
 H^2&=\frac{8\pi G}{3}
        \bigl(\rho_{\mathrm m}-\rho_{\mathrm s}\bigr)
        =:F_0(a),&
 \dot H&=-4\pi G
        \bigl((1+w)\rho_{\mathrm m}-2\rho_{\mathrm s}\bigr).
 \label{x_reciprocidad:vii:eq:dinamica-homogenea}
\end{align}
The couplings belong to the dimensional section already constructed.
The expression for \(\rho_{\mathrm s}\) incorporates the sign of its
contribution in this gravitational realization.

\begin{theorem}[Existence and uniqueness of the transverse threshold]
\label{x_reciprocidad:vii:thm:rebote}
There exists a unique \(a_{\mathrm b}>0\) with \(F_0(a_{\mathrm b})=0\):
\[
 a_{\mathrm b}^{\,3(1-w)}=\frac{s_0}{\rho_0}.
\]
The region allowed by \(H^2\geq0\) is \(a\geq a_{\mathrm b}\).
At the threshold, with
\(\rho_{\mathrm b}=\rho_{\mathrm m}(a_{\mathrm b})
                  =\rho_{\mathrm s}(a_{\mathrm b})>0\),
\[
 \dot H_{\mathrm b}=4\pi G(1-w)\rho_{\mathrm b}>0.
\]
The solution crossing \(H=0\) has a strict minimum of \(a\),
with local expansion
\[
 a(t)=a_{\mathrm b}\left[
 1+2\pi G(1-w)\rho_{\mathrm b}(t-t_{\mathrm b})^2
 +O((t-t_{\mathrm b})^3)\right].
\]
\end{theorem}

\begin{proof}
The factorization
\[
 F_0(a)=\frac{8\pi G}{3}a^{-6}
             \bigl(\rho_0a^{3(1-w)}-s_0\bigr)
\]
reduces its zeros to that of a strictly increasing function of
\((0,\infty)\) on \((-s_0,\infty)\). This proves existence,
uniqueness and the sign of \(F_0\) on both sides of the threshold.
The second equation of \eqref{x_reciprocidad:vii:eq:dinamica-homogenea}, evaluated
at equality of densities, gives the formula for \(\dot H_{\mathrm b}\).
Since \(\dot a=aH\), at the threshold
\(\ddot a_{\mathrm b}=a_{\mathrm b}\dot H_{\mathrm b}>0\).
The equations are smooth for \(a>0\), so their Taylor
expansion gives the indicated expansion. The sign change of \(H\) is
transverse: it passes from contraction to expansion.
\end{proof}

For \(w=1\), both terms have the same exponent and equality of
amplitudes determines a degenerate case. For \(w>1\), the sign of
\(\dot H\) at a coincidence of densities is negative.
The domain \(w<1\) of the theorem expresses exactly the order between the
two dilution exponents.

\section{Local persistence under controlled perturbations}
\label{x_reciprocidad:vii:sec:persistencia}

Persistence of the threshold uses two independent estimates:
control of the sign at the endpoints and of the derivative on the
interval. Let \(I=[a_-,a_+]\subset(0,\infty)\), with
\(a_-<a_{\mathrm b}<a_+\), be chosen so that
\[
 F_0(a_-)<0<F_0(a_+),\qquad
 m_I:=\inf_{a\in I}F_0'(a)>0.
\]
Such an interval exists because
\[
 F_0'(a_{\mathrm b})
   =\frac{8\pi G}{a_{\mathrm b}}(1-w)\rho_{\mathrm b}>0.
\]
Define
\(\delta_I=\min\{-F_0(a_-),F_0(a_+)\}>0\).

\begin{theorem}[Local persistence of the transverse bounce]
\label{x_reciprocidad:vii:thm:persistencia}
Let \(R\in C^1(I)\) be a perturbation satisfying
\[
 \|R\|_{\infty,I}<\delta_I,\qquad
 \|R'\|_{\infty,I}<m_I.
\]
Then \(F_R=F_0+R\) has a unique zero \(a_R\in(a_-,a_+)\), and
\(F_R'(a_R)>0\). The realization \(H^2=F_R(a)\) admits a local
transverse bounce, with
\[
 \dot H_R=\frac{a_R}{2}F_R'(a_R)>0
\]
at the threshold. If \(R(a,\eta)\) is jointly \(C^1\), the persistent
zeros depend locally in a \(C^1\) manner on the parameter \(\eta\).
\end{theorem}

\begin{proof}
The first bound preserves
\(F_R(a_-)<0<F_R(a_+)\); the intermediate value theorem provides
an interior zero. The second gives
\(F_R'(a)\geq m_I-\|R'\|_{\infty,I}>0\) on \(I\), so the zero
is unique and simple.

At \(a>a_R\) near the zero,
\(F_R(a)=F_R'(a_R)(a-a_R)+o(a-a_R)\).
The integral
\[
 t-t_{\mathrm b}=\int_{a_R}^{a}
          \frac{d\xi}{\xi\sqrt{F_R(\xi)}}
\]
is finite and strictly increasing. Its inverse defines the expanding
branch; time reflection defines the contracting branch.
The two join with \(\dot a(t_{\mathrm b})=0\) and
\[
 a(t)-a_R
 =\frac{a_R^2F_R'(a_R)}4(t-t_{\mathrm b})^2
   +o((t-t_{\mathrm b})^2).
\]
Away from the threshold, differentiating \(H^2=F_R(a)\) gives
\(\dot H=aF_R'(a)/2\); its continuous limit gives the positive expression
at \(a_R\). The \(C^1\) regularity of \(F_R\) guarantees continuity
of acceleration at this junction. Finally,
\(\partial_aF_R(a_R)>0\) allows application of the implicit function
theorem when the perturbation depends jointly on \((a,\eta)\).
\end{proof}

The first bound controls existence; the second ensures local
uniqueness and transversality. The proved stability corresponds to
that threshold in the selected interval. The torsional amplitude
enters its location, while the relation between
exponents determines the temporal orientation of the crossing.

% END OWNER


% BEGIN OWNER /Users/ruben/Documents/New project/output/EXTREMA_MEDIA_RAZON_RECIPROCIDAD_ES_EN_20260918/ARTICULO/ES/source/antecedentes_vii/51_solucion_homogenea.tex
% FORMALIZACION_NUEVA de la realización homogénea preexistente.
% Prioridad de la fórmula integrada: procedencia histórica pendiente de cotejo.
% Conserva la ecuación y el dominio del capítulo de rebote; no selecciona s_0.
\chapter{Explicit homogeneous solution and causal completeness}
\label{x_reciprocidad:vii:sec:solucion-homogenea}
Determination of the amplitude allows a complete realization
of the bounce to be developed, in addition to the local argument. The spatially
flat chart of \eqref{x_reciprocidad:vii:eq:dinamica-homogenea} is retained, in units
\(c=1\), with \(\rho_0,s_0,G>0\), and the material sector is specialized
to dust, \(w=0\), with \(\Lambda_{\mathrm{eff}}=0\).
The coefficient \(s_0\) is that obtained from the material
state in the corresponding realization; the following integration
uses that coefficient and preserves its domain.

\begin{theorem}[Exact integration of the dust bounce]
\label{x_reciprocidad:vii:thm:rebote-polvo}
Set
\[
 a_{\rm b}^3=\frac{s_0}{\rho_0},\qquad
 \tau^2=\frac{s_0}{6\pi G\rho_0^2},\qquad u=t-t_{\rm b}.
\]
The solution of the homogeneous system with
\(a(t_{\rm b})=a_{\rm b}\) and \(H(t_{\rm b})=0\) is, for every
\(t\in\mathbb R\),
\begin{equation}
 a(t)=a_{\rm b}\left(1+\frac{u^2}{\tau^2}\right)^{1/3},
 \qquad H(t)=\frac{2u}{3(\tau^2+u^2)},\qquad
 \dot H(t)=\frac{2(\tau^2-u^2)}{3(\tau^2+u^2)^2}.
 \label{x_reciprocidad:vii:eq:solucion-polvo}
\end{equation}
It has a unique minimum, at \(t_{\rm b}\); it contracts before and
expands afterward. The evolution is smooth and satisfies
\(a(t)\geq a_{\rm b}>0\).
\end{theorem}
\begin{proof}
Let \(y=a^3\). The Friedmann equation gives
\[
 \dot y^{\,2}=24\pi G(\rho_0y-s_0).
\]
The function \(y=s_0/\rho_0+6\pi G\rho_0u^2\) satisfies that
identity and the initial conditions. Its derivatives give the
three expressions of \eqref{x_reciprocidad:vii:eq:solucion-polvo}. To check
the evolution equation as well, let us write
\[
 \rho_{\rm b}=\frac{\rho_0^2}{s_0},\qquad
 x=\frac{u}{\tau}.
\]
Then
\[
 \rho_{\rm m}=\frac{\rho_{\rm b}}{1+x^2},\qquad
 \rho_{\rm s}=\frac{\rho_{\rm b}}{(1+x^2)^2},\qquad
 4\pi G\rho_{\rm b}=\frac{2}{3\tau^2}.
\]
Substituting into \(\dot H=-4\pi G(\rho_{\rm m}-2\rho_{\rm s})\)
gives exactly the third formula. Thus the complete
system \((\dot a,\dot H)\) is satisfied, whose field is smooth
for \(a>0\). Its local uniqueness, continued along the explicit
solution, fixes the evolution through those initial conditions.
The denominator is positive for every \(t\), and the statements
about the minimum and the sign follow from \(H\).
\end{proof}

\begin{figure}[htbp]
\centering
\begin{tikzpicture}[x=1.45cm,y=1.15cm,>=Latex]
\draw[->] (-3.25,0)--(3.3,0) node[right] {\(\displaystyle u/\tau\)};
\draw[->] (0,0)--(0,2.8) node[above] {\(\displaystyle a/a_{\rm b}\)};
\draw[densely dashed,gray] (-3.1,1)--(3.1,1);
\draw[very thick,ink,domain=-3:3,samples=121,smooth]
 plot (\x,{(1+\x*\x)^(1/3)});
\fill[accent] (0,1) circle (1.6pt);
\node[anchor=north east] at (-.08,.98) {\(1\)};
\node[ink] at (-2,2.65) {Contraction};
\node[ink] at (2,2.65) {Expansion};
\foreach \t in {-3,-2,-1,1,2,3}
{\draw (\t,.045)--(\t,-.045) node[below] {\(\t\)};}
\end{tikzpicture}
\caption{Dimensionless profile of the proved solution:
\(a/a_{\rm b}=(1+(u/\tau)^2)^{1/3}\).
The parameters \(a_{\rm b}\) and \(\tau\) are calculated from
\(\rho_0,s_0,G\); the drawing represents the entire rescaled family.}
\label{x_reciprocidad:vii:fig:rebote-polvo}
\end{figure}

\section{Finite curvature and extension of geodesics}
The metric realized in this chart is
\(g=-dt^2+a(t)^2\,d\mathbf x^2\) on
\(\mathbb R\times\mathbb R^3\). With the convention in which
the flat curvature scalar is \(R=6(\dot H+2H^2)\), the solution gives
\[
 R(t)=\frac{4\tau^2+\frac43u^2}{(\tau^2+u^2)^2}.
\]
In particular, \(R(t_{\rm b})=4/\tau^2\) is finite.
In an orthonormal frame, the Riemann components depend on
\(H^2\) and on \(\dot H+H^2\), which are bounded functions on
\(\mathbb R\). The same holds for polynomial contractions
of those components; for example,
\[
 R_{\mu\nu\rho\sigma}R^{\mu\nu\rho\sigma}
 =12\bigl[(\dot H+H^2)^2+H^4\bigr].
\]

\begin{proposition}[Completeness of causal geodesics]
\label{x_reciprocidad:vii:prop:completitud-causal-polvo}
The preceding metric realization is geodesically complete
for the timelike and null geodesics of its Levi--Civita
connection, in both directions.
\end{proposition}
\begin{proof}
Spatial translations provide the constants
\(p_i=a^2\,dx^i/d\lambda\). Causal normalization gives
\[
 \left|\frac{dt}{d\lambda}\right|
 =\sqrt{\epsilon+\frac{|\mathbf p|^2}{a(t)^2}},
 \qquad \epsilon=1\ \text{in the timelike case},\quad
 \epsilon=0\ \text{in the null case}.
\]
A nonconstant null geodesic has \(|\mathbf p|>0\).
Since \(a\geq a_{\rm b}\), in the timelike case
\[
 \left|\frac{d\lambda}{dt}\right|
 \geq\left(1+\frac{|\mathbf p|^2}{a_{\rm b}^2}\right)^{-1/2}>0,
\]
and in the null case
\(\left|d\lambda/dt\right|=a/|\mathbf p|
 \geq a_{\rm b}/|\mathbf p|>0\).
Thus, reaching \(t=\pm\infty\) requires infinite affine length.
On a compact time interval, \(a\) and its derivatives are
smooth and \(a\) is bounded away from zero. The preceding explicit
equations and \(dx^i/d\lambda=p_i/a^2\) allow the
geodesic to be continued through its finite endpoints. Both
possibilities of continuation are covered.
\end{proof}

The completeness result belongs to this homogeneous dust
realization, with its spatial, matter and amplitude conditions. Its
local persistence under perturbations of the threshold preserves the bounds
of theorem~\ref{x_reciprocidad:vii:thm:persistencia}; that local property and the
global completeness of the explicit solution are distinct conclusions.
The minimum scale comes from the relation between material density and
torsional response. Memory thus enters a verifiable geometric
evolution, with a regular realization of contraction
and expansion.

% END OWNER


% BEGIN OWNER /Users/ruben/Documents/New project/output/EXTREMA_MEDIA_RAZON_RECIPROCIDAD_ES_EN_20260918/ARTICULO/ES/source/antecedentes_vii/60_tensor_residual_y_balance.tex
% Artículo VII. Recuperación por dependencia, 10-09-2026.
% RESULTADO_RECUPERADO: integral activo 2249, fuente/colaboracion/
% partes_iv_v/tex/residencias/86_cosmologia_cerrada.tex (S04),
% líneas 2--44 (cartas seleccionadas), 174--248 (tensor y aceleración).
% Fuente S04 leída íntegramente, SHA-256:
% feb2582c3316c8b30829cfb0476bfa2f4e95889dabfba16fe43f2d835a226673.
% Se recuperan definición, unicidad, conservación, descomposición,
% intercambio y aceleración, con sus condiciones explícitas.
% FORMALIZACION_REUNIDA: relación con la fuente efectiva de 31;
% se distingue conservación total de conservación material separada.
% CERTIFICADO_NUEVO: expansión de las pruebas de traza, balance y signos;
% cálculo local de Ricci para las tres cartas ya seleccionadas por S04.
% No se modifica el rebote de 50 ni se determina aquí una amplitud s_0.
%
% DEPENDENCIAS DE INTEGRACIÓN:
% - vii:sec:realizacion-cartan: la métrica procede de la cotetrada realizada.
% - vii:sec:acoplamientos: G_int y el reloj t_0 del mismo estado HMT;
%   c_int=1 es una elección de unidades posterior a su construcción.
% - vii:eq:ecuacion-metrica-efectiva y vii:eq:balance-metrico-total, de 31:
%   se usan únicamente para identificar el residual sobre una solución.
% - T_b es el tensor publicado por el transductor material declarado;
%   su conservación separada es una condición expresamente enunciada.
% - El selector TRIT de la carta y t_0(x)>0 son antecedentes de la
%   especialización de Sitter, no datos cosmológicos observacionales.
% El presente fragmento no certifica la clausura del conjunto de VII.

\chapter{Residual tensor, trace decomposition, and covariant balance}
\label{x_reciprocidad:vii:sec:tensor-residual-balance}

The geometric realization of the enriched state provides a metric,
while the material transducer publishes the source to which the
reading refers. The residual tensor preserves the difference between the two publications
in the same tensor space. Its decomposition distinguishes the
sector proportional to the metric from the traceless sector, and determines
the exchange current corresponding to that partition.

\section{Chart data and construction of the residual tensor}
\label{x_reciprocidad:vii:subsec:datos-tensor-residual}

Let us fix an HMT state \(x\) and a connected chart \(\mathcal U\) of its
geometric realization. On \(\mathcal U\), the nondegenerate co-tetrad
of section~\ref{x_reciprocidad:vii:sec:realizacion-cartan} determines a smooth metric
\(g=g_x\) of signature \((-+++)\). We denote by
\(\mathring\nabla\) its Levi--Civita connection and write
\[
 \mathsf G_{\mu\nu}[g]
 :=\operatorname{Ric}_{\mu\nu}[g]-\tfrac12R[g]g_{\mu\nu}.
\]
This distinguishes the Einstein tensor \(\mathsf G[g]\) from
the gravitational coupling \(G_{\rm int}\). The derivative used
in the balances of this section is \(\mathring\nabla\),
consistently with the effective metric equation of
Section~\ref{x_reciprocidad:vii:subsec:fuente-metrica-efectiva}.

We work in the dimensional chart \(c_{\rm int}=1\). We assume
\(G_{\rm int}(x)>0\) and \(\Lambda_{\rm ref}(x)\) constant on
\(\mathcal U\), and put \(\kappa_x=8\pi G_{\rm int}(x)\).
The symbol \(x\) identifies the state supplying those readers;
covariant differentiation acts on the points of \(\mathcal U\).
The genealogy of the couplings remains in
section~\ref{x_reciprocidad:v:ant:sec:gravity}.

Let \(T_b\) be the smooth symmetric tensor published by the chart's material
transducer; in the baryonic reading, it is its baryonic source. We define
\begin{equation}
 \mathcal E^{\rm ref}_{\mu\nu}
 :=\kappa_x^{-1}\bigl(\mathsf G_{\mu\nu}[g]
                            +\Lambda_{\rm ref}g_{\mu\nu}\bigr),
 \qquad
 T^{\rm res}_{\mu\nu}:=\mathcal E^{\rm ref}_{\mu\nu}-T_{b,\mu\nu}.
 \label{x_reciprocidad:vii:eq:residual-definicion}
\end{equation}
The tensor \(T^{\rm res}\) is the effective residual tensor, also
denoted \(T_{\rm osc}\) when interpreting the geometric sector remaining
to be published relative to \(T_b\). The construction preserves
\(g_x,G_{\rm int},\Lambda_{\rm ref}\) and \(T_b\) as effective arguments.

\begin{theorem}[Uniqueness of the residual and balance with the material source]
\label{x_reciprocidad:vii:thm:residual-unicidad-balance}
For the preceding data there exists a unique symmetric tensor
\(T^{\rm res}\) satisfying
\(\mathcal E^{\rm ref}=T_b+T^{\rm res}\). Its divergence is
\begin{equation}
 \mathring\nabla^\mu T^{\rm res}_{\mu\nu}
 =-\mathring\nabla^\mu T_{b,\mu\nu}.
 \label{x_reciprocidad:vii:eq:balance-residual-material}
\end{equation}
In particular, if the material source is conserved separately,
\(\mathring\nabla^\mu T_{b,\mu\nu}=0\), then
\(T^{\rm res}\) is also conserved.
\end{theorem}
\begin{proof}
The difference of two tensors satisfying the stated equation is
zero; expression~\eqref{x_reciprocidad:vii:eq:residual-definicion} provides its
existence. The contracted Bianchi identity and Levi--Civita metric
compatibility give
\[
 \mathring\nabla^\mu\mathsf G_{\mu\nu}=0,
 \qquad \mathring\nabla g=0.
\]
Since \(\kappa_x\) and \(\Lambda_{\rm ref}\) are constant on the
chart, \(\mathring\nabla^\mu\mathcal E^{\rm ref}_{\mu\nu}=0\).
Differentiating the difference defining \(T^{\rm res}\) proves
\eqref{x_reciprocidad:vii:eq:balance-residual-material}, including the statement under
separate material conservation.
\end{proof}

The spatial and temporal constancy of the couplings in the chart fixes
the domain of this balance. A family of states may publish
different couplings in different charts; each preceding identity
applies to the chart and state specifying it.

\section{Canonical decomposition and exchange current}
\label{x_reciprocidad:vii:subsec:traza-corriente-residual}

Define the trace function and the two tensors
\begin{equation}
 \theta:=g^{\mu\nu}T^{\rm res}_{\mu\nu},\qquad
 T^{\Lambda}_{\mu\nu}:=\tfrac14\theta g_{\mu\nu},\qquad
 T^{0}_{\mu\nu}:=T^{\rm res}_{\mu\nu}-T^{\Lambda}_{\mu\nu}.
 \label{x_reciprocidad:vii:eq:descomposicion-traza-residual}
\end{equation}
The superscript \(\Lambda\) identifies the component proportional to the
metric; \(\Lambda_{\rm ref}\) is the reference parameter of
\eqref{x_reciprocidad:vii:eq:residual-definicion}. Their functions are distinct.

\begin{proposition}[Trace projectors and uniqueness of the decomposition]
\label{x_reciprocidad:vii:prop:proyectores-traza-residual}
On symmetric covariant tensors of order two, the maps
\[
 P_{\rm tr}(S)=\tfrac14\operatorname{tr}_g(S)g,
 \qquad P_0(S)=S-P_{\rm tr}(S)
\]
satisfy
\begin{equation}
 P_{\rm tr}^2=P_{\rm tr},\quad P_0^2=P_0,\quad
 P_{\rm tr}P_0=P_0P_{\rm tr}=0,\quad P_{\rm tr}+P_0=I.
 \label{x_reciprocidad:vii:eq:proyectores-traza-residual}
\end{equation}
In particular, \(T^{\rm res}=T^\Lambda+T^0\),
\(\operatorname{tr}_gT^0=0\), and this decomposition is unique.
\end{proposition}
\begin{proof}
In dimension four, \(g^{\mu\nu}g_{\mu\nu}=4\). Therefore,
\(\operatorname{tr}_g(P_{\rm tr}S)=\operatorname{tr}_gS\), which
proves idempotence of \(P_{\rm tr}\), vanishing of the trace
of \(P_0S\) and, by expansion, the other identities.
If \(S=fg+Q\) with \(\operatorname{tr}_gQ=0\), taking the trace
gives \(f=\operatorname{tr}_gS/4\). This also determines
\(Q=S-fg\), and with it uniqueness.
\end{proof}

\begin{theorem}[Covariant exchange between the two sectors]
\label{x_reciprocidad:vii:thm:intercambio-traza-residual}
On the domain of \eqref{x_reciprocidad:vii:eq:residual-definicion}, let
\(b_\nu:=\mathring\nabla^\mu T_{b,\mu\nu}\). The trace-exchange
current is the one-form
\begin{equation}
 \mathcal J^{\rm tr}_\nu
 :=\tfrac14\partial_\nu\theta
 =\mathring\nabla^\mu T^\Lambda_{\mu\nu}.
 \label{x_reciprocidad:vii:eq:corriente-intercambio-traza}
\end{equation}
The complete balances are
\begin{equation}
 \mathring\nabla^\mu T^0_{\mu\nu}
       =-b_\nu-\mathcal J^{\rm tr}_\nu,
 \qquad
 \mathring\nabla^\mu(T^\Lambda+T^0)_{\mu\nu}=-b_\nu.
 \label{x_reciprocidad:vii:eq:balances-sectores-residuales}
\end{equation}
If \(b=0\), the two sectors exchange exactly opposite currents.
In that case they are conserved separately if and only if \(\theta\) is
constant in the connected chart.
\end{theorem}
\begin{proof}
Metric compatibility allows the calculation
\[
 \mathring\nabla^\mu\bigl(\tfrac14\theta g_{\mu\nu}\bigr)
 =\tfrac14g^{\mu\alpha}(\partial_\alpha\theta)g_{\mu\nu}
 =\tfrac14\partial_\nu\theta.
\]
Subtracting this equality from \eqref{x_reciprocidad:vii:eq:balance-residual-material} produces
\eqref{x_reciprocidad:vii:eq:balances-sectores-residuales}. With \(b=0\), separate
conservation is equivalent to \(\mathrm d\theta=0\); a smooth
function with zero differential is constant on each connected component.
\end{proof}

The one-form \(\mathcal J^{\rm tr}\) measures the exchange of the tensor
partition. It retains a type distinct from the transport covector of
section~\ref{x_reciprocidad:vii:sec:variacion-memoria} and from the spin current
\(\sigma_{IJ}\in\Omega^3(\mathcal U;\Lambda^2E^*)\) of
\eqref{x_reciprocidad:vii:eq:corriente-variacional}.

\section{Link to the effective Cartan--Holst source}
\label{x_reciprocidad:vii:subsec:residual-fuente-cartan}

On a solution of \eqref{x_reciprocidad:vii:eq:ecuacion-metrica-efectiva}, in the
same units \(c=1\), the residual has expression
\begin{equation}
 T^{\rm res}
 =T^{\rm phys,LC}+U^{\rm phys,spin}-T_b
   +\frac{\Lambda_{\rm ref}-\Lambda_{\rm int}}{\kappa_x}\,g.
 \label{x_reciprocidad:vii:eq:residual-composicion-cartan}
\end{equation}
Indeed, substituting
\(\mathsf G+\Lambda_{\rm int}g
  =\kappa_x(T^{\rm phys,LC}+U^{\rm phys,spin})\)
into \eqref{x_reciprocidad:vii:eq:residual-definicion} gives equality term by term.
When the material transducer publishes precisely
\(T_b=T^{\rm phys,LC}\) and \(\Lambda_{\rm ref}=\Lambda_{\rm int}\),
we obtain \(T^{\rm res}=U^{\rm phys,spin}\).
The separate conservation condition of \(T_b\) in
theorem~\ref{x_reciprocidad:vii:thm:residual-unicidad-balance} retains its own
function: \eqref{x_reciprocidad:vii:eq:balance-metrico-total} establishes the balance of the
total source, and the partition among its components depends on the material
action used.

\section{Trace sector and acceleration of the selected charts}
\label{x_reciprocidad:vii:subsec:residual-aceleracion-ds}

The TRIT reader of state \(x\) publishes \(s=\tau(x)\in\{+1,0,-1\}\),
and its positive clock publishes
\begin{equation}
 H_*=\frac{2\pi}{108t_0(x)}>0.
 \label{x_reciprocidad:vii:eq:reloj-carta-ds}
\end{equation}
Let \(h_s\) be a three-dimensional metric of constant sectional curvature
\(s\). The three charts of the selected realization are written
\begin{equation}
 g_s=-\mathrm dt^2+a_s(t)^2h_s,\qquad
 a_+(t)=H_*^{-1}\cosh(H_*t),\quad
 a_0(t)=\exp(H_*t),\quad
 a_-(t)=H_*^{-1}\sinh(H_*t).
 \label{x_reciprocidad:vii:eq:cartas-ds-residual}
\end{equation}
The closed and flat charts are considered for \(t\in\mathbb R\);
the open chart, for \(t>0\). The sign \(s\) and the clock \(H_*\)
retain their antecedents in the state selecting the realization.

\begin{lemma}[Curvature of the three charts of the same clock]
\label{x_reciprocidad:vii:lem:curvatura-cartas-ds}
For the functions in \eqref{x_reciprocidad:vii:eq:cartas-ds-residual},
\begin{equation}
 \frac{\ddot a_s}{a_s}=H_*^2,\qquad
 \left(\frac{\dot a_s}{a_s}\right)^2+\frac{s}{a_s^2}=H_*^2,
 \qquad
 \operatorname{Ric}[g_s]=3H_*^2g_s.
 \label{x_reciprocidad:vii:eq:curvatura-cartas-ds}
\end{equation}
Consequently, \(R[g_s]=12H_*^2\) and
\(\mathsf G[g_s]=-3H_*^2g_s\).
\end{lemma}
\begin{proof}
The first two equalities are obtained by differentiating the three functions:
in the closed and open charts we use
\(\cosh^2z-\sinh^2z=1\), and in the flat chart
\(\dot a_0=H_*a_0\). To verify the tensor identification,
the connection coefficients involving the time coordinate are
\[
 \Gamma^0{}_{ij}=a_s\dot a_s(h_s)_{ij},\qquad
 \Gamma^i{}_{0j}=(\dot a_s/a_s)\delta^i_j,
\]
and the purely spatial ones coincide with those of \(h_s\).
Contraction of the curvature gives
\[
 \operatorname{Ric}_{00}=-3\ddot a_s/a_s,\qquad
 \operatorname{Ric}_{0i}=0,\qquad
 \operatorname{Ric}_{ij}
  =(a_s\ddot a_s+2\dot a_s^2+2s)(h_s)_{ij}.
\]
Substituting the first two identities of
\eqref{x_reciprocidad:vii:eq:curvatura-cartas-ds} produces
\(\operatorname{Ric}_{00}=-3H_*^2\) and
\(\operatorname{Ric}_{ij}=3H_*^2a_s^2(h_s)_{ij}\).
Taking the trace and forming the Einstein tensor proves the other
statements, with the curvature convention of
\eqref{x_reciprocidad:vii:eq:ecuacion-metrica-efectiva}.
\end{proof}

\begin{theorem}[Density, pressure, and sign of the selected acceleration]
\label{x_reciprocidad:vii:thm:residual-ds-aceleracion}
In the preceding charts, for \(\Lambda_{\rm ref}=0\) and \(T_b=0\),
the residual tensor is
\begin{equation}
 T^{\rm res}=-\rho_*g_s,\qquad
 \rho_*:=\frac{3H_*^2}{8\pi G_{\rm int}}>0,
 \qquad T^\Lambda=T^{\rm res},\quad T^0=0.
 \label{x_reciprocidad:vii:eq:residual-ds-puro}
\end{equation}
For a unit comoving observer \(u=\partial_t\), the density and
isotropic pressure are \(\rho_*\) and \(p_*=-\rho_*\).
The acceleration equation evaluates to
\begin{equation}
 \frac{\ddot a_s}{a_s}
 =-\frac{4\pi G_{\rm int}}3(\rho_*+3p_*)=H_*^2>0.
 \label{x_reciprocidad:vii:eq:aceleracion-residual-ds}
\end{equation}
\end{theorem}
\begin{proof}
The preceding lemma and \eqref{x_reciprocidad:vii:eq:residual-definicion} give
\(T^{\rm res}=-(3H_*^2/\kappa_x)g_s\).
Since \(g_s(u,u)=-1\),
\(T^{\rm res}_{\mu\nu}u^\mu u^\nu=\rho_*\).
The spatial projector \(q^{\mu\nu}=g_s^{\mu\nu}+u^\mu u^\nu\)
satisfies \(q^{\mu\nu}(g_s)_{\mu\nu}=3\), and therefore
\[
 p_*:=\tfrac13q^{\mu\nu}T^{\rm res}_{\mu\nu}=-\rho_*.
\]
The trace equals \(-4\rho_*\), so
\eqref{x_reciprocidad:vii:eq:descomposicion-traza-residual} produces the two sectors
of \eqref{x_reciprocidad:vii:eq:residual-ds-puro}. Finally,
\(\rho_*+3p_*=-2\rho_*\), and substitution into the right-hand side
of \eqref{x_reciprocidad:vii:eq:aceleracion-residual-ds} gives
\(\kappa_x\rho_*/3=H_*^2\), equal to the geometric value already calculated.
Positivity uses exactly \(G_{\rm int}>0\) and \(H_*>0\).
\end{proof}

In this realization, the trace sector publishes a constant density
and an opposite pressure, and the internal clock fixes positive acceleration.
In a general chart, \(T^\Lambda\) describes the pointwise isotropic
component, while \(T^0\) preserves the traceless residual and both
obey the exchange of \eqref{x_reciprocidad:vii:eq:balances-sectores-residuales}.
Identification of these sectors with observational dark sources
requires a single realization to reproduce expansion,
lensing, rotation, structure growth and the cosmic background jointly. The
present construction fixes the tensor, its balances and the acceleration of
the selected charts entering that comparison.

% END OWNER


% BEGIN OWNER /Users/ruben/Documents/New project/output/EXTREMA_MEDIA_RAZON_RECIPROCIDAD_ES_EN_20260918/ARTICULO/ES/source/antecedentes_vii/70_horizontes_e_informacion.tex
% Recuperacion de 09c_entropia_sectorial.tex, 11_confluencia.tex y
% 85_agujero_negro_doble_hoja.tex; originales en fuentes_conservadas/
% horizontes_informacion. Procedencia: desarrollo/HORIZONTES_INFORMACION_FUENTES.md.
% Las referencias a antecedentes se acuerdan con el editor de VII.
% No incorpora una corriente material ni evalua la amplitud torsional.

\chapter{Pure bipartite states, reduced entropy and horizon area laws}
\label{x_reciprocidad:vii:sec:horizontes-informacion}

The enriched state preserves routes and their records during
evolution; a subsystem reading determines which part of that information
remains accessible. The construction of this section uses the angular
channels of section~\ref{x_reciprocidad:iv:angular:section}, the incidence and the
area operator of section~\ref{x_reciprocidad:v:ant:sec:area}, and the action,
area and temperature units of
sections~\ref{x_reciprocidad:v:ant:sec:gravity}
and~\ref{x_reciprocidad:v:ant:sec:ii-kb-aut-desarrollo}. These are prior publications of the
same APP--TRIT--TPK construction. Each here transmits its domain,
orientation, normalization and realization conditions.

Three spaces are distinguished: the binary sheet record, the tritic
factors of sectorial incidence, and the geometric realization of the
exterior regions of a horizon. Purification, transport of observables,
and dimensional conversion specify their relationships. This
distinction preserves the structure accompanying each entropic value.

\section{Sheet record and bipartite quantum purification}
\label{x_reciprocidad:vii:subsec:hojas-purificacion}

Write \(a=A_{\rm rad}\) and \(y=C^*_{\rm rad}\) for the
angular coordinates already constructed. The channels and their probabilities are
\begin{equation}
 \mathfrak q_\pm=e^{-a\mp y},\qquad
 p_\pm=\frac{\mathfrak q_\pm}{\mathfrak q_++\mathfrak q_-}
       =\frac{e^{\mp y}}{2\cosh y}.
 \label{x_reciprocidad:vii:eq:probabilidades-hoja}
\end{equation}
The signs index these two channels. On
\(\mathcal H_{\rm h}=\operatorname{span}\{|+\rangle,|-\rangle\}\),
define
\[
 \rho_{\rm h}(y)=p_+|+\rangle\langle+|+p_-|-\rangle\langle-|.
\]
The auxiliary space
\(\mathcal H_{\rm ar}=\operatorname{span}
\{|\mathrm{adv}\rangle,|\mathrm{ret}\rangle\}\)
carries two orthonormal labels. With a phase \(\chi_0\) fixed, the joint
state is
\begin{equation}
 |\Psi_{\rm ar}\rangle
 =\sqrt{p_+}|+\rangle\otimes|\mathrm{adv}\rangle
  +e^{i\chi_0}\sqrt{p_-}|-\rangle\otimes|\mathrm{ret}\rangle.
 \label{x_reciprocidad:vii:eq:purificacion-hoja}
\end{equation}

\begin{proposition}[Accessible information in the sheet record]
\label{x_reciprocidad:vii:prop:purificacion-hoja}
The vector~\eqref{x_reciprocidad:vii:eq:purificacion-hoja} is normalized, and its partial
trace over \(\mathcal H_{\rm ar}\) is \(\rho_{\rm h}(y)\).
Its entanglement entropy, expressed in nats, is
\begin{equation}
 s_{\rm h}(y)=\log(2\cosh y)-y\tanh y.
 \label{x_reciprocidad:vii:eq:entropia-hoja}
\end{equation}
Exchange of the labels \(+\) and \(-\) conjugates
\(\rho_{\rm h}(y)\) to \(\rho_{\rm h}(-y)\) and preserves
\(s_{\rm h}\).
\end{proposition}
\begin{proof}
The squared norm is \(p_++p_-=1\). Orthogonality of the
auxiliary labels eliminates the cross terms in the partial trace.
The squared Schmidt coefficients are precisely \(p_\pm\).
By~\eqref{x_reciprocidad:vii:eq:probabilidades-hoja},
\(p_+-p_-=-\tanh y\); substituting
\(\log p_\pm=\mp y-\log(2\cosh y)\) into
\(-\sum_\pm p_\pm\log p_\pm\) proves the formula.
Exchange of signs permutes the two eigenvalues and preserves entropy.
\end{proof}

The common factor \(e^{-a}\) cancels upon normalization of the binary record.
The relative phase \(\chi_0\) remains in the joint state even though this
diagonal reduction does not express it. Logical reversibility concerns
transport of the complete record, with its labels and correlations;
phase return preserves the memory advance described in
Section~\ref{x_reciprocidad:vii:sec:retorno-memoria-gravedad}. A thermal implementation
additionally has its own energy and dissipation law. In particular,
\(s_{\rm h}\) measures this binary subsystem and its concrete bipartition.

\section{Quantum purification of the incidence sectors}
\label{x_reciprocidad:vii:subsec:purificacion-sectorial}

The boundary incidence constructed earlier distinguishes two sets
of eighteen links, with sectorial labels \(m=90,120\). Each link
carries \(\mathbb C^3\), with basis \(|0\rangle,|1\rangle,|2\rangle\)
and occupation operator \(\operatorname{diag}(0,1,2)\). Let
\(\mathsf N_m\) be their sectorial sums, and
\begin{equation}
 \mathsf N=\mathsf N_{90}+\mathsf N_{120},\qquad
 \mathsf D_\partial=90\mathsf N_{90}+120\mathsf N_{120}.
 \label{x_reciprocidad:vii:eq:ocupaciones-sectoriales}
\end{equation}
In the positive branch, \(\gamma_{\rm BI}\) denotes the Barbero--Immirzi
parameter already obtained, and \(a_0>0\) its angular areal factor.
This notation distinguishes the parameter from the connection and its holonomies.
Area is published as
\begin{equation}
 \widehat{\mathcal A}_m=\ell_P^2\gamma_{\rm BI}a_0\mathsf N_m,
 \qquad
 \widehat{\mathcal A}=\widehat{\mathcal A}_{90}
                         +\widehat{\mathcal A}_{120}.
 \label{x_reciprocidad:vii:eq:area-sectores}
\end{equation}

For the finite parameter \(\Delta=\Delta_{\rm mod}\in\mathbb R\) of the
reduced state, set
\begin{align}
 x_m&=e^{-m\Delta},& z_m&=1+x_m+x_m^2,&
 Z(\Delta)&=z_{90}^{18}z_{120}^{18},\\
 \rho_\partial(\Delta)&=Z(\Delta)^{-1}e^{-\Delta\mathsf D_\partial}.
 \label{x_reciprocidad:vii:eq:estado-sectorial}
\end{align}
The trace in this section is the ordinary matrix trace, and the states
have trace one. The parameter \(\Delta\) preserves the selection and
normalization of the antecedent state.

\begin{proposition}[Tensor purification of the boundary]
\label{x_reciprocidad:vii:prop:purificacion-sectorial}
For a finite real weight \(w\), let
\begin{equation}
 |\operatorname{TFD}(w)\rangle
 =\frac{|0\rangle|0\rangle+e^{-w/2}|1\rangle|1\rangle
                   +e^{-w}|2\rangle|2\rangle}
        {\sqrt{1+e^{-w}+e^{-2w}}}.
 \label{x_reciprocidad:vii:eq:tfd-local}
\end{equation}
The product
\begin{equation}
 |\Psi_\Delta\rangle
 =|\operatorname{TFD}(90\Delta)\rangle^{\otimes18}
  \otimes|\operatorname{TFD}(120\Delta)\rangle^{\otimes18}
 \label{x_reciprocidad:vii:eq:tfd-sectorial}
\end{equation}
is a normalized purification of \(\rho_\partial(\Delta)\).
\end{proposition}
\begin{proof}
At each link, the sum of the squares of the three coefficients is one.
The partial trace eliminates terms with different indices and returns
\(\operatorname{diag}(1,e^{-w},e^{-2w})/(1+e^{-w}+e^{-2w})\).
Partial trace commutes with the tensor product. Multiplying the
thirty-six reductions gives the exponential and partition function
of~\eqref{x_reciprocidad:vii:eq:estado-sectorial}.
\end{proof}

\section{Informational equation and areal fluctuations}
\label{x_reciprocidad:vii:subsec:ecuacion-informacional}

The occupation means and variances per link are
\begin{equation}
 f_m=\frac{x_m+2x_m^2}{z_m},\qquad
 v_m=\frac{x_m+4x_m^2}{z_m}-f_m^2.
 \label{x_reciprocidad:vii:eq:momentos-sectoriales}
\end{equation}
Write \(s(\rho)=-\operatorname{Tr}(\rho\log\rho)\), with
\(0\log0=0\). In the preceding purification, \(s(\rho_\partial)\)
is its entanglement entropy.

\begin{proposition}[Mean area, entropy, and sectorial response]
\label{x_reciprocidad:vii:prop:respuesta-sectorial}
For real finite \(\Delta\), one has
\begin{align}
 \langle\mathsf N\rangle_\Delta&=18(f_{90}+f_{120}),\\
 \langle\mathsf D_\partial\rangle_\Delta&=18(90f_{90}+120f_{120}),\\
 \langle\widehat{\mathcal A}\rangle_\Delta
 &=18\ell_P^2\gamma_{\rm BI}a_0(f_{90}+f_{120}),\\
 s(\Delta)&=\log Z(\Delta)+\Delta\langle\mathsf D_\partial\rangle_\Delta,
 \label{x_reciprocidad:vii:eq:entropia-sectorial}\\
 \frac{\mathrm d\langle\mathsf D_\partial\rangle_\Delta}{\mathrm d\Delta}
 &=-\operatorname{Var}_\Delta(\mathsf D_\partial)
  =-18(90^2v_{90}+120^2v_{120}),\\
 \frac{\mathrm ds}{\mathrm d\Delta}
 &=-\Delta\operatorname{Var}_\Delta(\mathsf D_\partial),\\
 \frac{\mathrm d\langle\widehat{\mathcal A}\rangle_\Delta}{\mathrm d\Delta}
 &=-18\ell_P^2\gamma_{\rm BI}a_0(90v_{90}+120v_{120}).
 \label{x_reciprocidad:vii:eq:respuestas-area-entropia}
\end{align}
\end{proposition}
\begin{proof}
The local weights are \((1,x_m,x_m^2)/z_m\), whose first two
moments give \(f_m\) and \(v_m\). Additivity of occupation
operators and factorization of the state give the first three formulas.
Taking the trace of
\(-\rho_\partial\log\rho_\partial
=\rho_\partial(\log Z+\Delta\mathsf D_\partial)\)
gives the entropy. Differentiation of the finite sums yields
\(\partial_\Delta\log Z=-\langle\mathsf D_\partial\rangle_\Delta\)
and \(\partial_\Delta^2\log Z
=\operatorname{Var}_\Delta(\mathsf D_\partial)\).
Independent factors add their variances. In the derivative of
\eqref{x_reciprocidad:vii:eq:entropia-sectorial}, the two first-moment
terms cancel. Finally, \(\partial_\Delta f_m=-mv_m\)
gives the area response.
\end{proof}

Loss of uniformity is quantified in dimension \(d_\partial=3^{36}\):
\begin{equation}
 I_\partial(\Delta)
 =36\log3-s(\Delta)
 =\mathcal D\bigl(\rho_\partial(\Delta)\Vert I/d_\partial\bigr)\ge0.
 \label{x_reciprocidad:vii:eq:deficit-informacional}
\end{equation}
Here \(\mathcal D\) denotes relative entropy, defined and justified
below. Substituting \(\log(I/d_\partial)=-36\log3\,I\) proves the
identity. Convexity of \(t\log t\) gives positivity; equality
corresponds to the uniform distribution, which in this family occurs
exactly at \(\Delta=0\). The deficit measures dimensionless information;
the scale \(\gamma_{\rm BI}a_0\ell_P^2\) measures area per occupation.

\begin{theorem}[Occupation inversion and complementary area]
\label{x_reciprocidad:vii:thm:inversion-sectorial}
Let \(\mathcal R_\partial\) be the involution exchanging
\(|0\rangle\) and \(|2\rangle\) and fixing \(|1\rangle\) in each
factor. Then
\begin{align}
 \mathcal R_\partial\mathsf N\mathcal R_\partial^*&=72I-\mathsf N,\\
 \mathcal R_\partial\mathsf D_\partial\mathcal R_\partial^*
 &=7560I-\mathsf D_\partial,\\
 \rho_\partial(-\Delta)&=\mathcal R_\partial\rho_\partial(\Delta)
                          \mathcal R_\partial^*,\\
 s(-\Delta)&=s(\Delta),\\
 \langle\widehat{\mathcal A}\rangle_{-\Delta}
 +\langle\widehat{\mathcal A}\rangle_\Delta
 &=72\gamma_{\rm BI}a_0\ell_P^2,\\
 \mathcal D\bigl(\rho_\partial(\Delta)\Vert\rho_\partial(-\Delta)\bigr)
 &=\Delta\bigl(7560-2\langle\mathsf D_\partial\rangle_\Delta\bigr).
 \label{x_reciprocidad:vii:eq:inversion-sectorial}
\end{align}
\end{theorem}
\begin{proof}
Each occupation transforms by \(n\mapsto2-n\). In each sector,
\(\mathsf N_m\mapsto36I-\mathsf N_m\), so the constant
term of the weighted reader is \(36(90+120)=7560\).
Moreover, \(z_m(-\Delta)=e^{2m\Delta}z_m(\Delta)\) and
\(Z(-\Delta)=e^{7560\Delta}Z(\Delta)\). Substitution into the
normalized state proves conjugation. Its spectrum preserves entropy;
the transformation of \(\mathsf N\) gives complementary area. Subtracting
the logarithms of the two states and taking the expectation in
\(\rho_\partial(\Delta)\) gives the final equality.
\end{proof}

At \(\Delta=0\), \(s=36\log3\). As \(\Delta\to+\infty\),
zero occupation concentrates the state and entropy tends to zero;
inversion carries that limit to maximal occupation. For \(\Delta>0\),
entropy decreases strictly because the weighted reader has more
than one eigenvalue and the state has full rank. The complete configuration,
area, and entropy thus preserve differentiated roles.

\section{Modular law and finite area variations}
\label{x_reciprocidad:vii:subsec:ley-modular-finita}

Fix \(\rho_0=\rho_\partial(\Delta_0)\), with
real finite \(\Delta_0\), and hold
\(\Delta_0,\gamma_{\rm BI},a_0,\ell_P\) constant. The normalized areas
\(\mathcal A_m^{\rm n}=\gamma_{\rm BI}a_0\mathsf N_m\)
express the modular operator as
\begin{equation}
 \mathsf K_0=-\log\rho_0
 =c_0I+\sum_{m=90,120}\beta_m\mathcal A_m^{\rm n},
 \quad \beta_m=\frac{m\Delta_0}{\gamma_{\rm BI}a_0},
 \quad c_0=\log Z(\Delta_0).
 \label{x_reciprocidad:vii:eq:operador-modular-areal}
\end{equation}
One has \(\beta_{120}=\tfrac43\beta_{90}\), including when both
are zero. The ratio is \(4/3\) when \(\Delta_0\ne0\).
\(\mathsf K_0\) is the logarithm of this reference state; its type
distinguishes it from temporal transport and the gravitational connection.

\begin{theorem}[First modular law and finite correction]
\label{x_reciprocidad:vii:thm:ley-modular-finita}
For any normalized state \(\sigma\) on the same space, let
\(\Delta_\sigma\langle B\rangle
=\operatorname{Tr}[(\sigma-\rho_0)B]\). Then
\begin{equation}
 s(\sigma)-s(\rho_0)
 =\sum_{m=90,120}\beta_m\Delta_\sigma\langle\mathcal A_m^{\rm n}\rangle
  -\mathcal D(\sigma\Vert\rho_0),
 \label{x_reciprocidad:vii:eq:ley-modular-finita}
\end{equation}
where
\(\mathcal D(\sigma\Vert\rho_0)
=\operatorname{Tr}[\sigma(\log\sigma-\log\rho_0)]\)
is finite and nonnegative. If \(\sigma(\varepsilon)\) is differentiable,
normalized, and \(\sigma(0)=\rho_0\), one obtains
\begin{equation}
 \left.\frac{\mathrm ds(\sigma(\varepsilon))}{\mathrm d\varepsilon}\right|_0
 =\sum_{m=90,120}\beta_m
   \left.\frac{\mathrm d\langle\mathcal A_m^{\rm n}\rangle_{
                 \sigma(\varepsilon)}}{\mathrm d\varepsilon}\right|_0.
 \label{x_reciprocidad:vii:eq:primera-ley-modular}
\end{equation}
\end{theorem}
\begin{proof}
By definition,
\(\mathcal D(\sigma\Vert\rho_0)=-s(\sigma)
+\operatorname{Tr}(\sigma\mathsf K_0)\), whereas
\(s(\rho_0)=\operatorname{Tr}(\rho_0\mathsf K_0)\).
Subtraction and~\eqref{x_reciprocidad:vii:eq:operador-modular-areal} prove the finite
identity, since normalization cancels the term \(c_0I\).

For positivity, diagonalize \(\sigma\) with eigenvalues
\(p_i\) and vectors \(u_i\), and \(\rho_0\) with eigenvalues
\(q_j>0\) and vectors \(v_j\). If
\(r_i=\langle u_i,\rho_0u_i\rangle>0\), concavity of the
logarithm gives
\[
 \sum_j|\langle u_i,v_j\rangle|^2\log q_j\le\log r_i.
\]
Consequently,
\(\mathcal D(\sigma\Vert\rho_0)
\ge\sum_i p_i\log(p_i/r_i)\ge0\).
The final inequality is convexity of \(t\log t\) with weights
\(r_i\), whose sum is one. Faithfulness of \(\rho_0\) keeps its
logarithms finite; zero eigenvalues of \(\sigma\) are
interpreted through \(0\log0=0\).

For the infinitesimal formula, in a neighbourhood of the faithful state \(\rho_0\),
the derivative of the trace of \(-\sigma\log\sigma\) is
\(-\operatorname{Tr}[\dot\sigma(\log\rho_0+I)]\).
This identity follows from the spectral derivative of the trace.
The condition \(\operatorname{Tr}\dot\sigma=0\) leaves
\(\operatorname{Tr}(\dot\sigma\mathsf K_0)\), and substitution of
\eqref{x_reciprocidad:vii:eq:operador-modular-areal} completes the proof.
\end{proof}

For physical areas, the coefficients are \(\beta_m/\ell_P^2\).
Applying the common transducer \(k_B\) gives \(S=k_Bs\) and
preserves the term \(k_B\mathcal D\) in full. The derivative in
\eqref{x_reciprocidad:vii:eq:respuestas-area-entropia} varies \(\Delta\); law
\eqref{x_reciprocidad:vii:eq:primera-ley-modular} holds \(\Delta_0\) fixed and
varies the state around the reference. These are different operations.

\section{Action, decision energy, and areal unit}
\label{x_reciprocidad:vii:subsec:conversion-informacion-area}

We fix the same oriented action section and its gravitational
and thermal sections, on the positive branch. The antecedents provide
\begin{equation}
 \begin{aligned}
 \ell_P^2&=\frac{\hbar G}{c^3},& t_P&=\frac{\ell_P}{c},\\
 T_{108}&=2\pi t_P,& h&=2\pi\hbar,\\
 E_P&=\frac{\hbar}{t_P}=k_B\Theta_{\rm clk}\log3.&&
 \end{aligned}
 \label{x_reciprocidad:vii:eq:unidad-accion-informacion}
\end{equation}
The thermal equality preserves the previously constructed pair and
normalization; changing its dimensional basis transports
\(k_B\) and \(\Theta_{\rm clk}\) jointly. The elementary area increment is
\(\delta\mathcal A=\gamma_{\rm BI}a_0\ell_P^2\).

\begin{proposition}[Conversion coefficient between decision and area]
\label{x_reciprocidad:vii:prop:tension-informacion-area}
The ratio of decision energy to the areal increment satisfies
\begin{equation}
 \begin{aligned}
 \mathcal T_{I/A}:=\frac{E_P}{\delta\mathcal A}
 &=\frac{k_B\Theta_{\rm clk}\log3}{\gamma_{\rm BI}a_0\ell_P^2}\\
 &=\frac{c^3}{\gamma_{\rm BI}a_0G t_P}
  =\frac{c^4}{\gamma_{\rm BI}a_0G\ell_P}.
 \end{aligned}
 \label{x_reciprocidad:vii:eq:tension-informacion-area}
\end{equation}
Its dimension is energy per area.
\end{proposition}
\begin{proof}
Substitution of \(E_P=\hbar/t_P\) and
\(\ell_P^2=\hbar G/c^3\) cancels the common action section.
The identity \(\ell_P=ct_P\) gives the final expression. Positivity
of the factors ensures that all quotients are defined.
\end{proof}

The capacity \(81\) of the cubic cut and the bipartition of the thirty-
six sectoral links correspond to different domains. In the
cubic graph \(\{0,\ldots,N-1\}^3\), the \(N^2\) parallel lines
between two opposite faces are edge-disjoint: every cut
intercepts the \(N^2\), and a transverse layer attains that bound.
For \(N=9\), the uniform product state of the trits of that layer
has \(s_{\rm corte}=81\log3\). Assigning a decision of the
same cycle to each link publishes
\begin{equation}
 S_{\rm corte}^{\rm fis}=81k_B\log3,\qquad
 E_{\rm corte}^{(1)}=81E_P
                  =\Theta_{\rm clk}S_{\rm corte}^{\rm fis}.
 \label{x_reciprocidad:vii:eq:corte-energia-informacion}
\end{equation}
The proof combines the minimum cut, additivity of product-state
entropy, and~\eqref{x_reciprocidad:vii:eq:unidad-accion-informacion}. In the sectorial
state, the relevant finite equation remains
\eqref{x_reciprocidad:vii:eq:ley-modular-finita}, with its relative entropy.

\section{Cosmological normalization of the horizon}
\label{x_reciprocidad:vii:subsec:balance-horizonte}

For an areal radius \(R>0\), the cosmological realization considered
uses the triad
\begin{equation}
 E_\Lambda(R)=\frac{c^4R}{2G},\qquad
 S_{\rm H}(R)=\frac{k_B\pi R^2}{\ell_P^2},\qquad
 T_{\rm cos}(R)=\frac{\hbar c}{2\pi k_BR}.
 \label{x_reciprocidad:vii:eq:horizonte-cosmologico}
\end{equation}
The factor \(2\pi\) belongs to this temperature normalization.
The area \(\mathcal A_{\rm H}=4\pi R^2\) expresses the same entropy
as \(S_{\rm H}/k_B=\mathcal A_{\rm H}/(4\ell_P^2)\).

\begin{theorem}[Gravitational half-action cell]
\label{x_reciprocidad:vii:thm:celda-media-accion}
In~\eqref{x_reciprocidad:vii:eq:horizonte-cosmologico},
\begin{align}
 T_{\rm cos}(R)S_{\rm H}(R)&=E_\Lambda(R),\\
 T_{\rm cos}(\ell_P)S_{\rm H}(\ell_P)
 &=\frac{E_P}{2}=\frac{k_B\Theta_{\rm clk}\log3}{2},\\
 E_\Lambda(\ell_P)t_P&=\frac{\hbar}{2},&
 E_\Lambda(\ell_P)T_{108}&=\frac h2.
 \label{x_reciprocidad:vii:eq:media-accion-horizonte}
\end{align}
Keeping the chart constants fixed and varying \(R\),
\begin{equation}
 T_{\rm cos}\,\mathrm dS_{\rm H}=2\,\mathrm dE_\Lambda,
 \qquad S_{\rm H}\,\mathrm dT_{\rm cos}=-\mathrm dE_\Lambda.
 \label{x_reciprocidad:vii:eq:diferencial-horizonte}
\end{equation}
\end{theorem}
\begin{proof}
Multiplying temperature and entropy gives
\(\hbar cR/(2\ell_P^2)=c^4R/(2G)\).
At \(R=\ell_P\), the result is
\(\hbar/(2t_P)=E_P/2\). The two durations in
\eqref{x_reciprocidad:vii:eq:unidad-accion-informacion} give the action relations.
For the differentials,
\(\mathrm dE_\Lambda=c^4\mathrm dR/(2G)\),
\(\mathrm dS_{\rm H}=2\pi k_BR\mathrm dR/\ell_P^2\), and
\(\mathrm dT_{\rm cos}=-T_{\rm cos}\mathrm dR/R\).
Substitution proves both identities.
\end{proof}

The product \(E_\Lambda=T_{\rm cos}S_{\rm H}\) expresses a
state identity. Its radial variations preserve both terms
of~\eqref{x_reciprocidad:vii:eq:diferencial-horizonte}; a complete thermodynamic interpretation
additionally specifies the corresponding time generator, work and
material realization. In this chart, action and
gravitation fix the unit \(\ell_P^2\), while Boltzmann converts
information into dimensional entropy. Transport of sections with
\(G\hbar\) fixed preserves that unit when the section \(c\) is
kept fixed, and preserves \(S_{\rm H}/k_B\) at fixed radius.

In the cell \(R=\ell_P\), \(S_{\rm H}/k_B=\pi\), so
the effective multiplicity is \(\Omega_{\rm eff}=e^\pi\).
For the tritic hyperbolic involution \(H_{\rm trit}\), with spectrum
\(\{1,-1\}\), diagonalization gives
\(\operatorname{spec}(e^{\pi H_{\rm trit}})=\{e^\pi,e^{-\pi}\}\).
This spectral relation specifies the expansive value; a microscopic
identification of both spaces additionally preserves their realization
map and the state.

\section{Kruskal chart and exchange of the exteriors}
\label{x_reciprocidad:vii:subsec:kruskal-dos-hojas}

Let us fix a published state \(x\), its radius \(R=R(x)>0\) and the
positive sections \(c,G\). In the following chart, \(x,R,c,G\)
are fixed parameters; \(t,r\) are the spacetime coordinates.
The exterior realization of the two sheets is specified by
\begin{equation}
 f(r)=1-\frac Rr,\qquad
 \mathrm ds^2=-f(r)c^2\mathrm dt^2+f(r)^{-1}\mathrm dr^2
                 +r^2\mathrm d\Omega_2^2,\qquad r>R.
 \label{x_reciprocidad:vii:eq:metrica-exterior}
\end{equation}
For the sheet reader \(\sigma\in\{+1,-1\}\), define
\begin{align}
 r_*&=r+R\log\left(\frac rR-1\right),\\
 \mathcal U&=-\sigma\exp\left(-\frac{ct-r_*}{2R}\right),&
 \mathcal V&=\sigma\exp\left(\frac{ct+r_*}{2R}\right).
 \label{x_reciprocidad:vii:eq:carta-kruskal}
\end{align}
The symbols \(\mathcal U,\mathcal V\) are coordinates and are
distinguished from route-transport operators.

\begin{theorem}[Geometric realization and sheet involution]
\label{x_reciprocidad:vii:thm:kruskal-hojas}
Chart~\eqref{x_reciprocidad:vii:eq:carta-kruskal} realizes the two sheets as the
two exteriors \(\mathcal U\mathcal V<0\) of the Kruskal
extension. We have
\begin{equation}
 -\mathcal U\mathcal V=\left(\frac rR-1\right)e^{r/R},\qquad
 \mathrm ds^2=-\frac{4R^3}{r}e^{-r/R}\,
                    \mathrm d\mathcal U\,\mathrm d\mathcal V
                    +r^2\mathrm d\Omega_2^2.
 \label{x_reciprocidad:vii:eq:metrica-kruskal}
\end{equation}
The involution
\begin{equation}
 \mathcal C_{\rm H}:(\sigma,\mathcal U,\mathcal V)
 \longmapsto(-\sigma,-\mathcal U,-\mathcal V)
 \label{x_reciprocidad:vii:eq:involucion-horizonte}
\end{equation}
is isometric, fixes the areal radius, exchanges the exterior regions, and preserves
the horizon set \(\mathcal U\mathcal V=0\).
The metric is regular at \(r=R\) and satisfies the vacuum equations
on the regular domain \(r>0\).
\end{theorem}
\begin{proof}
Multiplying the coordinates gives the first identity. Since
\(\mathrm dr_*/\mathrm dr=f^{-1}\),
\[
 \mathrm d\mathcal U=\frac{\mathcal U}{2R}
            (-c\,\mathrm dt+f^{-1}\mathrm dr),\qquad
 \mathrm d\mathcal V=\frac{\mathcal V}{2R}
            ( c\,\mathrm dt+f^{-1}\mathrm dr).
\]
Their product reproduces the radial and temporal part of the metric.
The function \(z\mapsto(z-1)e^z\) has derivative \(ze^z>0\)
for \(z>0\), determining \(r\) from the product.
In the exterior region one recovers
\(\sigma=\operatorname{sign}\mathcal V\) and
\(t=(R/c)\log(-\mathcal V/\mathcal U)\).
At \(r=R\), the radial coefficient of the extended metric is
\(-4R^2/e\), finite and nonzero. Inversion preserves the
product, radius, and \(\mathrm d\mathcal U\mathrm d\mathcal V\),
and exchanges the two sign choices.

To check the vacuum in the exterior chart, the independent mixed
components of the Einstein tensor reduce to
\[
 \mathsf G^t{}_t=\mathsf G^r{}_r
 =\frac{rf'+f-1}{r^2},\qquad
 \mathsf G^\theta{}_\theta=\mathsf G^\phi{}_\phi
 =\frac{f'}r+\frac{f''}2.
\]
With \(f'=R/r^2\) and \(f''=-2R/r^3\), all vanish.
The regular expression~\eqref{x_reciprocidad:vii:eq:metrica-kruskal} prolongs the
result through the horizons on the domain \(r>0\).
\end{proof}

This realization preserves
\(E_{\rm H}=c^4R/(2G)\) and
\(S_{\rm H}=k_B\pi R^2/\ell_P^2\), which depend on the areal radius.
The temperature \(T_{\rm cos}\) of
\eqref{x_reciprocidad:vii:eq:horizonte-cosmologico} belongs to the cosmological
realization declared there. The exterior chart
\eqref{x_reciprocidad:vii:eq:metrica-exterior} preserves its own temporal normalization.
The geometric map establishes the exchange of exteriors; comparison
with formation by collapse, radiation spectra or quasinormal modes
additionally uses their respective physical conditions.

\section{Transport of states, observables, and refinements}
\label{x_reciprocidad:vii:subsec:transporte-informacion-horizonte}

Let \(F_x\) denote the chart of the two exteriors just constructed.
In each regular angular chart, its temporal and radial Jacobian is
\begin{equation}
 \left|\frac{\partial(\mathcal U,\mathcal V)}{\partial(t,r)}\right|
 =\frac{cr}{2R^3}e^{r/R}>0.
 \label{x_reciprocidad:vii:eq:jacobiano-kruskal}
\end{equation}
The identity follows from the two differentials of the preceding proof.
The remaining state records are transported as fiber
coordinates; the chart preserves its identity and memory.

\begin{proposition}[Unitary equivalence of sheet and exterior readings]
\label{x_reciprocidad:vii:prop:transporte-unitario-horizonte}
Let \(\mu\) be a sheet measure and \(\nu=(F_x)_*\mu\) its
transported measure. Change of variables defines a unitary map
\begin{equation}
 W_{\rm H}:L^2(\mu)\longrightarrow L^2(\nu),\qquad
 (W_{\rm H}\psi)(z)=\psi(F_x^{-1}z).
 \label{x_reciprocidad:vii:eq:unitaria-horizonte}
\end{equation}
With densities relative to fixed coordinate measures, the same
map incorporates the corresponding Jacobian factor.
Conjugation
\begin{equation}
 \mathfrak I_{\rm H}(B)=W_{\rm H}BW_{\rm H}^*
 \label{x_reciprocidad:vii:eq:mapa-observables-horizonte}
\end{equation}
is a unital completely positive \(*\)-isomorphism. It transports the
state, its expectations, and its modular functional calculus. For the
symmetric measure of the two sheets, it also intertwines their involutions.
\end{proposition}
\begin{proof}
The definition of pushforward measure gives
\(\int|\psi\circ F_x^{-1}|^2\,\mathrm d\nu
=\int|\psi|^2\,\mathrm d\mu\).
The chart is invertible in the exteriors, so the isometry is
surjective. Unitary conjugation preserves product, adjoint and
identity; in each matrix amplification it conjugates by
\(I_n\otimes W_{\rm H}\), proving complete positivity.
For a state \(\rho\), cyclicity of the trace gives
\(\operatorname{Tr}(W_{\rm H}\rho W_{\rm H}^*
\,\mathfrak I_{\rm H}(B))=\operatorname{Tr}(\rho B)\).
Functional calculus implies
\(g(W_{\rm H}\rho W_{\rm H}^*)=W_{\rm H}g(\rho)W_{\rm H}^*\)
on its domain, in particular for the logarithm and modular flow of
a state faithful on its support. The equality
\(F_x\circ\mathcal C_{\rm hoja}=\mathcal C_{\rm H}\circ F_x\)
proves the intertwining of involutions; the symmetric measure realizes them
unitarily.
\end{proof}

% REMATE_LOCAL_X_HORIZONTE
\Needspace{25\baselineskip}
If the refinements preserve \((\sigma,t,r,R)\) and use the same
isometry on the fiber registers, the inclusions
\(j_m^{\rm hoja}\) and \(j_m^{\rm ext}\) satisfy
\begin{equation}
 W_{{\rm H},m+1}j_m^{\rm hoja}
 =j_m^{\rm ext}W_{{\rm H},m}.
 \label{x_reciprocidad:vii:eq:naturalidad-horizonte}
\end{equation}
Indeed, change of variables acts on the preserved geometric coordinates
and inclusion on the same fibre coordinates; evaluating
both compositions gives the same function and the same transported density.
The equality extends the map to the refined system without resetting the records.

For a bipartition \(\mathcal H_A\otimes\mathcal H_B\), a
factorized transport \(U_A\otimes U_B\) satisfies
\[
 \rho'_A=U_A\rho_AU_A^*,\qquad s(\rho'_A)=s(\rho_A).
\]
The first identity follows from partial trace and the second from the spectrum.
For a general unitary, the equivalent formulation also transports
the subsystem algebra:
\(\mathcal B(\mathcal H_A)\otimes I\mapsto
U(\mathcal B(\mathcal H_A)\otimes I)U^*\).
This preserves the same operational question concerning accessible
information. Entropy equalities apply to states with the relevant
finite entropies.

The relation between information and gravitation is expressed through
distinct composable operations: the state and bipartition fix
accessible information; incidence and Barbero--Immirzi fix areal
resolution; action and gravitation construct the unit \(\ell_P^2\);
Boltzmann realizes entropy dimensionally; the horizon chart
transports geometry and its observables. The finite laws preserve
relative entropy, and the geometric realizations preserve their measures,
temporal normalization and domains.

% END OWNER


% BEGIN OWNER /Users/ruben/Documents/New project/output/EXTREMA_MEDIA_RAZON_RECIPROCIDAD_ES_EN_20260918/ARTICULO/ES/source/antecedentes_vii/80_observador_y_respuesta_relacional.tex
% Recuperación y articulación de los propietarios documentados en
% PROCEDENCIA_OBSERVADOR_RESPUESTA.md. Originales conservados íntegros.
% Antecedentes de integración: núcleo APP--TRIT--TPK, realización del
% transporte de borde, lectores areal/volumétrico y reloj energético.
% Las condiciones de representación se enuncian antes de cada uso.

\chapter{Observer, memory, and relational response}
\label{x_reciprocidad:vii:sec:observador-respuesta}

The reading operation acts on the transport published by the
APP--TRIT--TPK enriched state. Its domain preserves route, sheet, orientation,
carry and memory; the result accessible to the observer is a map
of that domain. This distinction allows three operations to be described
together: choosing an oriented representative of transport,
inscribing its record in an apparatus and the local response induced
by the global boundary. Each operation transmits different information.
The following construction specifies their maps and the conditions
allowing them to be composed.

\section{Boundary transport and intrinsic reading}
\label{x_reciprocidad:vii:sec:observador-ciclo}

The finite publication of transport provides an oriented cycle
\(C_n\), with \(n\geq1\), vertices \(0,\ldots,n-1\) and edges
\(e_j:j\longrightarrow j+1\pmod n\). We fix the lattice realization
\(L\) of the record, its vector space
\(W=L\otimes_{\mathbb Z}\mathbb R\), a positive inner product and
the nonzero primitive representative \(v\in L\) of the distinguished residual
class, in the declared chamber. The APP sum and product sheets
provide the transport polarizations; TRIT orientation and
TPK updating precede this linear publication. The cycle,
the lattice realization and the residual class are the data this
section receives from the transport constructed previously.

A \(0\)-cochain is a function on the vertices with values in
\(W\); a \(1\)-cochain assigns a vector to each oriented edge.
The coboundary acts by
\[
 (d\Phi)(e_j)=\Phi(j+1)-\Phi(j).
\]
We denote by \(\mathbf1_E\) the constant scalar cochain on the
edges. For a transport \(\omega\in C^1(C_n;W)\), its accumulated
publication \(\Gamma_\omega\) is obtained through partial sums and
affine interpolation; its closure defect is
\begin{equation}
 B(\omega)=\Gamma_\omega(1)-\Gamma_\omega(0)
 =\sum_{j=0}^{n-1}\omega(e_j).
 \label{x_reciprocidad:vii:eq:observador-defecto}
\end{equation}

\begin{theorem}[Exact decomposition and harmonic mode of the cycle]
\label{x_reciprocidad:vii:thm:observador-descomposicion}
Every transport has a unique decomposition
\[
 \omega=d\Phi+u\mathbf1_E,\qquad \Phi(0)=0,
 \qquad u=\frac1n\sum_j\omega(e_j).
\]
In particular, \(B(d\Phi)=0\) and \(B(\omega)=nu\).
\end{theorem}

\begin{proof}
Let \(y_j=\omega(e_j)-u\). The sum of the \(y_j\) is zero.
We define \(\Phi(k)=\sum_{j=0}^{k-1}y_j\); the value for \(k=n\)
coincides with \(\Phi(0)=0\), so the function is well defined
on the cycle, including its final edge. We have \(d\Phi(e_j)=y_j\).
If \(d\Phi+u\mathbf1_E=d\Psi+w\mathbf1_E\), the sum over edges gives
\(nu=nw\). Then \(d(\Phi-\Psi)=0\); connectivity of the cycle
and normalization at the initial vertex imply \(\Phi=\Psi\).
The identity \(\sum_j(\Phi(j+1)-\Phi(j))=0\) directly proves
telescopic cancellation.
\end{proof}

\begin{definition}[Admissible transport and normalized reading]
\label{x_reciprocidad:vii:def:observador-admisible}
Transport is admissible when its harmonic mode belongs to
\(\mathbb Rv\). On the space \(\mathcal T_{\rm adm}\) of those
transports we define the global reading and its density per phase:
\begin{equation}
 \mathfrak a_v(\omega)
 =\frac{\langle B(\omega),v\rangle}{\langle v,v\rangle},
 \qquad
 \mu_v(\omega)=\frac{\mathfrak a_v(\omega)}n.
 \label{x_reciprocidad:vii:eq:observador-lectura}
\end{equation}
We have \(B(\omega)=\mathfrak a_v(\omega)v\).
\end{definition}

\begin{theorem}[Determination of the intrinsic reading]
\label{x_reciprocidad:vii:thm:observador-unicidad}
The quotient
\(\mathcal T_{\rm adm}/dC^0(C_n;W)\) is one-dimensional, generated
by the class of \(v\mathbf1_E\). The map \(\mu_v\) is the
unique linear functional on \(\mathcal T_{\rm adm}\) satisfying
\[
 F(\omega+d\Phi)=F(\omega),\qquad F(v\mathbf1_E)=1.
\]
\end{theorem}

\begin{proof}
The preceding decomposition writes each admissible transport as
\(d\Phi+\lambda v\mathbf1_E\). Its class modulo exact cochains
is determined by \(\lambda\). The class of \(v\mathbf1_E\)
is nonzero because its sum is \(nv\neq0\). By linearity and
invariance, every functional with the indicated characteristics equals
\(F(\omega)=\lambda F(v\mathbf1_E)=\lambda\). The formula
\eqref{x_reciprocidad:vii:eq:observador-lectura} gives precisely \(\mu_v=\lambda\).
The global reading is \(\mathfrak a_v=n\mu_v\).
\end{proof}

\begin{definition}[Edge Observer]
\label{x_reciprocidad:vii:def:observador-borde}
Given the polarization of the published transport, a boundary observer
is a pair \(\mathcal O=(\varepsilon,\Psi)\), where
\(\varepsilon\in\{+1,-1\}\) and \(\Psi\in C^0(C_n;W)\).
Its reading action is
\begin{equation}
 \omega^{\mathcal O}=\varepsilon\omega+d\Psi.
 \label{x_reciprocidad:vii:eq:observador-accion}
\end{equation}
The sign fixes the orientation and the cochain fixes the trivialization.
Two polarization choices are compared through their published
transport; the relation \eqref{x_reciprocidad:vii:eq:observador-accion} describes
exactly the representatives differing by orientation and coboundary.
\end{definition}

\begin{proposition}[Covariance of the reading]
\label{x_reciprocidad:vii:prop:observador-covarianza}
The action preserves admissibility and satisfies
\[
 B(\omega^{\mathcal O})=\varepsilon B(\omega),\quad
 \mathfrak a_v(\omega^{\mathcal O})=\varepsilon\mathfrak a_v(\omega),
 \quad \mu_v(\omega^{\mathcal O})=\varepsilon\mu_v(\omega).
\]
Trivializations with the same orientation publish the same closure.
\end{proposition}

\begin{proof}
The sum of the coboundary \(d\Psi\) is zero. Therefore, the harmonic mode
transforms into \(\varepsilon u\), which still belongs to
\(\mathbb Rv\), and the defect transforms into \(\varepsilon B\).
The two scalar equalities are obtained by normalized projection.
\end{proof}

The unit of this reading is the residual representative \(v\), with its
internal normalization. Trivialization freedom acts on the
cochain representative; closure information resides in its
class. The description allows observers to be compared while keeping
genealogical transport, its linear publication and the scalar
evaluating the residual class separate.

\section{Boundary identity and cellular interior}
\label{x_reciprocidad:vii:sec:observador-stokes}

Let \(K\) be a finite oriented two-dimensional cell complex,
equipped with a fundamental \(2\)-chain \(c_K\) such that
\(\partial c_K=[C_n]\). This condition expresses that interior
incidences cancel with their multiplicities and orientations.
Let \(\Omega\in C^1(K;W)\) be an extension of \(\omega\) and
\(F_\Omega=d\Omega\) its additive discrete curvature.

\begin{theorem}[Cellular holographic identity of the reading]
\label{x_reciprocidad:vii:thm:observador-stokes}
With the preceding data,
\begin{equation}
 B(\omega)=\langle\Omega,\partial c_K\rangle
          =\langle d\Omega,c_K\rangle,
 \qquad
 \mathfrak a_v(\omega)
 =\frac{\langle\langle d\Omega,c_K\rangle,v\rangle}
        {\langle v,v\rangle}.
 \label{x_reciprocidad:vii:eq:observador-stokes}
\end{equation}
In particular, \(F_\Omega=0\) implies \(\mathfrak a_v(\omega)=0\).
If closure is nonzero, every extension of transport on that complex
has nonzero discrete curvature.
\end{theorem}

\begin{proof}
The pairing between chains and cochains satisfies
\(\langle d\Omega,c_K\rangle=\langle\Omega,\partial c_K\rangle\).
When it is expanded, each interior edge appears with the sum of the incidence
numbers of the adjoining cells; the condition on
\(\partial c_K\) leaves exactly the oriented boundary cycle.
The restriction of \(\Omega\) to that cycle is \(\omega\), whose sum
is \(B(\omega)\). Projection onto \(v\) gives the second formula
and the flatness implication.
\end{proof}

The equality compares two evaluations of the same extended cochain:
the accumulated boundary defect and the coboundary integrated over the interior.
Its curvature takes values in \(W\) and uses the additive cellular coboundary.
Realization by a geometric connection additionally preserves
its own coefficient space and transport law.

\section{Isometric inscription of the record in the apparatus}
\label{x_reciprocidad:vii:sec:observador-aparato}

Let \(\mathfrak M_n\) be a nonempty finite set of records
published at level \(n\), with result, route, carry, sheet and memory.
Its record space is
\(\mathcal H_n=\ell^2(\mathfrak M_n)\), with basis \(|m\rangle\).
The visible update \(F_n\) has a record lift
\[
 \widehat F_n(m)=(F_n(m),m).
\]
We take as records of the next level a set containing those
pairs. The second coordinate preserves the antecedent and makes the lift
injective. This construction specifies the memory retained by
inscription, including when the visible update merges values.

In a finite apparatus space \(\mathcal K_{A,n}\), we fix an
orthonormal family \((|a_m\rangle)_{m\in\mathfrak M_n}\) and define
\[
 \iota_n:\mathcal H_n\longrightarrow\mathcal K_{A,n},
 \quad \iota_n|m\rangle=|a_m\rangle,
 \qquad
 V_n:\mathcal H_n\longrightarrow\mathcal H_{n+1},
 \quad V_n|m\rangle=|\widehat F_n(m)\rangle.
\]
Let \(\mathcal P_n=\iota_n\mathcal H_n\) be the pointer subspace.

\begin{theorem}[Compatibility between updating and inscription]
\label{x_reciprocidad:vii:thm:observador-entrelazador}
The maps \(\iota_n\) and \(V_n\) are isometries. There exists a unique
isometry \(W_n^0:\mathcal P_n\longrightarrow\mathcal K_{A,n+1}\)
satisfying
\begin{equation}
 W_n^0\iota_n=\iota_{n+1}V_n.
 \label{x_reciprocidad:vii:eq:observador-cuadrado}
\end{equation}
If \(\dim\mathcal K_{A,n+1}\geq\dim\mathcal K_{A,n}\),
it extends to an isometry of all of \(\mathcal K_{A,n}\).
Through auxiliary complements equalizing dimensions, a unitary
extension between the enlarged spaces is obtained.
\end{theorem}

\begin{proof}
The basis images under \(\iota_n\) are orthonormal; the
images under \(V_n\) are also orthonormal by injectivity of
\(\widehat F_n\). For \(z=\iota_n\psi\) we define
\(W_n^0z=\iota_{n+1}V_n\psi\). The representation of \(z\) is unique,
and the three maps preserve its norm. The formula determines the map on
a basis of \(\mathcal P_n\) and proves the square.
If the target dimension is sufficient, the orthogonal complement of
\(W_n^0\mathcal P_n\) contains an orthonormal family of the size of
\(\mathcal P_n^\perp\); sending a basis of the latter to that family
constructs the extension. Finally both spaces are enlarged to
the same dimension and the two prescribed families are completed to orthonormal
bases. The map between those bases is unitary.
\end{proof}

When the record includes labels of a finite instrument, this same
inscription realizes its environment: if the operators
\(K_{y,a}:\mathcal H_S\to\mathcal H'_S\) satisfy
\(\sum_{y,a}K_{y,a}^*K_{y,a}=I\), the map
\[
 V_{\rm inst}\psi=\sum_{y,a}K_{y,a}\psi\otimes|y,a\rangle
\]
is isometric, since its squared norm is
\(\sum_{y,a}\langle\psi,K_{y,a}^*K_{y,a}\psi\rangle=\|\psi\|^2\).
The inscription \(\iota\) transports each label to its corresponding
pointer. The model thus fixes the premeasurement record; a
material device additionally requires its interaction and calibration.


% BEGIN OWNER /Users/ruben/Documents/New project/output/EXTREMA_MEDIA_RAZON_RECIPROCIDAD_ES_EN_20260918/ARTICULO/ES/source/antecedentes_vii/80b_salida_y_condicionamiento.tex
\section{Outcome map and record conditioning}
\label{x_reciprocidad:vii:sec:salida-condicionamiento}

The preceding inscription is composed with the apparatus reader on the
complete genealogical state. Coupling preserves the operation that
produces the correlation; the output map publishes the result;
conditioning restricts the prolongations compatible with that register.

Let \((\Omega_S,\Sigma_S)\) and \((\Omega_M,\Sigma_M)\) be the
measurable spaces of genealogical states of the system and apparatus,
obtained from their compatible cylinders and records. A measurement
context is written \(\mathfrak M_a=(U_a,q_a,\mathcal C_a,m_0)\):
\(m_0\) is the apparatus preparation, \(\mathcal C_a\) specifies
compatibilities and the extension horizon, and
\[
 U_a:\Omega_S\times\Omega_M\longrightarrow\Omega_S\times\Omega_M,
 \qquad q_a:\Omega_M\longrightarrow Y_a
\]
are, respectively, the measurable coupling and the measurable readout into
the outcome space \((Y_a,\Sigma_{Y_a})\). In a reversible regime,
\(U_a\) is bijective and bimeasurable; the corresponding unitarity is
also required in the Hilbert-space realization. The individual outcome is the
composition
\begin{equation}
 \operatorname{Out}_a(x)
 =q_a\bigl(\operatorname{pr}_M U_a(x,m_0)\bigr).
 \label{x_reciprocidad:vii:eq:out-genealogico}
\end{equation}
It is determined by the complete state, apparatus preparation, and
context. The statistical description of a preparation is another readout
of the set of compatible states.

\begin{proposition}[Outcome measure and restriction to a fibre]
\label{x_reciprocidad:vii:prop:medida-salida}
Let \(\Omega_a\in\Sigma_S\) be the set of states compatible with the
preparation and context, endowed with a probability measure \(\mu_a\)
on \((\Omega_a,\Sigma_S|_{\Omega_a})\).
The map \(\operatorname{Out}_a|_{\Omega_a}\) is measurable and induces
the outcome probability measure
\begin{equation}
 \nu_a(B)=\mu_a\bigl(\operatorname{Out}_a^{-1}(B)\cap\Omega_a\bigr),
 \qquad B\in\Sigma_{Y_a}.
 \label{x_reciprocidad:vii:eq:medida-imagen-salida}
\end{equation}
For an outcome \(y\) whose singleton is measurable, the fibre
\(\Omega_{a,y}=\{x\in\Omega_a:\operatorname{Out}_a(x)=y\}\)
is measurable. Whenever \(\mu_a(\Omega_{a,y})>0\),
\begin{equation}
 \mu_{a,y}(E)=
 \frac{\mu_a(E\cap\Omega_{a,y})}{\mu_a(\Omega_{a,y})}
 \label{x_reciprocidad:vii:eq:condicionamiento-salida}
\end{equation}
for \(E\in\Sigma_S|_{\Omega_a}\) defines a probability measure
concentrated on that fibre.
\end{proposition}
\begin{proof}
The inclusion \(x\mapsto(x,m_0)\), coupling, apparatus projection, and
readout are measurable; their composition is therefore measurable.
Preimages preserve disjoint unions and the whole space, so
\eqref{x_reciprocidad:vii:eq:medida-imagen-salida} is countably additive with total mass
one. The fibre is the preimage of the singleton \(\{y\}\) within
\(\Omega_a\). Restricting \(\mu_a\) to that fibre gives a positive,
countably additive measure; division by its positive mass normalizes it.
\end{proof}

For a prefix \(s\) of depth \(n\), the extensions following the record
form
\[
 \operatorname{Fut}_{a,y}(s)
 =\{x\in\Omega_{a,y}:x|_n=s\}.
\]
The restriction preserves the prefix already traversed and the coupling
register; it modifies the family of admitted prolongations. The output
\eqref{x_reciprocidad:vii:eq:out-genealogico} acts on complete states; the outcome law
\eqref{x_reciprocidad:vii:eq:medida-imagen-salida} is obtained by transporting
the preparation measure through that map. Representation by a
density matrix expresses the statistical
level of the quantum realization and retains its conditions of
correspondence with the genealogical fibers.

% END OWNER


\section{Accessible algebra and pointer complement}
\label{x_reciprocidad:vii:sec:observador-expectativa}

Define \(P_m=|a_m\rangle\langle a_m|\),
\(P=\sum_mP_m\), and \(P_\perp=I-P\). The algebra of the pointer
subspace is \(P\mathcal B(\mathcal K_{A,n})P\), whose unit is \(P\).
On it, the diagonal reader is
\[
 \mathcal D_P(A)=\sum_mP_mAP_m.
\]

\begin{proposition}[Record expectation and unital extension]
\label{x_reciprocidad:vii:prop:observador-pinzamiento}
The map \(\mathcal D_P\) is a completely positive,
trace-preserving and idempotent conditional expectation onto the
pointer algebra, with \(\mathcal D_P(P)=P\). Its image is
\(\bigoplus_m\mathbb CP_m\).
On the algebra of the complete apparatus, the map
\begin{equation}
 \mathcal E_P(A)=\sum_mP_mAP_m+P_\perp AP_\perp
 \label{x_reciprocidad:vii:eq:observador-complemento}
\end{equation}
is a unital, completely positive, trace-preserving
conditional expectation. Its image is
\[
 \left(\bigoplus_m\mathbb CP_m\right)
 \oplus\mathcal B(\mathcal P_n^\perp).
\]
Both maps are orthogonal projections for the Hilbert--Schmidt
inner product. Their spectra are contained in \(\{0,1\}\).
When eliminated components exist, the separation between the
fixed subspace and the kernel is one.
\end{proposition}

\begin{proof}
The projectors \(P_m\) are mutually orthogonal. In the complete
apparatus, the family that adds \(P_\perp\) sums to \(I\). Each term
\(A\mapsto P_mAP_m\) is completely positive: after any
matrix amplification it acts by congruence with \(I\otimes P_m\).
The same holds for \(P_\perp\). The sum preserves that property.
Cyclicity of the trace and the sum of the projectors prove
trace preservation on each domain. Orthogonality annihilates the
cross terms when the map is iterated, so it is idempotent.
Its fixed elements are precisely the stated blocks and the
map is bimodular with respect to that subalgebra.
Moreover,
\(\operatorname{Tr}(A^*P_mBP_m)
 =\operatorname{Tr}((P_mAP_m)^*B)\),
from which self-adjointness for Hilbert--Schmidt follows. A self-adjoint
projection has only the spectral values zero and one when
both subspaces are nonzero. If only the fixed subspace remains,
the map is the identity.
\end{proof}

The sum \(\sum_mP_mAP_m\), applied without the complement to the entire
apparatus, sends \(I\) to \(P\). Its natural unit belongs to the pointer
algebra. The formula \eqref{x_reciprocidad:vii:eq:observador-complemento} also preserves
the auxiliary block; an orthonormal resolution of that block
provides, when desired, diagonalization of the entire apparatus.

If \(U\) is a unitary transformation of the apparatus and
\(P'_m=UP_mU^*\), then
\[
 \mathcal E_{P'}(UAU^*)=U\mathcal E_P(A)U^*.
\]
The equality is obtained by substituting each projector into the sum.
Likewise, for \(A\in\mathcal B(\mathcal P_n)\), the square
\eqref{x_reciprocidad:vii:eq:observador-cuadrado} implies
\[
 \mathcal D_{P_{n+1}}(W_n^0A(W_n^0)^*)
 =W_n^0\mathcal D_{P_n}(A)(W_n^0)^*.
\]
Indeed, each matrix unit \(|a_m\rangle\langle a_{m'}|\)
is transported to the unit corresponding to the distinct records
\(\widehat F_n(m),\widehat F_n(m')\); the diagonal reader preserves
exactly the terms with \(m=m'\).

\section{Factorization of the response through the effective boundary}
\label{x_reciprocidad:vii:sec:observador-factorizacion}

Let \(\mathcal X_\infty^{\rm enr}\) be the domain of compatible enriched
states and let
\(b:\mathcal X_\infty^{\rm enr}\longrightarrow B_\infty\) be its global
boundary trace. We denote by \(B_{\rm ef}=\operatorname{im}b\)
the effectively published boundary. For a response space
\(\mathcal V_v\), we consider
\(I_v:\mathcal X_\infty^{\rm enr}\longrightarrow\mathcal V_v\).

\begin{theorem}[Relational response criterion]
\label{x_reciprocidad:vii:thm:observador-factorizacion}
There exists a unique map \(\overline I_v:B_{\rm ef}\to\mathcal V_v\)
such that \(I_v=\overline I_v\circ b\) if and only if
\begin{equation}
 b(x)=b(x')\quad\Longrightarrow\quad I_v(x)=I_v(x').
 \label{x_reciprocidad:vii:eq:observador-fibra}
\end{equation}
Uniqueness extends to all of \(B_\infty\) when \(b\) is
surjective.
\end{theorem}

\begin{proof}
A composition through \(b\) is constant on each fiber. Conversely,
for \(\beta\in B_{\rm ef}\), we define
\(\overline I_v(\beta)=I_v(x)\) with \(b(x)=\beta\).
Condition \eqref{x_reciprocidad:vii:eq:observador-fibra} guarantees independence
of the representative; every value of \(\overline I_v\) is determined
by some state of that fiber. If \(b\) is surjective,
\(B_{\rm ef}=B_\infty\). Otherwise, the composition equation
determines only the values on \(B_{\rm ef}\).
\end{proof}

This criterion is set-theoretic. For a continuous factorization, the
topologies and the quotient-map condition
of \(b:\mathcal X_\infty^{\rm enr}\to B_{\rm ef}\) are also preserved; with it,
continuity of \(I_v\) is equivalent to that of its factor. This specification
separates determination of a reading from its regularity properties.

Let \(\ell_v:\mathcal X_\infty^{\rm enr}\to\mathcal L_v\) be the
restriction accessible to a local neighborhood. The effective global character
is verified through two states satisfying
\begin{equation}
 \ell_v(x)=\ell_v(x'),\qquad I_v(x)\neq I_v(x').
 \label{x_reciprocidad:vii:eq:observador-testigo-global}
\end{equation}
If \(I_v\) factors through \(b\), these two states have distinct
boundaries. Moreover, the response cannot be written as a function of
\(\ell_v\), since that function would assign the same value to both states.
The difference of boundaries thus acquires a verifiable operator effect.

\section{Ambivalence weights and self-inertial response}
\label{x_reciprocidad:vii:sec:observador-ambivalencia}

For the general formula we consider positive readers
\(A_\partial,V_\partial\) and an energy reader \(Q\)
defined on the declared effective boundary domain.
A realization on
\(B_{\rm av}\subseteq B_{\rm ef}\) is made explicit below. Memory remains as a boundary
coordinate \(\mathsf M_\partial(\beta)\). The ambivalence
hinge provides the two orientations
\[
 r_{\rm av}=\frac{\pi_{\rm HMT}}{\varphi_{\rm HMT}^2},\qquad
 r_{\rm av}^{J}=\frac{\varphi_{\rm HMT}}{\pi_{\rm HMT}^2}.
\]
The superscript \(J\) expresses interchange of the generating pair.
\subsubsection{Normalized realization of the boundary readers.}
The APP--TRIT--TPK modal incidence, developed in the regional
construction, provides the path envelope with matrix
\[
 F_{\rm av}=\begin{pmatrix}0&1\\1&1\end{pmatrix}.
\]
The positive vector \(v=(1,\varphi_{\rm HMT})^{\mathsf T}\)
satisfies \(F_{\rm av}v=\varphi_{\rm HMT}v\). Its normalized
transition and stationary weights are
\begin{equation}
 P_{ij}=\frac{(F_{\rm av})_{ij}v_j}{\varphi_{\rm HMT}v_i},
 \qquad
 P=\begin{pmatrix}0&1\\
 \varphi_{\rm HMT}^{-2}&\varphi_{\rm HMT}^{-1}\end{pmatrix},
 \qquad
 (\mu_{\rm ar},\mu_{\rm vol})
 =\frac{(1,\varphi_{\rm HMT}^{2})}{1+\varphi_{\rm HMT}^{2}}.
 \label{x_reciprocidad:vii:eq:lectores-parry}
\end{equation}
The rows of \(P\) sum to one by
\(\varphi^{-2}+\varphi^{-1}=1\); the stationary equation
\(\mu_{\rm ar}=\mu_{\rm vol}\varphi^{-2}\), together with
normalization, determines the stated weights. This Parry measure
weights the paths of the envelope; the chronological frequency
\((1/3,2/3)\) counts the instants of the modal calendar and retains
its distinct function.

The common memory-amplitude realization is defined on the sector
\(\mathcal X_{\rm av}\subset\mathcal X_\infty^{\rm enr}\) in which
the readers of the two sheets have the factorization
\begin{equation}
 \mathcal L_{\rm vol}(x)=\mu_{\rm vol}\mathcal M(x),\qquad
 \mathcal L_{\rm ar}(x)=\pi_{\rm HMT}\mu_{\rm ar}\mathcal M(x),
 \qquad \mathcal M(x)>0.
 \label{x_reciprocidad:vii:eq:lectores-amplitud-comun}
\end{equation}
Both values belong to the same positive memory-amplitude
line. The areal publication carries the regional mark
\(\pi_{\rm HMT}\) of the return. The dimensionless reading is obtained
by dividing both sections by \(\mathcal L_{\rm vol}(x)\).
Thus the sums of the weights are performed between quantities of the
same type; a possible realization as physical areas and volumes
is then composed with their respective dimensional transducers.

\begin{proposition}[Descent of the normalized readers and the energy clock]
\label{x_reciprocidad:vii:prop:lectores-frontera-explicitos}
On \(B_{\rm av}=b(\mathcal X_{\rm av})\), the rules
\begin{equation}
 A_\partial(b(x))
 =\frac{\mathcal L_{\rm ar}(x)}{\mathcal L_{\rm vol}(x)}
 =r_{\rm av},\qquad V_\partial(b(x))=1
 \label{x_reciprocidad:vii:eq:lectores-normalizados}
\end{equation}
define two positive readers independent of the representative.
Given an orientation \(a\) of the action section, the decision
energy of \eqref{x_reciprocidad:v:ant:eq:ii-kb-aut-energia} determines on states
\begin{equation}
 q_a(x)=\frac{h_a(x)}{108t_0(x)}
       =\frac{\hbar_a(x)}{t_P(x)}\in\mathcal L_E^+,
 \qquad t_P=\frac{54}{\pi_{\rm HMT}}t_0.
 \label{x_reciprocidad:vii:eq:lector-q-accion}
\end{equation}
We write
\[
 b_{\rm av}:=b|_{\mathcal X_{\rm av}}:
 \mathcal X_{\rm av}\longrightarrow B_{\rm av}.
\]
There exists a unique reader \(Q:B_{\rm av}\to\mathcal L_E^+\) such that
\(q_a|_{\mathcal X_{\rm av}}=Q\circ b_{\rm av}\)
if and only if, for any \(x,x'\in\mathcal X_{\rm av}\),
\begin{equation}
 b(x)=b(x')\ \Longrightarrow\
 \frac{h_a(x)}{t_0(x)}=\frac{h_a(x')}{t_0(x')}.
 \label{x_reciprocidad:vii:eq:descenso-reloj}
\end{equation}
In particular, a chart of fixed action and duration satisfies this
condition and publishes \(Q=h_a/(108t_0)\).
\end{proposition}
\begin{proof}
The quotient of \eqref{x_reciprocidad:vii:eq:lectores-amplitud-comun} cancels the
common amplitude and equals
\(\pi_{\rm HMT}\mu_{\rm ar}/\mu_{\rm vol}
=\pi_{\rm HMT}/\varphi_{\rm HMT}^2\).
The result is the same for all representatives of a fiber
of \(b_{\rm av}\) within \(\mathcal X_{\rm av}\) and is positive.
Dividing both readers by the same positive factor
preserves the weighted quotients in which they enter.
For the clock, the factorization criterion of
theorem~\ref{x_reciprocidad:vii:thm:observador-factorizacion}, applied to \(q_a\),
is exactly equivalent to \eqref{x_reciprocidad:vii:eq:descenso-reloj}. Fixed sections
satisfy it; so do variable sections whose
quotient descends through \(b_{\rm av}\). Uniqueness follows from the definition
of \(B_{\rm av}\) as the effective image.
\end{proof}

The value of the reader \(Q\) belongs to the energy line
\(\mathcal L_E=\mathcal L_{\rm act}\otimes\mathcal L_T^*\).
With \(E_1,E_2\) on that same line,
\(E_1E_2/Q^2\) is dimensionless and independent of the positive energy
basis. The symbol \(Q\) in this section exclusively designates
that clock energy; the quantity of heat in the Landauer
realization belongs to the corresponding thermal balance.
Common-amplitude factorization specifies a concrete
realization of the general readers. Every other domain preserves its
boundary maps and its own verification of descent.
Complete memory remains in the enriched history and in its
boundary publication, even though it cancels in the normalized quotient.


We define
\[
 w_A(\beta)=
 \frac{A_\partial(\beta)r_{\rm av}}
 {A_\partial(\beta)r_{\rm av}+V_\partial(\beta)r_{\rm av}^{J}},
 \qquad w_V(\beta)=1-w_A(\beta).
\]
In realization \eqref{x_reciprocidad:vii:eq:lectores-normalizados}, the weights
admit the exact evaluation
\begin{equation}
 w_A=\frac{r_{\rm av}^{\,2}}{r_{\rm av}^{\,2}+r_{\rm av}^{J}}
     =\frac{\pi_{\rm HMT}^{4}}
            {\pi_{\rm HMT}^{4}+\varphi_{\rm HMT}^{5}},\qquad
 w_V=\frac{\varphi_{\rm HMT}^{5}}
            {\pi_{\rm HMT}^{4}+\varphi_{\rm HMT}^{5}}.
 \label{x_reciprocidad:vii:eq:pesos-amplitud-comun}
\end{equation}
Substitution preserves both factors of the definition: the normalized
areal reader and its subsequent weighting by \(r_{\rm av}\).
Generation and regional marking of the return are performed once.
Multiplying numerator
and denominator by \(\varphi_{\rm HMT}^4\pi_{\rm HMT}^2\)
gives the last expression. These weights remain constant in the
common-amplitude sector. With fixed projectors and energies, a
variation of the response may come from the reader \(Q\);
the witness to global dependence additionally requires the states with
equal local restriction and different response of
\eqref{x_reciprocidad:vii:eq:observador-testigo-global}.


Let \(P_A,P_V\) be complementary orthogonal projections of a
response space \(\mathcal H_v\), and let \(E_1,E_2>0\) be
two internal energies fixed for comparison. The response is
\begin{equation}
 \overline I_v(\beta)=\frac{E_1E_2}{Q(\beta)^2}
       \bigl(w_A(\beta)P_A+w_V(\beta)P_V\bigr),
 \qquad I_v=\overline I_v\circ b.
 \label{x_reciprocidad:vii:eq:observador-respuesta-ponderada}
\end{equation}
Energies that are explicit readers of
\(\beta\) are also admitted; in that case they are preserved as such in the formula.

\begin{proposition}[Positivity and covariance of the response]
\label{x_reciprocidad:vii:prop:observador-respuesta-positiva}
The response \eqref{x_reciprocidad:vii:eq:observador-respuesta-ponderada} is
self-adjoint, positive definite and constant on the fibers of \(b\).
Its eigenvalues in the nonzero sectors are
\[
 \lambda_A=E_1E_2Q^{-2}w_A,\qquad
 \lambda_V=E_1E_2Q^{-2}w_V.
\]
The complete exchange
\((A_\partial,r_{\rm av},P_A)
 \leftrightarrow(V_\partial,r_{\rm av}^{J},P_V)\)
permutes the two summands and preserves the operator. A unitary
transport of the response space conjugates that operator and preserves
its spectrum with multiplicities.
\end{proposition}

\begin{proof}
Positivity of the readers gives \(0<w_A,w_V<1\). For
\(z=z_A+z_V\), with \(z_A=P_Az\), \(z_V=P_Vz\),
\[
 \langle z,\overline I_vz\rangle
 =\frac{E_1E_2}{Q^2}
   (w_A\|z_A\|^2+w_V\|z_V\|^2)>0
\]
if \(z\neq0\). The orthogonal decomposition proves the spectral
expressions. Exclusive dependence on \(\beta\) gives constancy
on the fibers and factorization. Exchanging the triples
exchanges numerators, weights and projectors at once, so it
preserves the sum. For a unitary transport \(U\), substitution
\(P_A\mapsto UP_AU^*\), \(P_V\mapsto UP_VU^*\) produces
\(U\overline I_vU^*\), with the same spectrum.
\end{proof}

An exchange implemented within the same space by a unitary
sending \(P_A\) to \(P_V\) requires the two sectors to have
equal dimension. Exchanging the two triples as labels of the
formula is valid with each sector's own dimension.

Condition \eqref{x_reciprocidad:vii:eq:observador-testigo-global} requires comparing
the responses. For example, simultaneously multiplying
\(A_\partial,V_\partial\) by the same positive factor preserves
both weights; if \(E_1,E_2,Q\) are also preserved, the response
remains identical. In contrast, with equal weights and energies, replacing
\(Q\) by \(qQ\), \(q>0\), multiplies the response by \(q^{-2}\).
If two states with equal local restriction realize that change with
\(q\neq1\), they provide the global witness. These are operator
separation conditions; existence of the states is checked
in the boundary domain being realized.

\section{Connection of the joint state and acceleration reading}
\label{x_reciprocidad:vii:sec:observador-autoinercia}

In a differentiable realization of the joint state
\(\Gamma=(x,\mathcal O,m)\), let \(N\) be its realization space,
\(\nabla^{\rm tot}\) its connection and \(\pi_S:N\to N_S\) a
differentiable projection onto the subsystem, equipped with connection
\(\nabla^S\). The connection and projection preserve the provenance
of the transports and readers they realize. The observed frame is
obtained through
\[
 \mathfrak F_{\mathcal O}(\Gamma)
 =\operatorname{Res}_{\mathcal O}(\Gamma,\nabla^{\rm tot}).
\]
For a twice-differentiable trajectory \(\Gamma(t)\),
write \(\gamma_S=\pi_S\circ\Gamma\).

\begin{definition}[Self-inertia and projection defect]
\label{x_reciprocidad:vii:def:observador-autoinercia}
A trajectory is self-inertial with respect to the declared connection
when \(\nabla^{\rm tot}_{\dot\Gamma}\dot\Gamma=0\).
The defect of its projection is
\begin{equation}
 \mathcal K_{\pi_S}(\dot\Gamma,\dot\Gamma)
 =\nabla^S_{\dot\gamma_S}\dot\gamma_S
  -d\pi_S\bigl(\nabla^{\rm tot}_{\dot\Gamma}\dot\Gamma\bigr).
 \label{x_reciprocidad:vii:eq:observador-defecto-proyeccion}
\end{equation}
\end{definition}

\begin{proposition}[Subsystem acceleration]
\label{x_reciprocidad:vii:prop:observador-aceleracion}
For a self-inertial trajectory,
\[
 a_S=\nabla^S_{\dot\gamma_S}\dot\gamma_S
     =\mathcal K_{\pi_S}(\dot\Gamma,\dot\Gamma).
\]
In coordinates \(x^B\) on \(N\) and \(y^a\) on \(N_S\), the
defect is quadratic in velocity, with components
\begin{equation}
 (\mathcal K_{\pi_S})^a{}_{BC}
 =\partial_B\partial_C\pi_S^a
  +(\Gamma^S)^a{}_{bc}\,
        \partial_B\pi_S^b\partial_C\pi_S^c
  -\partial_D\pi_S^a(\Gamma^{\rm tot})^D{}_{BC}.
 \label{x_reciprocidad:vii:eq:observador-defecto-coordenadas}
\end{equation}
\end{proposition}

\begin{proof}
The chain rule gives
\(\ddot\gamma_S^a=\partial_B\pi_S^a\ddot\Gamma^B
 +\partial_B\partial_C\pi_S^a\dot\Gamma^B\dot\Gamma^C\).
Adding the contribution of \(\nabla^S\) and subtracting the image of
the total acceleration, the terms containing \(\ddot\Gamma\)
cancel. The remaining ones are
\eqref{x_reciprocidad:vii:eq:observador-defecto-coordenadas} contracted with
\(\dot\Gamma^B\dot\Gamma^C\). When total acceleration vanishes,
the stated equality results.
\end{proof}

The geometric response factors through \(b\) when the connection,
projection and kinematic datum used in evaluation are
determined by \(\beta=b(x)\), or when those kinematic data are
kept fixed in the comparison. Under those conditions,
\eqref{x_reciprocidad:vii:eq:observador-defecto-coordenadas} gives the same result
for any representatives of a fiber;
theorem~\ref{x_reciprocidad:vii:thm:observador-factorizacion} constructs its unique
factor on \(B_{\rm ef}\). If the kinematic datum varies, it
forms part of the domain of the response being compared.

The self-inertial formulation describes a geodesic trajectory of the
joint state and an effective acceleration of its subsystem due
to the projection defect. The gravitational realization then identifies
that connection, the frame and the response with their physical
observables. Observer covariance, memory inscription and
factorization through the boundary thus retain complementary functions:
the first determines the reading class, the second retains the
genealogy and the third specifies the relational dependence of the
response.

% END OWNER


% BEGIN OWNER /Users/ruben/Documents/New project/output/EXTREMA_MEDIA_RAZON_RECIPROCIDAD_ES_EN_20260918/ARTICULO/ES/source/antecedentes_vii/81_propagacion_bidireccional.tex
% Recuperación de 13_accion_interaccion_absorbedor y83_operador_absorbedor.
% Desarrollo del canal cronológico desde el transporte de rutas.
% Procedencia y condiciones en PROCEDENCIA_PROPAGACION_BIDIRECCIONAL.md.
% Los originales y la sección80 permanecen íntegros. Sin contenido de radión.

\chapter{Bidirectional propagation and absorption condition at the boundary}
\label{x_reciprocidad:vii:sec:propagacion-bidireccional}

APP--TRIT--TPK transport preserves the route orientation and memory
distinguishing its outward journey from its return. On that antecedent, three
successive operations are constructed: the action of an enriched history, propagation
of a source through the realized transports and imposition of a
common condition at the boundary. Inscription of the record in an apparatus
space, established in section~\ref{x_reciprocidad:vii:sec:observador-aparato}, allows
those operations to be expressed through linear maps without confusing the
observable publication with the state preserving its antecedents.

Temporal direction, sheet conjugation and spatial reflection
have distinct domains. Their relations are first fixed on the record;
the intertwinings used by a dynamical realization are then
specified. This order determines which information remains in the propagators
and which condition is introduced by choosing an absorbing boundary.

\section{Route action and record involutions}
\label{x_reciprocidad:vii:subsec:accion-bidireccional}

A finite route contains its positions, cardinal directions
\((d_1,\ldots,d_N)\), sheets \((\mu_1,\ldots,\mu_N)\), with
\(\mu_t\in\{\Sigma,\Pi\}\), and the increments of its enriched
record. For \(\lambda\geq0\), we define
\begin{equation}
 \mathcal A_\lambda(\gamma)
 =\sum_{t=2}^N
 \bigl(1-\delta_{d_t,d_{t-1}}
       +\lambda(1-\delta_{\mu_t,\mu_{t-1}})\bigr).
 \label{x_reciprocidad:vii:eq:accion-ruta-bidireccional}
\end{equation}
The first term counts changes of direction; the second counts sheet
changes with the declared weight. For \(N\leq1\), the sum is zero.
The sector's general action additionally preserves its boundary and
memory terms:
\begin{equation}
 \mathcal S[\gamma]
 =\sum_{a\subset\gamma}\mathcal L_a
  +\mathcal B_{\rm in}(\partial_-\gamma)
  +\mathcal B_{\rm out}(\partial_+\gamma)
  +\mathcal M(\gamma).
 \label{x_reciprocidad:vii:eq:accion-enriquecida-frontera}
\end{equation}
The terms of this expression specify a variational class with
typed endpoints and preserved memory. The sector's dynamics determines
\(\mathcal L_a\), the admissible conditions and the term \(\mathcal M\).

On the records we define \(\mathcal C_{\rm r}\) as the global
exchange of sheets; \(\mathcal P_{\rm r}\) as a reflection of the spatial
chart and its directions; and \(\mathcal T_{\rm r}\) as reversal
of temporal order, replacement of each direction by its opposite and
exchange of endpoints. Cardinal reflection is an involutive permutation.
Inversion of the signed record is
\(C_M(a_1,\ldots,a_N)=(-a_N,\ldots,-a_1)\).
The sheets, quotients and memory are transported with the complete route;
the scalar reading of a word constitutes a subsequent operation.

\begin{proposition}[Invariance of the local route action]
\label{x_reciprocidad:vii:prop:invariancia-ruta-bidireccional}
On the preceding domain of records,
\[
 \mathcal A_\lambda(\gamma)
 =\mathcal A_\lambda(\mathcal C_{\rm r}\gamma)
 =\mathcal A_\lambda(\mathcal P_{\rm r}\gamma)
 =\mathcal A_\lambda(\mathcal T_{\rm r}\gamma)
 =\mathcal A_\lambda(\mathcal C_{\rm r}\mathcal P_{\rm r}
                              \mathcal T_{\rm r}\gamma).
\]
\end{proposition}
\begin{proof}
A global permutation of the direction alphabet preserves equality
between consecutive directions. Exchanging the two sheets preserves
equality between consecutive sheets. Reversing the order establishes a
bijection between adjacent pairs, and replacement by opposite directions
also preserves their equalities. Each operation preserves the two counts
of \eqref{x_reciprocidad:vii:eq:accion-ruta-bidireccional}; their composition preserves them.
\end{proof}

The proposition determines the symmetry of the local functional. For a
transformation also to act on a dynamical problem with fixed boundaries,
it must transport its class of admissible prolongations and the other three
terms of \eqref{x_reciprocidad:vii:eq:accion-enriquecida-frontera}. In particular, two
routes with the same final publication may retain distinct variational
classes through their memory.

An interaction between two sections realized at a boundary \(B\)
uses a coupling map
\[
 \mathcal I_B:E_B^{(1)}\otimes E_B^{(2)}\longrightarrow E_B^{(12)}
\]
and a term \(\mathcal S_{\rm int}\) in the same action. Its covariance
is expressed by
\[
 \rho_{12}(g)\mathcal I_B
 =\mathcal I_B\bigl(\rho_1(g)\otimes\rho_2(g)\bigr).
\]
When composing the histories, the shared boundary record is counted
once; memories on both sides are transported to the composite section.
The coupling and this covariance condition are data of the concrete
interaction, distinct from the record involutions.

\section{Reversal of the realized evolution}
\label{x_reciprocidad:vii:subsec:reversion-evolucion}

The realization of the record provides a real Hilbert space
\(\mathcal H_{\mathbb R}\), an orthogonal complex structure
\(J_E\), with \(J_E^2=-I\), and a self-adjoint Hamiltonian \(H\).
Its domain \(\mathcal D(H)\) is invariant under \(J_E\); we assume
\(HJ_E=J_EH\) on that domain. Let \(\mathcal C\) be an orthogonal
involution preserving \(\mathcal D(H)\) and satisfying
\begin{equation}
 \mathcal C J_E\mathcal C^{-1}=-J_E,\qquad
 \mathcal C H\mathcal C^{-1}=H.
 \label{x_reciprocidad:vii:eq:dominio-reversion-evolucion}
\end{equation}
The internal action constant \(\hbar_{\rm int}>0\) retains its
prior generation; here it normalizes the evolution group.

\begin{theorem}[Temporal conjugation of the group]
\label{x_reciprocidad:vii:thm:reversion-evolucion}
The generator \(A=-J_EH/\hbar_{\rm int}\), with domain \(\mathcal D(H)\),
is skew-adjoint. Its orthogonal group, unitary with respect to \(J_E\), satisfies
\begin{equation}
 U(t)=\exp(-J_EHt/\hbar_{\rm int}),\qquad
 \mathcal C U(t)\mathcal C^{-1}=U(-t).
 \label{x_reciprocidad:vii:eq:reversion-grupo-bidireccional}
\end{equation}
\end{theorem}
\begin{proof}
Orthogonality and \(J_E^2=-I\) imply \(J_E^*=-J_E\).
Domain invariance and commutation with \(H\) allow
\(H\) to be treated as a complex self-adjoint operator in the space defined by
\(J_E\). By functional calculus, \(-J_EH/\hbar_{\rm int}\) is
skew-adjoint and generates the stated group. In finite dimension, the same
statements follow from the exponential series and from \(A^*=-A\).
The hypotheses on \(\mathcal C\) give
\(\mathcal C A\mathcal C^{-1}=-A\). Conjugating the group, or its functional
calculus, produces \(\exp(-tA)=U(-t)\).
\end{proof}

The map realizing the signed involution \(C_M\) as \(\mathcal C\)
preserves temporal orientation and the preceding domain. Charge conjugation of
a field representation additionally preserves its own
characters and operators. The equality
\eqref{x_reciprocidad:vii:eq:reversion-grupo-bidireccional} fixes the time reversal
of the group under the stated conditions.

\section{Propagators of the chronological channel of a route}
\label{x_reciprocidad:vii:subsec:green-cronologico}

Isometric inscription of the record admits the unitary extensions
of theorem~\ref{x_reciprocidad:vii:thm:observador-entrelazador}. Let us fix a bilateral
realization of a chronological channel with fibers \(\mathcal H_n\),
\(n\in\mathbb Z\), and unitary transports
\(U_n:\mathcal H_n\to\mathcal H_{n+1}\).
The choice includes auxiliary spaces when required; it preserves
the route and its preceding records. We denote by \(U(n,k)\) the transport
from \(k\) to \(n\), with
\[
 U(k,k)=I,\qquad U(n,k)U(k,l)=U(n,l),\qquad
 U(n,k)^*=U(k,n).
\]
For \(n>k\), it is the product \(U_{n-1}\cdots U_k\); for \(n<k\),
it is the inverse of the transport from \(n\) to \(k\).

The source space \(\mathcal J_c\) contains finitely supported
sections. The response space \(\mathcal F\) contains all
sections of those fibers. The route's quadratic action, in internal
units of this realization, uses
\[
 (D\psi)_n=\psi_{n+1}-U_n\psi_n,\qquad
 E[\psi]=\sum_n\|(D\psi)_n\|^2
\]
for finitely supported sections. The kinetic operator with the sign
convention \(L=-D^*D\) is
\begin{equation}
 (L\psi)_n=U_n^*\psi_{n+1}-2\psi_n+U_{n-1}\psi_{n-1}.
 \label{x_reciprocidad:vii:eq:diferencia-covariante-bidireccional}
\end{equation}
Indeed, \((D^*z)_n=z_{n-1}-U_n^*z_n\), and substitution produces
that formula. The variation of \(E\) is
\(-2\operatorname{Re}\langle\delta\psi,L\psi\rangle\).

\begin{theorem}[Exact Green operators of the covariant difference]
\label{x_reciprocidad:vii:thm:green-cronologico}
Writing \(x_+=\max(x,0)\), the maps
\begin{align}
 (G_{\rm ret}j)_n&=\sum_k(n-k)_+U(n,k)j_k,\label{x_reciprocidad:vii:eq:green-ret-cronologico}\\
 (G_{\rm adv}j)_n&=\sum_k(k-n)_+U(n,k)j_k\label{x_reciprocidad:vii:eq:green-adv-cronologico}
\end{align}
from \(\mathcal J_c\) to \(\mathcal F\) satisfy
\[
 LG_{\rm ret}j=LG_{\rm adv}j=j.
\]
The retarded response vanishes before and at the first position of the
source; the advanced one vanishes after and at the last. They are the unique
solutions vanishing sufficiently far in the respective direction.
For finitely supported \(\psi\), we also have
\(G_{\rm ret}L\psi=G_{\rm adv}L\psi=\psi\).
With respect to the pairing with finitely supported sources,
\(G_{\rm adv}=G_{\rm ret}^{\dagger}\).
\end{theorem}
\begin{proof}
Each sum has as many terms as source positions.
Composition of transports reduces application of \(L\) to
the scalar identity
\[
 (n+1-k)_+-2(n-k)_++(n-1-k)_+=\delta_{nk}.
\]
The function \((k-n)_+\) satisfies the same identity. This proves
the two equations and the supports. For uniqueness, a homogeneous
solution vanishing at two consecutive positions is determined
as zero at the next by \eqref{x_reciprocidad:vii:eq:diferencia-covariante-bidireccional};
the recurrence, using invertible transports, extends in both
directions. The difference of two solutions with the same causal support
vanishes sufficiently far away and is therefore zero.
Applying this argument to \(\psi\) and to \(G_{\rm ret}L\psi\), or to the
advanced case, gives the right-hand identities on finite support.
Finally, for two compactly supported sources, the double sum of the pairing
is finite and
\[
 \bigl((k-n)_+U(k,n)\bigr)^*=(k-n)_+U(n,k).
\]
Interchanging the indices proves the stated adjoint relation.
\end{proof}

The responses belong to \(\mathcal F\); in general they grow
linearly outside the source and do not belong to \(\ell^2\).
The finite-support domain and its pairing specify the Green adjunction
used here. The result constructs propagation of the chronological
channel from its action; a realization of spatial fields additionally preserves
its own operator, geometry and propagation domain.

Let \(C_0\) be an isometric involution, linear or antilinear, of
\(\mathcal H_0\). Transport from the reference fiber constructs
\[
 C_n=U(-n,0)C_0U(0,n):\mathcal H_n\longrightarrow\mathcal H_{-n},
 \qquad (\mathcal C\psi)_n=C_{-n}\psi_{-n}.
\]
We have \(C_{-n}C_n=I\) and
\(C_{n+1}U_n=U_{-n-1}^*C_n\), by composition.
In particular, \(\mathcal C^2=I\), and the change \((n,k)\mapsto(-n,-k)\)
in the preceding sums proves
\begin{equation}
 \mathcal C L=L\mathcal C,\qquad
 \mathcal C G_{\rm ret}\mathcal C^{-1}=G_{\rm adv}.
 \label{x_reciprocidad:vii:eq:intercambio-green-cronologico}
\end{equation}
When \(C_0\) is the declared realization of the record involution,
these formulas transport that involution throughout the channel, without identifying
phase return with a memory reset.

\section{Causal interchange and decomposition of the response}
\label{x_reciprocidad:vii:subsec:green-realizacion-campos}

In a field realization, let
\(L:\mathcal D(L)\subset\mathcal H_\Phi\to\mathcal H_J\)
be an operator with unique retarded and advanced Green operators on the domains of
admissible sources. We denote by \(J^+(S)\) and \(J^-(S)\) the
future and past causal domains of a support \(S\). The conditions are
\[
 \begin{gathered}
 LG_{\rm ret}=LG_{\rm adv}=I,\\
 \operatorname{supp}(G_{\rm ret}j)\subset J^+(\operatorname{supp}j),\\
 \operatorname{supp}(G_{\rm adv}j)\subset J^-(\operatorname{supp}j).
 \end{gathered}
\]
Let \(C_\Phi,C_J\) be isometric involutions, both linear or both
antilinear, that transport these domains,
reverse their causal orientations, and satisfy
\(LC_\Phi=C_JL\). A Green pairing is retained for which
\(G_{\rm adv}=G_{\rm ret}^{\dagger}\).

\begin{proposition}[Conjugation of the Green operators and parities]
\label{x_reciprocidad:vii:prop:green-paridades}
Under the preceding conditions,
\[
 C_\Phi G_{\rm ret}C_J^{-1}=G_{\rm adv}.
\]
On the real space of response operators, define
\[
 \Theta(K)=C_\Phi K C_J^{-1},\qquad
 \mathsf P_\pm(K)=\tfrac12(K\pm\Theta(K)).
\]
These are complementary projections for the involution \(\Theta\), and
\begin{equation}
 G_{\rm sym}=\tfrac12(G_{\rm ret}+G_{\rm adv}),\qquad
 G_{\rm rad}=\tfrac12(G_{\rm ret}-G_{\rm adv})
 \label{x_reciprocidad:vii:eq:descomposicion-green}
\end{equation}
satisfy \(\Theta(G_{\rm sym})=G_{\rm sym}\),
\(\Theta(G_{\rm rad})=-G_{\rm rad}\),
\(G_{\rm sym}^{\dagger}=G_{\rm sym}\), and
\(LG_{\rm rad}=0\). The support of the symmetric response is contained
in the union of the two causal domains.
\end{proposition}
\begin{proof}
The conjugated operator solves the source equation, since
\[
 LC_\Phi G_{\rm ret}C_J^{-1}=C_JLG_{\rm ret}C_J^{-1}=I.
\]
Orientation reversal takes its support to the advanced domain;
its uniqueness identifies it with \(G_{\rm adv}\). The involutions give
\(\Theta^2=I\); expanding \((I\pm\Theta)^2\) and
\((I+\Theta)(I-\Theta)\) proves the statements about projections.
Green adjunction proves symmetry and the difference of the two
source equations gives \(LG_{\rm rad}=0\). The sum of two responses
has support contained in the union of their supports.
\end{proof}

If the involutions are antilinear, \(\Theta\) is linear over
real scalars; this is the operator space used in the proposition.
Conjugation of the temporal group, the support condition, and the
Green adjoint relation are three distinct properties, transmitted by the
corresponding realization maps.

The causal reduction of the bilateral update, developed in the section
on conservation and recovery, identifies another precise origin of
retardation: elimination of a memory component that continues to
participate in the joint evolution. Its discrete kernel preserves the
order of interactions and the contribution of the initial memory.
The Green operators constructed here instead solve a propagation
equation with support conditions. The two constructions can be composed
once the realization map identifying their domains and operators is
specified; their respective causality conditions remain explicit.

\section{Projection of absorbed currents and variation of the action}
\label{x_reciprocidad:vii:subsec:proyeccion-absorcion}

Let \(\gamma_\partial:\mathcal H_\Phi\to\mathcal H_\partial\)
be the boundary reading, distinct from the route \(\gamma\). We define
\[
 \mathcal A_\partial=\gamma_\partial G_{\rm rad}.
\]
The finite window fixes the source space and boundary reading;
the Green operators preserve their response domain. At a finite level with
fixed inner products, its Moore--Penrose
inverse determines
\begin{equation}
 P_{\rm abs}=I-\mathcal A_\partial^+\mathcal A_\partial.
 \label{x_reciprocidad:vii:eq:proyector-absorcion}
\end{equation}

\begin{theorem}[Absorption condition and boundary response]
\label{x_reciprocidad:vii:thm:proyeccion-absorcion}
\(P_{\rm abs}\) is the orthogonal projection onto
\(\ker\mathcal A_\partial\). For any source \(j\), the source
\(j_{\rm abs}=P_{\rm abs}j\) satisfies
\begin{equation}
 \gamma_\partial G_{\rm ret}j_{\rm abs}
 =\gamma_\partial G_{\rm adv}j_{\rm abs}.
 \label{x_reciprocidad:vii:eq:igualdad-borde-absorbido}
\end{equation}
It is the absorbed source closest to \(j\) for the declared inner
product. Nonzero absorbed sources exist exactly when
\(\operatorname{rank}\mathcal A_\partial<\dim\mathcal H_J\).
\end{theorem}
\begin{proof}
The orthogonal decomposition
\(\mathcal H_J=\ker\mathcal A_\partial\oplus
 \operatorname{im}\mathcal A_\partial^*\) shows that
\(\mathcal A_\partial^+\mathcal A_\partial\) is the identity
on the second summand and zero on the first. Its complement is the
stated projector. By definition,
\(2\mathcal A_\partial j_{\rm abs}
 =\gamma_\partial(G_{\rm ret}-G_{\rm adv})j_{\rm abs}=0\).
For \(k\in\ker\mathcal A_\partial\), the orthogonal decomposition gives
\[
 \|j-k\|^2=\|(I-P_{\rm abs})j\|^2+\|P_{\rm abs}j-k\|^2.
\]
The unique minimum corresponds to \(k=P_{\rm abs}j\).
The final equivalence is the rank--nullity theorem.
\end{proof}

In the source-field pairing, the symmetric action is
\(\mathscr S[j]=\tfrac12\langle j,G_{\rm sym}j\rangle\).
The declared adjunction makes that quantity real. For real variations,
\[
 \mathrm d\mathscr S[j](h)
 =\operatorname{Re}\langle h,G_{\rm sym}j\rangle.
\]
The action restricted to absorbed sources is written on the ambient
space as
\begin{equation}
 \mathscr S_{\rm abs}[j]
 =\tfrac12\langle P_{\rm abs}j,G_{\rm sym}P_{\rm abs}j\rangle,
 \quad
 \mathrm d\mathscr S_{\rm abs}[j](h)
 =\operatorname{Re}\langle P_{\rm abs}h,G_{\rm sym}P_{\rm abs}j\rangle.
 \label{x_reciprocidad:vii:eq:variacion-restringida-absorcion}
\end{equation}
When the pairing identifies field and source through Riesz, the ambient gradient
is \(P_{\rm abs}G_{\rm sym}P_{\rm abs}j\).
The field response remains \(G_{\rm sym}j_{\rm abs}\); the
gradient projector expresses restriction of the variations.
The formulas are proved by expanding the quadratic term of
\(j+th\) and using self-adjointness. A material term
\(\mathscr S_{\rm matter}\) adds its own differential.

\section{Explicit absorbing condition in a chronological window}
\label{x_reciprocidad:vii:subsec:absorcion-momentos}

The channel of subsection~\ref{x_reciprocidad:vii:subsec:green-cronologico} allows
that condition to be evaluated without introducing a target response. Let there be a
source supported on \(0,\ldots,N\), with \(N\geq1\). Let us transport
its components to the reference fiber:
\[
 \widehat j_k=U(0,k)j_k,\qquad
 M_0(j)=\sum_{k=0}^N\widehat j_k,\qquad
 M_1(j)=\sum_{k=0}^Nk\widehat j_k.
\]
The identity \((n-k)_+-(k-n)_+=n-k\) gives
\begin{equation}
 (G_{\rm ret}-G_{\rm adv})j\big|_n
 =U(n,0)\bigl(nM_0(j)-M_1(j)\bigr).
 \label{x_reciprocidad:vii:eq:diferencia-green-momentos}
\end{equation}

\begin{proposition}[Absorption and two transported moments]
\label{x_reciprocidad:vii:prop:absorcion-momentos}
If the trace evaluates the field at two distinct positions \(a,b\),
equality of the advanced and retarded traces is equivalent to
\(M_0(j)=M_1(j)=0\). In a fibre of dimension \(d\), the space
of such sources has dimension \((N-1)d\). For \(N\geq2\),
it contains the nonzero sources with transported components
\((z,-2z,z,0,\ldots,0)\), \(z\neq0\).
\end{proposition}
\begin{proof}
Invertibility of \(U(a,0)\) and \(U(b,0)\) reduces the boundary
condition to \(aM_0-M_1=0\) and \(bM_0-M_1=0\).
Subtracting the equations gives \((a-b)M_0=0\), and then \(M_1=0\).
The two-moment map is surjective: for prescribed values
\((z_0,z_1)\), it suffices to take
\(\widehat j_0=z_0-z_1\), \(\widehat j_1=z_1\), and all other
components zero. Its rank is \(2d\); its nullity is
\((N+1)d-2d\). The indicated source has both moments zero.
\end{proof}

In the transported frame, write
\[
 B_N=\begin{pmatrix}1&1&\cdots&1\\0&1&\cdots&N\end{pmatrix}.
\]
The absorbing projector is evaluated by
\begin{equation}
 P_N=\bigl[I-B_N^{\mathsf T}(B_NB_N^{\mathsf T})^{-1}B_N\bigr]
       \otimes I_{\mathcal H_0}.
 \label{x_reciprocidad:vii:eq:proyector-momentos-cronologicos}
\end{equation}
The two rows of \(B_N\) are independent for \(N\geq1\), so
the inverse exists. Unitary transport between frames takes this
projector to the original fibers. For \(N=2\),
\[
 P_2=\frac16
 \begin{pmatrix}1&-2&1\\-2&4&-2\\1&-2&1\end{pmatrix}
 \otimes I_{\mathcal H_0}.
\]
This construction produces a nonzero absorbed class and its projection
operation in the chronological channel. Its amplitude and identification with a
material current preserve the rules of the corresponding physical sector.

\section{Refinement and continuity of the boundary condition}
\label{x_reciprocidad:vii:subsec:refinamiento-absorcion}

Let \(R_m^\Phi,R_m^J,R_m^\partial\) be the refinement maps of
fields, sources and boundaries. Suppose they intertwine the two Green operators and the
trace. Composition gives
\begin{equation}
 R_m^\Phi G_{{\rm sym},m}=G_{{\rm sym},m+1}R_m^J,\qquad
 R_m^\partial\mathcal A_{\partial,m}
 =\mathcal A_{\partial,m+1}R_m^J.
 \label{x_reciprocidad:vii:eq:refinamiento-lectura-absorcion}
\end{equation}
The second equality takes an absorbed source to another absorbed source.
Naturality of the projectors additionally preserves the orthogonal
structure of the realizations.

\begin{proposition}[Naturality of the absorbing projection]
\label{x_reciprocidad:vii:prop:naturalidad-proyeccion-absorcion}
For a source isometry \(R=R_m^J\), we have
\(P_{{\rm abs},m+1}R=RP_{{\rm abs},m}\) exactly when
\[
 R(\ker\mathcal A_{\partial,m})
       \subset\ker\mathcal A_{\partial,m+1},\qquad
 R((\ker\mathcal A_{\partial,m})^\perp)
       \subset(\ker\mathcal A_{\partial,m+1})^\perp.
\]
With these conditions at every level, the projectors define an
orthogonal projection in the completion of the isometric inductive limit
of the source spaces.
\end{proposition}
\begin{proof}
Each source is decomposed into its components in the kernel and its
orthogonal complement. The two inclusions imply that applying the
projection before or after refinement respectively preserves the
first component and annihilates the second. Conversely, equality of
operators applied to each component gives the two inclusions.
In the inductive limit, the equality allows an operator to be defined on the
union of the levels without dependence on the representative. All projectors
have norm at most one, so the operator extends to the
completion. Its self-adjointness and idempotence identities, valid on
the dense union, pass by continuity to the completion.
\end{proof}

To extend a Green response to a completion as well, its
domains and corresponding continuity estimates are preserved.
Finite-support chronological propagation and factorization of the
response through the boundary image, proved in
\ref{x_reciprocidad:vii:thm:observador-factorizacion}, thus retain their respective
domains during refinement.

The physical realization of emission and absorption composes these operators with
the material currents, action and causal conditions of the sector.
The intertwining must preserve the route involution, the source-field
pairing, the transported action and the boundary conditions at
each level. The preceding construction provides reversal of the group, the
Green operators of the chronological channel, temporal decomposition and an explicit
absorbing projection; spatial realizations use their own maps and
causality conditions.

% END OWNER

\chapter{Action realization and metrological comparison}

% BEGIN OWNER /Users/ruben/Documents/New project/output/EXTREMA_MEDIA_RAZON_RECIPROCIDAD_ES_EN_20260918/ARTICULO/ES/source/antecedentes_i/sections/revision_planck_contraste.tex
\section{Comparison of the reduced Planck constant}
\label{x_reciprocidad:sec:revision-planck-contraste}

The International System in force since 20 May 2019 fixes
\(h=6.62607015\times10^{-34}\,\mathrm{J\,s}\) exactly
\cite{x_reciprocidad:bipm2018}. Its reduced constant is therefore
\[
 \hbar_{\mathrm{SI}}=\frac{6.62607015}{2\pi}\,10^{-34}\,
 \mathrm{J\,s}
 =1.0545718176461563912\ldots\times10^{-34}\,\mathrm{J\,s}.
\]
The ellipsis corresponds to the expansion of an exact
expression, not to an experimental uncertainty of \(h\) in the current SI.
This reference value is used here after evaluation of
the HMT readers; its function is comparative.

The two return coordinates of the article are
\begin{align*}
 H_5&=1.0545718176461544696\ldots,\\
 \eta_{\mathrm{ret}}&=1.0545718169150390611\ldots.
\end{align*}
In identifying \(\mathcal U_S\) with \(\mathrm{J\,s}\), the
relative difference of the preceding section with respect to
\(\hbar_{\mathrm{SI}}\) is approximately \(-1.82\times10^{-15}\).
That of the subsequent section is approximately \(-6.93\times10^{-10}\).
The second difference contains the effect of the return used by the
electronic corrector; it is not equivalent to the evaluation error of the
first section.

The precision of these comparisons allows the quantitative proximity
of the construction of \(H_5\) to be recognized while distinguishing
proximity from exact identity. In this manuscript the displayed expressions
do not provide an exact equality with \(h/(2\pi)\).
The proved result concerning \(\Sigma_H\) is its decade \(-34\);
the analytic result concerning \(H_5\) is evaluation of the declared
character. Physical identification is examined in this chart and with the
explicit differences just given. The SI value does not enter
into the definition of the determinantal norm or the regional recurrence.

% END OWNER

